% Lesson 24 — Vocabulary: Words and expressions related to developing and sharing interest in a trade/job/profession — Anglais 1ère A — Cours signé OnBuch+
% Programme officiel MINESEC. Compiler avec : tectonic ENA24-vocabulary-words-and-expressions-related-to-developing.tex
\newcommand{\DOCMATIERE}{Anglais}
\newcommand{\DOCNIVEAU}{1ère A}
\newcommand{\DOCMODULE}{Chapter 2 : Economic Life and Occupations}
\newcommand{\DOCLECON}{ENA24}
\newcommand{\DOCTITRE}{Lesson 24 — Vocabulary: Words and expressions related to developing and sharing interest in a trade/job/profession}
% OnBuch+ — préambule « cours signés » ENRICHI (compilé avec tectonic / XeLaTeX)
% Système visuel « Pop » d'OnBuch+ : fond crème (jamais blanc), bordures encre
% épaisses, ombres « dures » décalées, titres Archivo Black en capitales.
% Version MATHÉMATIQUES : graphes (pgfplots/tikz), boîtes pédagogiques
% (définition, propriété, méthode, attention, le savais-tu, exercice résolu,
% exercice, TP, bilan), page de garde et signature OnBuch+ ancrée en bas.
%
% Macros attendues dans main.tex AVANT \input{preamble} :
%   \DOCMATIERE  \DOCNIVEAU  \DOCMODULE  \DOCLECON (numéro)  \DOCTITRE
\documentclass[11pt]{article}

\usepackage{fontspec}
\usepackage[english,french]{babel} % français = langue principale ; l'anglais via \eng{..} / textebox
\usepackage[a4paper,margin=2cm,top=3.4cm,bottom=2.8cm,headheight=1.4cm,footskip=1.3cm]{geometry}
\usepackage{amsmath,amssymb,amsfonts}
\usepackage{mathtools} % \xmapsto, \coloneqq... souvent utilisés par les modèles
\usepackage{cancel} % simplifications barrées (bilans de forces)
% \abs{z} (module, valeur absolue) et \norm{u} (norme), utilisés par les modèles
\providecommand{\abs}{}\let\abs\relax\DeclarePairedDelimiter{\abs}{\lvert}{\rvert}
\providecommand{\norm}{}\let\norm\relax\DeclarePairedDelimiter{\norm}{\lVert}{\rVert}
\usepackage[table]{xcolor} % [table] : fournit aussi \rowcolors (alternance de lignes)
\usepackage{pagecolor}
\usepackage{tikz}
\usetikzlibrary{arrows.meta,positioning,calc,shapes.geometric,shapes.misc,shapes.arrows,decorations.pathmorphing,decorations.pathreplacing,decorations.markings,patterns,shadows,mindmap,trees,fit,backgrounds,matrix,intersections,angles,quotes}
\usepackage{pgfplots}
\usepackage[version=4]{mhchem} % équations nucléaires \ce{^{235}_{92}U}
\usepackage[european,siunitx]{circuitikz} % schémas électriques
\pgfplotsset{compat=1.18}
\usepgfplotslibrary{fillbetween,groupplots} % aires entre courbes (calcul intégral)
\usepackage{enumitem}
\usepackage{fancyhdr}
% [most] charge toutes les bibliothèques tcolorbox (listings, minted, raster,
% documentation...) et gonfle inutilement la mémoire TeX sur un document long
% (~300 boîtes) : on ne charge que ce qui est réellement utilisé.
\usepackage[skins,breakable]{tcolorbox}
\usepackage{graphicx}
\usepackage{colortbl}
\usepackage{booktabs}
\usepackage{tabularx}
\usepackage{diagbox} % case d'en-tête diagonale des tableaux à double entrée
% Colonnes extensibles centrée / gauche / droite (C, L, R), souvent utilisées par les modèles
\newcolumntype{C}{>{\centering\arraybackslash}X}
\newcolumntype{L}{>{\raggedright\arraybackslash}X}
\newcolumntype{R}{>{\raggedleft\arraybackslash}X}
\usepackage{array}
\usepackage{multicol}
\usepackage{multirow}
\usepackage{siunitx}
% parse-numbers=false : accepte aussi \SI{2,5 \times 10^{-3}}{...}
\sisetup{locale=FR,output-decimal-marker={,},group-minimum-digits=4,parse-numbers=false}
\DeclareSIUnit\ohm{\ensuremath{\symup{\Omega}}} % U+2126 absent de la police
\DeclareSIUnit\division{div} % réglages d'oscilloscope (ms/div, V/div)
\DeclareSIUnit\tr{tr} % tours (vitesse de rotation, ex. \SI{1500}{\tr\per\minute})
\DeclareSIUnit\molar{M} % concentration molaire (ex. \SI{3}{\molar}, \SI{1}{\micro\molar})
\DeclareSIUnit\an{an} % année (le modèle confond parfois avec \year, un compteur TeX) — ex. \SI{5}{\kilo\meter\per\an}
\DeclareSIUnit\FCFA{FCFA} % franc CFA (ex. \SI{500}{\FCFA\per\kilo\gram})
\DeclareSIUnit\degreeBrix{\textdegree Brix} % teneur en sucre d'un sirop/jus (réfractométrie)
\DeclareSIUnit\month{mois} % le modèle utilise parfois \month, un compteur TeX, comme unité de durée
\usepackage{qrcode}
% \degree : commande fréquemment utilisée par les modèles pour un angle ou une
% température en dehors de \SI{}{} (non fournie par les paquets chargés).
\providecommand{\degree}{^{\circ}}
% \qed (fin de démonstration), utilisé par les modèles sans amsthm
\providecommand{\qed}{\hfill$\square$}
% \barre : double barre des valeurs interdites dans un tableau de variations (tkz-tab)
\providecommand{\barre}{\ensuremath{\|}}
% environnement proof (amsthm non chargé) utilisé par les modèles
\ifdefined\proof\else\newenvironment{proof}[1][Démonstration]{\par\noindent\textit{#1.}\ }{\qed\par}\fi
\usepackage{lastpage}
\usepackage{hyperref}
\hypersetup{hidelinks}

% ============================================================
% Palette « Pop » OnBuch+ — mêmes hex que la classe P (plus_ui.dart)
% ============================================================
\definecolor{poporange}{HTML}{F59321}
\definecolor{popdark}{HTML}{E07A0C}
\definecolor{popcream}{HTML}{FFF6E8}
\definecolor{popink}{HTML}{1C1714}
\definecolor{poppurple}{HTML}{7A5AE0}
\definecolor{popgreen}{HTML}{2E9E6A}
\definecolor{poppink}{HTML}{E0457B}
\definecolor{popblue}{HTML}{2F7DE1}
\definecolor{popgold}{HTML}{C98A12}
\definecolor{popmuted}{HTML}{8A7F76}
\definecolor{popline}{HTML}{EADFD2}
\definecolor{popgray}{HTML}{8A7F76}
\colorlet{popgrey}{popgray}
% Alias tolérants (le modèle nomme parfois les couleurs différemment)
\colorlet{popwhite}{white}
\colorlet{popblack}{popink}
\colorlet{popred}{poppink}
\colorlet{popyellow}{popgold}
% Teintes claires (fonds de tableaux/encadrés) — le modèle nomme parfois
% les couleurs avec un suffixe L (ex. popblueL) sans les avoir définies.
\colorlet{poporangeL}{poporange!20}
\colorlet{popdarkL}{popdark!20}
\colorlet{poppurpleL}{poppurple!20}
\colorlet{popgreenL}{popgreen!20}
\colorlet{poppinkL}{poppink!20}
\colorlet{popblueL}{popblue!20}
\colorlet{popgoldL}{popgold!20}
\colorlet{popredL}{popred!20}
\colorlet{popgrayL}{popgray!20}
% Teintes claires pour fonds de tableaux / zones
\colorlet{popblueL}{popblue!10!white}
\colorlet{poporangeL}{poporange!14!white}
\colorlet{popgreenL}{popgreen!12!white}
\colorlet{poppurpleL}{poppurple!12!white}
\colorlet{poppinkL}{poppink!12!white}
\colorlet{popgoldL}{popgold!14!white}

\pagecolor{popcream}

% ============================================================
% Typographie
% ============================================================
\defaultfontfeatures{Ligatures=TeX}
\setmainfont{PlusJakartaSans-Medium.ttf}[Path=../../../fonts/,BoldFont=PlusJakartaSans-Bold.ttf]
% Symboles, API et emojis absents de la police principale : repli sur DejaVu Sans (caractères actifs)
\newfontface\symfont{DejaVuSans.ttf}[Path=../../../fonts/]
\makeatletter
\newcommand\symactive[1]{\catcode#1=13 \begingroup\lccode`\~=#1 \lowercase{\endgroup\def~}{{\symfont\char#1}}}
\newcommand\symblank[1]{\catcode#1=13 \begingroup\lccode`\~=#1 \lowercase{\endgroup\def~}{}}
\newcommand\symhyph[1]{\catcode#1=13 \begingroup\lccode`\~=#1 \lowercase{\endgroup\def~}{\mbox{-}}}
\newcount\symc
\newcommand\symrange[2]{\symc=#1 \@whilenum\symc<#2 \do{\expandafter\symactive\expandafter{\the\symc}\advance\symc by 1 }}
\newcommand\symblankrange[2]{\symc=#1 \@whilenum\symc<#2 \do{\expandafter\symblank\expandafter{\the\symc}\advance\symc by 1 }}
\makeatother
\symrange{"0250}{"02FF}\symactive{"2713}\symactive{"2714}\symactive{"2717}\symactive{"2718}\symactive{"2610}\symactive{"2611}\symactive{"2612}
\symhyph{"2011}\symhyph{"2E1A}\symblankrange{"1F300}{"1FAFF}\symblank{"FE0F}
% Maths en sans-serif (Fira Math) avec chiffres et lettres droites de Plus
% Jakarta, pour que les formules se fondent dans le texte.
\usepackage[math-style=ISO,bold-style=ISO]{unicode-math}
\setmathfont{FiraMath-Regular.otf}[Path=../../../fonts/]
\setmathfont{PlusJakartaSans-Medium.ttf}[Path=../../../fonts/,range={up/num,up/{Latin,latin},\mathup}]
\usepackage{icomma} % virgule décimale sans espace en mode math
% Caractères mathématiques Unicode absents des polices de texte (Plus Jakarta,
% Archivo Black) : rendus par leur équivalent mathématique, en texte comme en
% formule — sinon ils s'affichent en carré vide (ex. « Arithmétique dans ℤ »).
\usepackage{newunicodechar}
\newunicodechar{ℝ}{\ensuremath{\mathbb{R}}}
\newunicodechar{ℤ}{\ensuremath{\mathbb{Z}}}
\newunicodechar{ℂ}{\ensuremath{\mathbb{C}}}
\newunicodechar{ℕ}{\ensuremath{\mathbb{N}}}
\newunicodechar{ℚ}{\ensuremath{\mathbb{Q}}}
\newunicodechar{𝔻}{\ensuremath{\mathbb{D}}}
\newunicodechar{α}{\ensuremath{\alpha}}
\newunicodechar{β}{\ensuremath{\beta}}
\newunicodechar{γ}{\ensuremath{\gamma}}
\newunicodechar{δ}{\ensuremath{\delta}}
\newunicodechar{ε}{\ensuremath{\varepsilon}}
\newunicodechar{θ}{\ensuremath{\theta}}
\newunicodechar{λ}{\ensuremath{\lambda}}
\newunicodechar{μ}{\ensuremath{\mu}}
\newunicodechar{σ}{\ensuremath{\sigma}}
\newunicodechar{τ}{\ensuremath{\tau}}
\newunicodechar{φ}{\ensuremath{\varphi}}
\newunicodechar{ψ}{\ensuremath{\psi}}
\newunicodechar{ω}{\ensuremath{\omega}}
\newunicodechar{Σ}{\ensuremath{\Sigma}}
% majuscules produites par \MakeUppercase dans les titres (α-aminés -> Α)
\newunicodechar{Α}{\ensuremath{\alpha}}
\newunicodechar{Β}{\ensuremath{\beta}}
\newunicodechar{Γ}{\ensuremath{\Gamma}}
\newunicodechar{Δ}{\ensuremath{\Delta}}
\newunicodechar{Ε}{\ensuremath{\varepsilon}}
\newunicodechar{Θ}{\ensuremath{\Theta}}
\newunicodechar{Λ}{\ensuremath{\Lambda}}
\newunicodechar{Μ}{\ensuremath{\mu}}
\newunicodechar{Τ}{\ensuremath{\tau}}
\newunicodechar{Φ}{\ensuremath{\Phi}}
\newunicodechar{Ψ}{\ensuremath{\Psi}}
\newunicodechar{Ω}{\ensuremath{\Omega}}
\newunicodechar{↦}{\ensuremath{\mapsto}}
\newunicodechar{∈}{\ensuremath{\in}}
\newunicodechar{∉}{\ensuremath{\notin}}
\newunicodechar{∧}{\ensuremath{\wedge}}
\newunicodechar{⇔}{\ensuremath{\Leftrightarrow}}
\newunicodechar{⟺}{\ensuremath{\Longleftrightarrow}}
\newunicodechar{⇒}{\ensuremath{\Rightarrow}}
\newunicodechar{⟹}{\ensuremath{\Longrightarrow}}
\newunicodechar{∘}{\ensuremath{\circ}}
\newunicodechar{∪}{\ensuremath{\cup}}
\newunicodechar{∩}{\ensuremath{\cap}}
\newunicodechar{≡}{\ensuremath{\equiv}}
\newunicodechar{⊂}{\ensuremath{\subset}}
\newunicodechar{⊥}{\ensuremath{\perp}}
\newunicodechar{∃}{\ensuremath{\exists}}
\newunicodechar{∀}{\ensuremath{\forall}}
\newunicodechar{∛}{\ensuremath{\sqrt[3]{\,}}}
\newunicodechar{∜}{\ensuremath{\sqrt[4]{\,}}}
\newunicodechar{⃗}{\ensuremath{{}^{\to}}}
\newunicodechar{ⁿ}{\ifmmode^{n}\else\textsuperscript{\ensuremath{n}}\fi}
\newunicodechar{ˣ}{\ifmmode^{x}\else\textsuperscript{\ensuremath{x}}\fi}
\newunicodechar{ᵇ}{\ifmmode^{b}\else\textsuperscript{\ensuremath{b}}\fi}
\newunicodechar{ʳ}{\ifmmode^{r}\else\textsuperscript{\ensuremath{r}}\fi}
\newunicodechar{ᵘ}{\ifmmode^{u}\else\textsuperscript{\ensuremath{u}}\fi}
\newunicodechar{ʸ}{\ifmmode^{y}\else\textsuperscript{\ensuremath{y}}\fi}
\newunicodechar{ᵃ}{\ifmmode^{a}\else\textsuperscript{\ensuremath{a}}\fi}
\newunicodechar{ᵉ}{\ifmmode^{e}\else\textsuperscript{\ensuremath{e}}\fi}
\newunicodechar{ᵏ}{\ifmmode^{k}\else\textsuperscript{\ensuremath{k}}\fi}
\newunicodechar{ᵖ}{\ifmmode^{p}\else\textsuperscript{\ensuremath{p}}\fi}
\newunicodechar{ᴺ}{\ifmmode^{N}\else\textsuperscript{\ensuremath{N}}\fi}
\newunicodechar{ᶜ}{\ifmmode^{c}\else\textsuperscript{\ensuremath{c}}\fi}
\newunicodechar{ʰ}{\ifmmode^{h}\else\textsuperscript{\ensuremath{h}}\fi}
\newunicodechar{ᵠ}{\ifmmode^{q}\else\textsuperscript{\ensuremath{q}}\fi}
\newunicodechar{ₙ}{\ifmmode_{n}\else\textsubscript{\ensuremath{n}}\fi}
\newunicodechar{ₐ}{\ifmmode_{a}\else\textsubscript{\ensuremath{a}}\fi}
\newunicodechar{ᵢ}{\ifmmode_{i}\else\textsubscript{\ensuremath{i}}\fi}
\newunicodechar{ᵣ}{\ifmmode_{r}\else\textsubscript{\ensuremath{r}}\fi}
\newunicodechar{ₖ}{\ifmmode_{k}\else\textsubscript{\ensuremath{k}}\fi}
\newunicodechar{ᵦ}{\ifmmode_{\beta}\else\textsubscript{\ensuremath{\beta}}\fi}
\newunicodechar{ₓ}{\ifmmode_{x}\else\textsubscript{\ensuremath{x}}\fi}
\newfontfamily\popdisplay{ArchivoBlack-Regular.ttf}[Path=../../../fonts/]
\newfontfamily\poplogo{SpaceGrotesk-Bold.ttf}[Path=../../../fonts/]
\newfontfamily\popsemi{PlusJakartaSans-SemiBold.ttf}[Path=../../../fonts/]
\newfontfamily\popxbold{PlusJakartaSans-ExtraBold.ttf}[Path=../../../fonts/]
\newfontfamily\popxboldls{PlusJakartaSans-ExtraBold.ttf}[Path=../../../fonts/,LetterSpace=3.0]

\setlist[enumerate]{leftmargin=1.7em,itemsep=5pt,topsep=4pt}
\setlist[itemize]{leftmargin=1.7em,itemsep=4pt,topsep=4pt,label={\color{poporange}\textbullet}}
\setlength{\parindent}{0pt}
\setlength{\parskip}{5pt}
\renewcommand{\arraystretch}{1.25}

% pgfplots : style Pop par défaut
\pgfplotsset{
  every axis/.append style={axis line style={popink,line width=1pt},tick style={popink},
    grid style={popline,line width=0.5pt},label style={font=\small},tick label style={font=\footnotesize}},
  popcurve/.style={poporange,line width=1.8pt,no markers,smooth},
}
% Même style disponible dans un \draw TikZ simple (hors pgfplots)
\tikzset{popcurve/.style={poporange,line width=1.8pt}}
\tikzset{popgrid/.style={popline,very thin}}
\colorlet{popcurve}{poporange} % le nom du style est parfois utilisé comme couleur

% ============================================================
% Pill / squiggle
% ============================================================
\newcommand{\poppill}[3]{%
  \tikz[baseline=-0.55ex]\node[draw=popink,line width=1.2pt,rounded corners=2.6pt,%
    fill=#2,inner xsep=6.5pt,inner ysep=3pt]{%
    \popxboldls\fontsize{7}{7}\selectfont\color{#3}\MakeUppercase{#1}};%
}
\newcommand{\popsquiggle}[1]{%
  \tikz[baseline=-0.4ex]{
    \draw[#1,line width=2.6pt,line cap=round]
      plot [smooth] coordinates {(0,0.09) (0.34,0.01) (0.68,0.21) (1.02,0.01) (1.36,0.10)};
  }%
}
% Mot-clé surligné (marqueur orange sous le texte)
\newcommand{\cle}[1]{{\popxbold\color{popdark}#1}}

% ============================================================
% En-tête / pied
% ============================================================
\pagestyle{fancy}
\fancyhf{}
\renewcommand{\headrulewidth}{0pt}
\renewcommand{\footrulewidth}{0pt}
\lhead{\raisebox{-0.15ex}{\poplogo\fontsize{13}{13}\selectfont\color{popink}OnBuch\color{poporange}+}}
\rhead{\poppill{\DOCMATIERE\ \textbullet\ \DOCNIVEAU\ \textbullet\ Leçon \DOCLECON}{white}{popink}}
\lfoot{\popxbold\fontsize{7.6}{7.6}\selectfont\color{popmuted}\MakeUppercase{Cours signé OnBuch+}\ \textbullet\ {\popsemi\DOCTITRE}}
\rfoot{\tikz[baseline=-0.6ex]\node[rounded corners=8.5pt,fill=popink,minimum height=17pt,minimum width=17pt,inner xsep=5pt,inner sep=0]{\popxbold\fontsize{8}{8}\selectfont\color{popcream}\thepage\,/\,\pageref*{LastPage}};}

% ============================================================
% Titres
% ============================================================
\newcommand{\chaptitre}[2]{%
  \begin{tcolorbox}[enhanced,colback=poporange,colframe=popink,boxrule=2.6pt,arc=11pt,
      drop shadow=popink,left=15pt,right=15pt,top=12pt,bottom=11pt,before skip=3pt,after skip=18pt]
    \poppill{#2}{popink}{popcream}\par\vspace{9pt}
    {\raggedright\hyphenpenalty=10000\exhyphenpenalty=10000\popdisplay\fontsize{22}{24}\selectfont\color{popcream}\MakeUppercase{#1}\par}
  \end{tcolorbox}
}
% Section numérotée : pastille orange avec le numéro + titre Archivo
\newcounter{popsec}
\newcommand{\coursec}[1]{%
  \stepcounter{popsec}\setcounter{popsub}{0}%
  \par\vspace{16pt}\noindent
  \begin{minipage}{\linewidth}
  \tikz[baseline=-0.7ex]\node[draw=popink,line width=1.6pt,fill=poporange,rounded corners=4pt,minimum size=22pt,inner sep=0]{\popdisplay\fontsize{12}{12}\selectfont\color{popcream}\thepopsec};%
  \hspace{8pt}{\popdisplay\fontsize{15}{17}\selectfont\color{popink}\MakeUppercase{#1}}\par
  \vspace{4pt}\noindent{\color{poporange}\rule{36pt}{3.6pt}}
  \end{minipage}\par\nopagebreak\vspace{8pt}
}
\newcounter{popsub}
\newcommand{\courssub}[1]{%
  \stepcounter{popsub}%
  \par\vspace{9pt}\noindent{\popxbold\fontsize{11.5}{13}\selectfont\color{popink}{\color{poporange}\thepopsec.\thepopsub}\hspace{6pt}#1}\par\nopagebreak\vspace{4pt}
}

% ============================================================
% Boîtes pédagogiques — même recette : fond blanc, bordure épaisse colorée,
% ombre dure encre, badge plein. Titre optionnel [#1].
% ============================================================
\tcbset{popbase/.style={breakable,enhanced,colback=white,boxrule=2.3pt,arc=9pt,
  shadow={3pt}{-3pt}{0pt}{popink},left=14pt,right=14pt,top=11pt,bottom=11pt,before skip=11pt,after skip=15pt,
  fontupper=\popsemi\fontsize{10.6}{14.5}\selectfont\color{popink},
  fontlower=\popsemi\fontsize{10.6}{14.5}\selectfont\color{popink}}}
% Coche dessinée (le glyphe ✓ n'existe pas dans les polices du document)
\newcommand{\popcheck}[1]{\tikz[baseline=-0.5ex]\draw[#1,line width=1.6pt,line cap=round,line join=round](-0.13,0.0)--(-0.04,-0.09)--(0.13,0.1);}
\providecommand{\coche}{\popcheck{popgreen}}
\providecommand{\resultat}[1]{\par\smallskip\noindent\fcolorbox{popgreen}{popgreenL}{\parbox{\dimexpr\linewidth-2\fboxsep-2\fboxrule\relax}{\textbf{Résultat.} #1}}\par\smallskip}
\newcommand{\popboxhead}[3]{\poppill{#1}{#2}{white}\ifx\relax#3\relax\else\hspace{6pt}{\popxbold\fontsize{10.6}{12}\selectfont\color{popink}#3}\fi\par\vspace{7pt}}

\newtcolorbox{aretenir}[1][]{popbase,colframe=popgreen,before upper={\popboxhead{À retenir}{popgreen}{#1}}}
% Texte en anglais (document de lecture, dialogue, extrait) : cadre
% d'étude. Un titre optionnel (source, auteur) s'affiche dans l'en-tête.
\newtcolorbox{textebox}[1][]{popbase,colframe=popblue,before upper={\popboxhead{Text}{popblue}{#1}\begin{otherlanguage}{english}},after upper={\end{otherlanguage}}}
% Marques ✓ / ✗ (forme correcte / fautive) souvent utilisées par les modèles
\providecommand{\cross}{\textcolor{poppink}{\ensuremath{\times}}}
\providecommand{\xmark}{\cross}\providecommand{\crossmark}{\cross}\providecommand{\cmark}{\ensuremath{\checkmark}}
% Ligne de réponse / justification à compléter (inventée par les modèles)
\providecommand{\justification}[1]{\par\noindent\textit{Justification~:} \dotfill\par}
% \textipa (package tipa non chargé) : la police ne couvre pas l'API -> simple passage
\providecommand{\textipa}[1]{#1}
% Soulignement (ulem non chargé) : \uline, \ul
\providecommand{\uline}[1]{\underline{#1}}\providecommand{\ul}[1]{\underline{#1}}
% \glossaire{\item ...} inventée par les modèles : liste de glossaire
\providecommand{\glossaire}[1]{\begin{itemize}#1\end{itemize}}
% \alert (beamer) inventée par les modèles : texte mis en évidence
\providecommand{\alert}[1]{\textbf{\textcolor{poppink}{#1}}}
% variantes ASCII d'ordinaux écrites en macro
\providecommand{\iere}{\textsuperscript{re}}\providecommand{\ieme}{\textsuperscript{e}}\providecommand{\ere}{\textsuperscript{re}}\providecommand{\eme}{\textsuperscript{e}}\providecommand{\er}{\textsuperscript{er}}\providecommand{\ie}{\textsuperscript{e}}
% Ordinaux français écrits en macro par les modèles : 1\ière, 3\ième
\newcommand{\ière}{\textsuperscript{re}}\newcommand{\ième}{\textsuperscript{e}}
% \exemple (inventée par les modèles) : étiquette « Exemple : »
\providecommand{\exemple}{\textbf{Exemple~:}\ }
% \square utilisable aussi hors mode math (cases à cocher)
\AtBeginDocument{\let\squareMath\square\renewcommand{\square}{\ensuremath{\squareMath}}}
% Mot ou phrase en anglais à l'intérieur d'un paragraphe français
\newcommand{\eng}[1]{\ifmmode\text{\textit{#1}}\else\begingroup\selectlanguage{english}\itshape #1\endgroup\fi}
\let\esp\eng % alias toléré
% flèches d'intonation inventées par les modèles
\providecommand{\textsearrow}{}\renewcommand{\textsearrow}{\ensuremath{\searrow}}\providecommand{\textnearrow}{}\renewcommand{\textnearrow}{\ensuremath{\nearrow}}
% \fr{..} : mot français (inventé par les modèles)
\providecommand{\fr}[1]{\textit{#1}}
\newtcolorbox{exemplebox}[1][]{popbase,colframe=poporange,before upper={\popboxhead{Exemple}{poporange}{#1}}}
\newtcolorbox{definition}[1][]{popbase,colframe=popblue,before upper={\popboxhead{Définition}{popblue}{#1}}}
\newtcolorbox{propriete}[1][]{popbase,colframe=poppurple,before upper={\popboxhead{Propriété}{poppurple}{#1}}}
\newtcolorbox{methode}[1][]{popbase,colframe=popgold,before upper={\popboxhead{Méthode}{popgold}{#1}}}
\newtcolorbox{attention}[1][]{popbase,colframe=poppink,before upper={\popboxhead{Attention}{poppink}{#1}}}
\newtcolorbox{experience}[1][]{popbase,colframe=popblue,colback=popblueL,before upper={\popboxhead{Expérience}{popblue}{#1}}}
% « Le savais-tu ? » : carte violette pleine (contexte, culture, Cameroun)
\newtcolorbox{savaistu}[1][]{popbase,colback=poppurple,colframe=popink,
  fontupper=\popsemi\fontsize{10.4}{14.2}\selectfont\color{white},
  before upper={\poppill{Le savais-tu\,?}{white}{poppurple}\ifx\relax#1\relax\else\hspace{6pt}{\popxbold\color{white}#1}\fi\par\vspace{7pt}}}
% Exercice résolu : énoncé en haut, \tcblower puis la solution sur fond crème
\newcounter{popexo}\newcounter{popexr}
\newtcolorbox{exoresolu}[1][]{popbase,colframe=poporange,colbacklower=popcream,
  segmentation style={popink,line width=1.2pt,dashed},
  before upper={\stepcounter{popexr}\popboxhead{Exercice résolu \thepopexr}{poporange}{#1}},
  before lower={\poppill{Solution}{popink}{popcream}\par\vspace{6pt}}}
% Exercice (non résolu en place) — la correction est dans la partie Corrigés
\newtcolorbox{exercice}[1][]{popbase,colframe=popink,
  before upper={\stepcounter{popexo}\popboxhead{Exercice \thepopexo}{popink}{#1}}}
% Corrigé
\newtcolorbox{corrige}[1][]{popbase,colframe=popgreen,colback=popgreenL,
  before upper={\popboxhead{Corrigé}{popgreen}{#1}}}
% Objectifs / prérequis / bilan
\newtcolorbox{objectifs}{popbase,colframe=popink,colback=poporangeL,
  before upper={\popboxhead{Objectifs de la leçon}{popink}{}}}
\newtcolorbox{prerequis}{popbase,colframe=popink,colback=popblueL,
  before upper={\popboxhead{Prérequis}{popblue}{}}}
\newtcolorbox{bilan}{popbase,colframe=popink,colback=white,boxrule=2.8pt,
  before upper={\popboxhead{Fiche bilan}{poporange}{L'essentiel à connaître pour l'examen}}}
% Figure encadrée avec légende
\newtcolorbox{popfigure}[1][]{enhanced,breakable=false,colback=white,colframe=popline,boxrule=1.4pt,arc=8pt,
  left=10pt,right=10pt,top=10pt,bottom=8pt,before skip=10pt,after skip=12pt,halign=center,
  lower separated=false,halign lower=center,
  fontlower=\popsemi\fontsize{9}{11.5}\selectfont\color{popmuted},
  #1}
\newcommand{\legende}[1]{\par\vspace{6pt}{\popsemi\fontsize{9}{11.5}\selectfont\color{popmuted}#1}}

% ============================================================
% Signature OnBuch+
% ============================================================
\newcommand{\poplogoinline}[1]{{\poplogo\fontsize{#1}{#1}\selectfont\color{popink}OnBuch\color{poporange}+}}
\newcommand{\signatureonbuch}{%
  \begin{tcolorbox}[enhanced,colback=popink,colframe=popink,boxrule=2.2pt,arc=10pt,
      drop shadow=poporange,left=14pt,right=14pt,top=10pt,bottom=10pt,before skip=0pt,after skip=0pt]
    \tikz[baseline=-0.7ex]\node[circle,fill=popgreen,draw=popcream,line width=1.3pt,minimum size=22pt,inner sep=0]{\popcheck{white}};%
    \hspace{9pt}\begin{minipage}[c]{0.62\linewidth}
      {\popxbold\fontsize{12}{13}\selectfont\color{popcream}Signé\ \textbullet\ {\poplogo OnBuch\color{poporange}+}}\par\vspace{2pt}
      {\raggedright\popsemi\fontsize{8}{10.5}\selectfont\color{popline}\MakeUppercase{Équipe pédagogique OnBuch+}\\\MakeUppercase{Conforme au programme officiel MINESEC}\par}
    \end{minipage}\hfill
    {\poplogo\fontsize{22}{22}\selectfont\color{popcream}OnBuch\color{poporange}+}
  \end{tcolorbox}%
}

% ============================================================
% Page de garde
% #1 titre · #2 matière · #3 niveau · #4 module · #5 n° leçon · #6 durée ·
% #7 description / objectifs (texte ou itemize)
% ============================================================
\newcommand{\pagegarde}[7]{%
  \thispagestyle{empty}
  \newgeometry{a4paper,margin=1.7cm,top=1.6cm,bottom=1.4cm}%
  \begin{tikzpicture}[remember picture,overlay]
    \fill[poporange!18] ([xshift=-3.2cm,yshift=-3.6cm]current page.north east) circle (4.6cm);
    \fill[poppurple!14] ([xshift=2.4cm,yshift=7.6cm]current page.south west) circle (3.8cm);
  \end{tikzpicture}%
  {\poplogo\fontsize{30}{30}\selectfont\color{popink}OnBuch\color{poporange}+}\hfill
  \raisebox{0.6ex}{\poppill{Cours signé \textbullet\ Premium}{popink}{popcream}}\par
  \vspace{1pt}\popsquiggle{poppurple}\par
  \vspace{8pt}
  \begin{tcolorbox}[enhanced,colback=poporange,colframe=popink,boxrule=2.8pt,arc=13pt,
      drop shadow=popink,left=18pt,right=18pt,top=12pt,bottom=13pt,after skip=10pt]
    \poppill{#2 \textbullet\ #3}{popink}{popcream}\hspace{5pt}\poppill{#4}{white}{popink}\par\vspace{8pt}
    {\popdisplay\fontsize{14}{14}\selectfont\color{popink}LEÇON #5}\par\vspace{4pt}
    {\raggedright\hyphenpenalty=10000\exhyphenpenalty=10000\popdisplay\fontsize{25}{27}\selectfont\color{popcream}\MakeUppercase{#1}\par}
    \vspace{8pt}
    {\popxbold\fontsize{9.5}{9.5}\selectfont\color{popink}\MakeUppercase{Durée conseillée : #6}}
  \end{tcolorbox}
  \begin{tcolorbox}[enhanced,colback=white,colframe=popink,boxrule=2.2pt,arc=10pt,
      drop shadow=popink,left=16pt,right=16pt,top=10pt,bottom=10pt,after skip=8pt]
    \poppill{Dans cette leçon}{poppurple}{white}\par\vspace{5pt}
    \popsemi\fontsize{9.3}{12.6}\selectfont\color{popink}#7
  \end{tcolorbox}
  \noindent\begin{minipage}[t]{0.487\linewidth}
    \begin{tcolorbox}[enhanced,colback=popgreen,colframe=popink,boxrule=2.2pt,arc=10pt,
        drop shadow=popink,left=11pt,right=11pt,top=10pt,bottom=10pt,height=3.1cm,valign=center]
      \tcbox[colback=white,colframe=popink,boxrule=1.3pt,arc=5pt,boxsep=0pt,nobeforeafter,
        left=3pt,right=3pt,top=3pt,bottom=3pt,valign=center]{\qrcode[height=1.9cm]{https://play.google.com/store/apps/details?id=com.luvvix.onbuch&hl=fr}}%
      \hspace{8pt}\begin{minipage}[c]{0.5\linewidth}
        {\popdisplay\fontsize{11}{12.5}\selectfont\color{white}\MakeUppercase{Télécharge OnBuch}}\par\vspace{4pt}
        {\raggedright\popsemi\fontsize{8.4}{11}\selectfont\color{white}Gratuit \textbullet\ Google Play\\Tuteur IA, annales, concours, résultats\par}
      \end{minipage}
    \end{tcolorbox}
  \end{minipage}\hfill
  \begin{minipage}[t]{0.487\linewidth}
    \begin{tcolorbox}[enhanced,colback=poppurple,colframe=popink,boxrule=2.2pt,arc=10pt,
        drop shadow=popink,left=11pt,right=11pt,top=10pt,bottom=10pt,height=3.1cm,valign=center]
      \tikz[baseline=-0.7ex]\node[circle,fill=white,draw=popink,line width=1.3pt,minimum size=1.6cm,inner sep=0]{\popdisplay\fontsize{24}{24}\selectfont\color{poppurple}+};%
      \hspace{8pt}\begin{minipage}[c]{0.6\linewidth}
        {\popdisplay\fontsize{11}{12.5}\selectfont\color{white}\MakeUppercase{Passe à OnBuch+}}\par\vspace{4pt}
        {\raggedright\popsemi\fontsize{8.4}{11}\selectfont\color{white}Tous les cours signés, Tuteur IA illimité, fiches PDF — onglet OnBuch+ de l'app\par}
      \end{minipage}
    \end{tcolorbox}
  \end{minipage}
  \vspace{8pt}
  \popcontenu
  \vfill
  \signatureonbuch
  \restoregeometry
  \newpage
}

% Rangée « Ce cours contient » (page de garde)
\newcommand{\poptile}[3]{% #1 couleur #2 gros texte #3 légende
  \begin{tcolorbox}[enhanced,colback=#1,colframe=popink,boxrule=2pt,arc=9pt,shadow={2.5pt}{-2.5pt}{0pt}{popink},
      left=6pt,right=6pt,top=9pt,bottom=9pt,halign=center,valign=center,height=1.75cm,width=\linewidth,nobeforeafter]
    {\popdisplay\fontsize{12.5}{13}\selectfont\color{white}\MakeUppercase{#2}}\par\vspace{3pt}
    {\centering\popsemi\fontsize{7.8}{9.5}\selectfont\color{white}#3\par}
  \end{tcolorbox}}
\newcommand{\popcontenu}{%
  \par\noindent{\popxboldls\fontsize{7.5}{7.5}\selectfont\color{popmuted}CE COURS SIGNÉ CONTIENT}\par\vspace{7pt}
  \noindent\begin{minipage}[t]{0.235\linewidth}\poptile{popblue}{Cours}{complet, illustré, pas à pas}\end{minipage}\hfill
  \begin{minipage}[t]{0.235\linewidth}\poptile{poporange}{Méthodes}{et exercices résolus}\end{minipage}\hfill
  \begin{minipage}[t]{0.235\linewidth}\poptile{poppink}{Exercices}{type examen + corrigés}\end{minipage}\hfill
  \begin{minipage}[t]{0.235\linewidth}\poptile{popgreen}{Fiche bilan}{l'essentiel à retenir}\end{minipage}\par}

% Sommaire
\newtcolorbox{sommaire}{enhanced,colback=white,colframe=popink,boxrule=2.2pt,arc=10pt,
  drop shadow=popink,left=18pt,right=18pt,top=13pt,bottom=13pt,before skip=6pt,after skip=20pt,
  before upper={\poppill{Sommaire}{poppurple}{white}\par\vspace{9pt}\popsemi\fontsize{10.6}{14.5}\selectfont\color{popink}}}

% Signature de fin de cours — poussée tout en bas de la dernière page
\newcommand{\findecours}{%
  \par\vfill\nopagebreak
  \begin{center}\popsquiggle{poporange}\end{center}\vspace{4pt}
  \signatureonbuch
}
\AtBeginDocument{\def\llbracket{\mathopen{[\mkern-3.5mu[}}\def\rrbracket{\mathclose{]\mkern-3.5mu]}}} % intervalles d'entiers (glyphe ⟦ absent de la police)
% Glyphes absents de Fira Math : redéfinis à partir de glyphes disponibles
\AtBeginDocument{%
  \def\setminus{\mathbin{\backslash}}\let\smallsetminus\setminus
  \def\vdots{\mathord{\vbox{\baselineskip3.2pt\lineskiplimit0pt\kern4pt\hbox{.}\hbox{.}\hbox{.}}}}%
  \def\ddots{\mathinner{\mkern1mu\raise6.5pt\hbox{.}\mkern2mu\raise3.6pt\hbox{.}\mkern2mu\raise0.7pt\hbox{.}\mkern1mu}}%
  \def\triangle{\mathord{\scalebox{0.9}{$\symup{\Delta}$}}}%
  \def\nabla{\mathord{\raisebox{\depth}{\scalebox{1}[-1]{$\symup{\Delta}$}}}}%
  \def\checkmark{\popcheck{popdark}}%
  \def\star{\mathbin{\ast}}\def\bigstar{\mathord{\ast}}%
  \def\diamond{\mathbin{\scalebox{0.6}{$\Diamond$}}}%
  \def\bigcup{\mathop{\mathchoice{\raisebox{-0.3em}{\scalebox{1.7}{$\cup$}}}{\scalebox{1.25}{$\cup$}}{\cup}{\cup}}\displaylimits}%
  \def\bigcap{\mathop{\mathchoice{\raisebox{-0.3em}{\scalebox{1.7}{$\cap$}}}{\scalebox{1.25}{$\cap$}}{\cap}{\cap}}\displaylimits}%
  \def\lesseqqgtr{\mathrel{\lessgtr}}\def\bigoplus{\mathop{\oplus}\displaylimits}%
}
\providecommand{\ssi}{si et seulement si} % abréviation fréquente
\AtBeginDocument{\providecommand{\wideparen}{\overparen}} % arc de cercle

\begin{document}

\pagegarde{Lesson 24 — Vocabulary: Words and expressions related to developing and sharing interest in a trade/job/profession}{Anglais}{1ère A}{Chapter 2 : Economic Life and Occupations}{ENA24}{2 h}{Cette leçon de vocabulaire dote l'élève du lexique nécessaire pour exprimer son intérêt pour un métier, un trade ou une profession, et pour échanger sur les parcours de formation. Champs lexicaux, collocations, faux amis et mises en contexte préparent directement aux épreuves du Probatoire et du Bac.
\vspace{4pt}
\begin{multicols}{2}\small
\begin{itemize}
\item À la fin de cette leçon, je sais nommer les principaux métiers, trades et professions en anglais
\item je sais distinguer les mots job, work, trade, profession, occupation et career
\item je sais employer le vocabulaire pour exprimer mon intérêt pour un métier (to be interested in, to be keen on, to take up...)
\item je sais utiliser les collocations courantes liées au monde du travail (to apply for a job, to earn a living, to undergo training...)
\item je sais reconnaître et éviter les faux amis (formation, stage, carrière, démission...)
\item je sais former des mots de la même famille (employ, employer, employee, employment, unemployment...)
\end{itemize}\end{multicols}}

\begin{sommaire}
\begin{multicols}{2}
\begin{enumerate}
\item Trades, jobs and professions : le champ lexical des métiers
\item Developing an interest : exprimer l'intérêt et la vocation
\item Sharing an interest : échanger sur les métiers
\item Collocations et expressions du monde du travail
\item Faux amis et pièges des francophones
\item Mots de la même famille et réemploi en contexte
\item Activité d'intégration
\item Méthodes et exercices résolus
\item Exercices
\item Corrigés des exercices
\item Fiche bilan
\end{enumerate}
\end{multicols}
\end{sommaire}

\begin{objectifs}
\begin{itemize}
\item À la fin de cette leçon, je sais nommer les principaux métiers, trades et professions en anglais
\item je sais distinguer les mots job, work, trade, profession, occupation et career
\item je sais employer le vocabulaire pour exprimer mon intérêt pour un métier (to be interested in, to be keen on, to take up...)
\item je sais utiliser les collocations courantes liées au monde du travail (to apply for a job, to earn a living, to undergo training...)
\item je sais reconnaître et éviter les faux amis (formation, stage, carrière, démission...)
\item je sais former des mots de la même famille (employ, employer, employee, employment, unemployment...)
\item je sais prononcer et orthographier correctement le lexique étudié
\item je sais réemployer ce lexique dans des phrases, des dialogues et un court texte
\end{itemize}
\end{objectifs}

\begin{prerequis}
\begin{itemize}
\item Vocabulaire de base des métiers vu au collège et en Seconde (teacher, farmer, doctor, driver...)
\item Présent simple et présent be + V-ing pour parler de ses activités et projets
\item Structures d'expression du goût : like + V-ing, want to + V, would like to + V
\item Notions de base sur les questions en wh- (What do you want to be? Why...?)
\end{itemize}
\end{prerequis}

\newpage

\coursec{Trades, jobs and professions : le champ lexical des métiers}

À Yaoundé comme à Buea, de nombreux élèves de Première hésitent encore sur leur orientation : menuisier, infirmier, ingénieur, commerçant ? Au Nigeria ou au Ghana, les \eng{career days} (« journées des carrières ») permettent aux lycéens de rencontrer des professionnels et de découvrir des métiers. Savoir dire en anglais que l'on s'intéresse à un métier, que l'on veut s'y former et en parler avec les autres est une compétence clé — et un thème fréquent à l'épreuve d'anglais du Probatoire et du Bac.

\textbf{Problématique :} comment nommer et classer correctement les métiers en anglais, sans confondre \eng{job}, \eng{work}, \eng{trade} et \eng{profession} ?

\courssub{Observer d'abord : un texte pour découvrir le vocabulaire}

Lisons d'abord ce court texte, rédigé dans le style d'un article de presse anglophone. Soulignez mentalement tous les mots qui désignent des métiers, des lieux de travail ou des personnes au travail.

\begin{textebox}[Career Day at Government High School, Buea]
Last Friday, our school organised its first career day, and the hall was full of curious students. Professionals from many fields came to talk about their jobs. Mr Ekane, a carpenter, explained that he had learnt his trade as an apprentice in a small workshop in Douala. “A trade is learnt with the hands,” he said, “but it also requires patience and honesty.” Next to him, Dr Manka, a surgeon at the regional hospital, described her long years at university before she could enter the medical profession. A young lawyer, a nurse, an engineer and an accountant also answered questions about their careers.

Not every speaker worked in an office. A trader from the market told us how she started her business with very little money, and a mechanic from Mile 16 showed us pictures of his garage. A fisherman from Limbe explained that his work depends on the sea and the seasons.

The headmaster reminded us that work is dignity, whether you are a mason on a construction site, a teacher in a classroom or a civil servant in a ministry. “Choose an occupation you love,” he concluded, “and you will never feel that you are working.” Many students left the hall with a dream: some want to become self-employed, others hope to join a profession after long studies. All of them understood that every job, big or small, helps to build the nation.
\end{textebox}

\textbf{Glossaire :}

\begin{center}
\begin{tabularx}{\linewidth}{|l|l|X|}
\hline
\rowcolor{popblueL} Mot anglais & Nature & Traduction française \\
\hline
career day & nom & journée des carrières (orientation) \\
\hline
trade & nom & métier manuel, artisanal \\
\hline
apprentice & nom & apprenti \\
\hline
workshop & nom & atelier \\
\hline
surgeon & nom & chirurgien \\
\hline
to enter a profession & expression & entrer dans une profession \\
\hline
trader & nom & commerçant(e) \\
\hline
self-employed & adjectif & qui travaille à son propre compte \\
\hline
occupation & nom & profession, activité professionnelle \\
\hline
\end{tabularx}
\end{center}

\textbf{Questions de compréhension :}
\begin{enumerate}
\item Where did Mr Ekane learn his trade?
\item What is the difference between Dr Manka's training and Mr Ekane's training?
\item Name three speakers who do not work in an office.
\item What advice does the headmaster give at the end?
\end{enumerate}

\begin{corrige}[Questions de compréhension]
1. \eng{He learnt his trade as an apprentice in a small workshop in Douala.}

2. \eng{Dr Manka studied for long years at university, whereas Mr Ekane learnt his trade through practice, as an apprentice.}

3. \eng{The trader (market), the mechanic (garage) and the fisherman (the sea) do not work in an office.} On peut aussi citer le maçon (\eng{construction site}).

4. \eng{He advises the students to choose an occupation they love.}
\end{corrige}

\courssub{Les mots-clés : job, work, occupation, trade, profession, career}

Le français utilise souvent le seul mot « métier » ou « travail » là où l'anglais distingue plusieurs termes. C'est la première difficulté pour un francophone.

\begin{definition}[Job]
\eng{A job} désigne un \cle{emploi précis}, un poste occupé par une personne. C'est un nom \textbf{dénombrable} : \eng{a job}, \eng{two jobs}. \eng{She found a job as a secretary.} = Elle a trouvé un emploi comme secrétaire.
\end{definition}

\begin{definition}[Work]
\eng{Work} désigne le \cle{travail} en général, l'activité, l'effort. C'est un nom \textbf{indénombrable} : jamais de \eng{a} devant, jamais de pluriel. \eng{Work is dignity.} = Le travail, c'est la dignité.
\end{definition}

\begin{definition}[Occupation]
\eng{Occupation} est le terme \cle{général} et un peu formel pour « profession, activité professionnelle ». On le rencontre surtout dans les formulaires : \eng{State your name, date of birth and occupation.} = Indiquez votre nom, votre date de naissance et votre profession.
\end{definition}

\begin{definition}[Trade]
\eng{A trade} is a job that needs special manual skills learnt through training or apprenticeship, e.g. a plumber or a tailor. = Un \cle{métier manuel} appris par la pratique, souvent en apprentissage auprès d'un maître artisan.
\end{definition}

\begin{definition}[Profession]
\eng{A profession} is a job that needs a high level of education and training, e.g. a doctor or a lawyer. = Une \cle{profession intellectuelle} exigeant de longues études et des diplômes.
\end{definition}

\begin{definition}[Career]
\eng{Career} désigne la \cle{carrière}, c'est-à-dire le parcours professionnel d'une personne sur la durée. \eng{He had a brilliant career in the army.} = Il a fait une brillante carrière dans l'armée.
\end{definition}

\begin{attention}[Work est indénombrable]
En anglais, \eng{work} est \textbf{indénombrable} : on dit \eng{I am looking for work} (je cherche du travail) et jamais \emph{a work}. Pour un emploi précis, on utilise \eng{a job} : \eng{I am looking for a job.} De même, on dit \eng{a piece of work} pour « un travail » (un devoir, une tâche), jamais \emph{a work}.
\end{attention}

\begin{exemplebox}[Les six mots en contexte]
\eng{My uncle is a carpenter; he learnt his trade in a workshop in Douala.} = Mon oncle est menuisier ; il a appris son métier dans un atelier à Douala.

\eng{She wants to enter the teaching profession.} = Elle veut entrer dans la profession enseignante.

\eng{He lost his job last month, so he is looking for work.} = Il a perdu son emploi le mois dernier, alors il cherche du travail.

\eng{On the form, write your occupation: “farmer”.} = Sur le formulaire, écris ta profession : « agriculteur ».

\eng{Teaching is not just a job; it is a career.} = Enseigner n'est pas seulement un emploi ; c'est une carrière.
\end{exemplebox}

\courssub{Trades : les métiers manuels et artisanaux}

Un \eng{trade} s'apprend par la pratique, souvent comme \eng{apprentice} (apprenti) auprès d'un artisan expérimenté. Voici les métiers manuels les plus courants, très présents dans la vie quotidienne camerounaise :

\begin{center}
\begin{tabularx}{\linewidth}{|l|X|l|}
\hline
\rowcolor{popblueL} English & Prononciation (approx.) & Français \\
\hline
carpenter & « KAR-pen-teur » & menuisier \\
\hline
mason & « MÉ-son » & maçon \\
\hline
plumber & « PLEU-meur » (b muet) & plombier \\
\hline
electrician & « è-lek-TRI-cheun » & électricien \\
\hline
tailor & « TÉ-leur » & tailleur, couturier \\
\hline
mechanic & « mé-KA-nik » & mécanicien \\
\hline
welder & « OUÈL-deur » & soudeur \\
\hline
hairdresser & « HÈR-drè-seur » & coiffeur, coiffeuse \\
\hline
baker & « BÉ-keur » & boulanger \\
\hline
shoemaker & « CHOU-mé-keur » & cordonnier \\
\hline
\end{tabularx}
\end{center}

\begin{exemplebox}[Trades en contexte]
\eng{The tailor in our quarter sews beautiful dresses.} = Le tailleur de notre quartier coud de belles robes.

\eng{A good mechanic is never unemployed in Douala.} = Un bon mécanicien n'est jamais au chômage à Douala.

\eng{The plumber repaired our tap in ten minutes.} = Le plombier a réparé notre robinet en dix minutes.

\eng{Her brother is an apprentice welder in a workshop near the market.} = Son frère est apprenti soudeur dans un atelier près du marché.
\end{exemplebox}

\courssub{Professions : les métiers diplômés}

Une \eng{profession} exige des études longues, souvent à l'université, et parfois une inscription à un ordre professionnel (médecins, avocats).

\begin{center}
\begin{tabularx}{\linewidth}{|l|X|l|}
\hline
\rowcolor{popblueL} English & Prononciation (approx.) & Français \\
\hline
doctor & « DOK-teur » & médecin, docteur \\
\hline
nurse & « NEURSS » & infirmier, infirmière \\
\hline
lawyer & « LO-yeur » & avocat, juriste \\
\hline
engineer & « en-dji-NIR » & ingénieur \\
\hline
teacher & « TI-tcheur » & enseignant, professeur \\
\hline
accountant & « a-KAON-tent » & comptable \\
\hline
architect & « AR-ki-tekt » & architecte \\
\hline
pharmacist & « FAR-ma-sist » & pharmacien \\
\hline
journalist & « DJEUR-na-list » & journaliste \\
\hline
\end{tabularx}
\end{center}

\begin{attention}[Prononciations piégeuses]
Certains noms de métiers ont une prononciation qui surprend les francophones :
\begin{itemize}
\item \eng{engineer} se prononce « en-dji-NIR » (accent sur la dernière syllabe, \eng{g} se dit « dj ») ;
\item \eng{lawyer} se prononce « LO-yeur » et non « lo-yer » à la française ;
\item \eng{carpenter} se prononce « KAR-pen-teur » (accent sur la première syllabe) ;
\item \eng{mechanic} se prononce « mé-KA-nik » (le \eng{ch} se dit « k ») ;
\item \eng{entrepreneur} se prononce « on-tre-preu-NEUR » (mot d'origine française, mais accent final en anglais).
\end{itemize}
\end{attention}

\courssub{Autres emplois courants au Cameroun, lieux de travail et statuts}

\textbf{Autres emplois fréquents :} \eng{trader / shopkeeper} (commerçant), \eng{farmer} (agriculteur, cultivateur), \eng{civil servant} (fonctionnaire), \eng{driver} (chauffeur), \eng{policeman} (policier), \eng{soldier} (soldat, militaire), \eng{fisherman} (pêcheur).

\textbf{Les lieux de travail (workplaces) :}

\begin{center}
\begin{tabularx}{\linewidth}{|X|X|}
\hline
\rowcolor{popblueL} English & Français \\
\hline
workshop & atelier \\
\hline
office & bureau \\
\hline
hospital & hôpital \\
\hline
school & école \\
\hline
construction site & chantier (de construction) \\
\hline
market & marché \\
\hline
factory & usine \\
\hline
farm & ferme \\
\hline
\end{tabularx}
\end{center}

\textbf{Les statuts (people at work) :} \eng{employer} (employeur), \eng{employee} (employé, salarié), \eng{worker} (travailleur, ouvrier), \eng{apprentice} (apprenti), \eng{trainee} (stagiaire), \eng{self-employed person} (travailleur indépendant), \eng{civil servant} (fonctionnaire), \eng{unemployed person} (chômeur, sans-emploi).

\begin{attention}[Employer ou employee ?]
Ne confondez pas \eng{employer} (celui qui emploie, le patron) et \eng{employee} (celui qui est employé, le salarié). Le suffixe \cle{-er} marque celui qui fait l'action ; le suffixe \cle{-ee} marque celui qui la subit : \eng{The employer pays the employee.} = L'employeur paie l'employé.
\end{attention}

\courssub{Carte mentale : le monde du travail}

\begin{popfigure}
\begin{tikzpicture}[mindmap, grow cyclic, every node/.style={concept, text=white}, root concept/.append style={concept color=popblue, font=\bfseries\small}, level 1/.append style={level distance=3.4cm, sibling angle=90}, level 2/.append style={level distance=2.1cm, sibling angle=45, concept color=popmuted, font=\scriptsize}]
\node[root concept]{WORLD OF WORK}
  child[concept color=poporange]{ node {Trades}
    child { node {carpenter} }
    child { node {tailor} }
    child { node {mechanic} }
    child { node {plumber} }
  }
  child[concept color=popgreen]{ node {Professions}
    child { node {doctor} }
    child { node {lawyer} }
    child { node {engineer} }
    child { node {teacher} }
  }
  child[concept color=poppurple]{ node {Workplaces}
    child { node {workshop} }
    child { node {office} }
    child { node {hospital} }
    child { node {market} }
  }
  child[concept color=poppink]{ node {People at work}
    child { node {employer} }
    child { node {employee} }
    child { node {apprentice} }
    child { node {trainee} }
  };
\end{tikzpicture}
\legende{Carte mentale du champ lexical : quatre familles de mots autour du monde du travail.}
\end{popfigure}

\courssub{Trade ou profession ? Le tableau comparatif}

\begin{center}
\begin{tabularx}{\linewidth}{|X|X|}
\hline
\rowcolor{popblueL} TRADE (métier manuel) & PROFESSION (métier diplômé) \\
\hline
Learnt through practice and apprenticeship & Learnt through long studies at university \\
\hline
Manual skills (hands) & Intellectual skills (mind) \\
\hline
Often self-employed, in a workshop & Often salaried, in an office or institution \\
\hline
carpenter, mason, tailor, plumber, mechanic & doctor, lawyer, engineer, teacher, accountant \\
\hline
\eng{He learnt his trade as an apprentice.} & \eng{She entered the legal profession.} \\
\hline
\end{tabularx}
\end{center}

\begin{aretenir}[L'essentiel]
\begin{itemize}
\item \eng{job} = emploi précis (dénombrable) ; \eng{work} = travail en général (indénombrable).
\item \eng{trade} = métier manuel appris en apprentissage ; \eng{profession} = métier intellectuel exigeant des diplômes.
\item \eng{occupation} = terme général et formel ; \eng{career} = parcours professionnel sur la durée.
\item \eng{employer} (patron) $\neq$ \eng{employee} (salarié).
\end{itemize}
\end{aretenir}

\begin{exoresolu}[Classer les métiers]
Énoncé : classez les métiers suivants dans la colonne qui convient (\emph{trade} ou \emph{profession}) : \eng{nurse, plumber, architect, welder, journalist, baker, accountant, hairdresser}.
\tcblower
Solution détaillée : on se demande pour chaque métier s'il s'apprend surtout par la pratique (trade) ou par de longues études diplômantes (profession).

Trades : \eng{plumber} (plombier), \eng{welder} (soudeur), \eng{baker} (boulanger), \eng{hairdresser} (coiffeur) — ces métiers s'apprennent en atelier, souvent comme apprenti.

Professions : \eng{nurse} (infirmier), \eng{architect} (architecte), \eng{journalist} (journaliste), \eng{accountant} (comptable) — ces métiers exigent une formation diplômante.

Phrase modèle : \eng{A welder learns a trade, but an architect enters a profession.} = Un soudeur apprend un métier manuel, mais un architecte entre dans une profession.
\end{exoresolu}

\begin{savaistu}[Blue collars and white collars]
Au Royaume-Uni et aux États-Unis, on distingue traditionnellement les \eng{blue-collar workers} (« cols bleus »), c'est-à-dire les travailleurs manuels — la couleur vient des bleus de travail des ouvriers — et les \eng{white-collar workers} (« cols blancs »), les employés de bureau, qui portaient des chemises blanches. On parle même aujourd'hui de \eng{pink-collar workers} pour les métiers du service traditionnellement féminins (secrétariat, soins).
\end{savaistu}

\begin{exercice}[À vous de jouer]
\begin{enumerate}
\item Complétez avec \eng{job}, \eng{work}, \eng{trade} ou \eng{profession} : (a) \eng{My father is a mason; he learnt his \ldots\ in Bamenda.} (b) \eng{I have too much \ldots\ to do this evening.} (c) \eng{Medicine is a noble \ldots .} (d) \eng{She got a \ldots\ as a secretary in Yaoundé.}
\item Traduisez : « Mon frère est apprenti chez un électricien ; il veut apprendre un métier. »
\item Trouvez le lieu de travail de : \eng{a nurse, a trader, a tailor, a teacher}.
\end{enumerate}
\end{exercice}

\begin{corrige}[Exercice « À vous de jouer »]
1. (a) \eng{trade} ; (b) \eng{work} (indénombrable) ; (c) \eng{profession} ; (d) \eng{job}.

2. \eng{My brother is an apprentice with an electrician; he wants to learn a trade.} (Attention : « un métier » manuel se traduit ici par \eng{a trade}, pas \emph{a job}.)

3. \eng{A nurse works in a hospital; a trader works at the market (or in a shop); a tailor works in a workshop; a teacher works in a school.}
\end{corrige}

\coursec{Developing an interest : exprimer l'intérêt et la vocation}

Savoir nommer les métiers ne suffit pas : il faut aussi pouvoir dire \emph{pourquoi} un métier nous attire, \emph{depuis quand} et \emph{comment} on compte s'y former. C'est exactement ce qu'on vous demandera à l'oral comme à l'écrit (composition, dialogue, lettre).

\courssub{Un dialogue pour observer}

\begin{textebox}[After the career day]
Amina: That career day was wonderful! I am really interested in nursing now.

Bobi: Really? Since when?

Amina: Since my aunt fell ill last year. I spent two weeks with her at the hospital, and I admired the way the nurses took care of her. That is when I developed an interest in the job.

Bobi: I see. And what attracts you most?

Amina: I like helping people, and I am good at science. Nursing combines both. What about you? Have you made up your mind?

Bobi: Not yet. I am fond of engines, so I am thinking of becoming a mechanic. My father says I should take up a trade instead of dreaming about university.

Amina: Why not? A skilled mechanic earns a good living in Douala. You could train as an apprentice in your uncle's garage.

Bobi: That is exactly my plan. I want to learn the trade properly, then open my own workshop one day.

Amina: Great! Whatever we choose, let us work hard for it.
\end{textebox}

\textbf{Glossaire :}

\begin{center}
\begin{tabularx}{\linewidth}{|l|l|X|}
\hline
\rowcolor{popblueL} Mot anglais & Nature & Traduction française \\
\hline
to be interested in & expression & s'intéresser à \\
\hline
to develop an interest in & expression & développer un intérêt pour \\
\hline
to take care of & phrasal verb & prendre soin de \\
\hline
to be fond of & expression & aimer beaucoup, être passionné de \\
\hline
to make up one's mind & expression idiomatique & se décider \\
\hline
to take up a trade & expression & se lancer dans un métier (manuel) \\
\hline
to earn a (good) living & expression & gagner (bien) sa vie \\
\hline
skilled & adjectif & qualifié, compétent \\
\hline
\end{tabularx}
\end{center}

\courssub{Les structures pour exprimer l'intérêt}

\begin{propriete}[Les constructions à connaître]
Toutes ces expressions se construisent avec la \textbf{préposition} indiquée, suivie d'un nom ou d'un \textbf{verbe en -ing} :
\begin{itemize}
\item \eng{to be interested in + nom / V-ing} : \eng{I am interested in engineering.} / \eng{I am interested in repairing cars.} = Je m'intéresse à la mécanique / à réparer les voitures.
\item \eng{to be fond of + nom / V-ing} : \eng{She is fond of children.} = Elle aime beaucoup les enfants.
\item \eng{to be keen on + nom / V-ing} : \eng{He is keen on computers.} = Il est passionné d'informatique.
\item \eng{to be good at + nom / V-ing} : \eng{I am good at science.} = Je suis bon en sciences.
\item \eng{to dream of + V-ing} : \eng{She dreams of becoming a surgeon.} = Elle rêve de devenir chirurgienne.
\item \eng{to think of / to plan to / to hope to / to want to + infinitif} : \eng{I am thinking of becoming a mechanic.} / \eng{I hope to enter the medical profession.}
\end{itemize}
\end{propriete}

\begin{attention}[La préposition change tout]
Le français « s'intéresser à » se traduit par \eng{to be interested \textbf{in}}, jamais \emph{interested to}. Et après une préposition, le verbe anglais prend toujours la forme en \cle{-ing} : \eng{interested in becom\textbf{ing}}, \eng{good at draw\textbf{ing}}, \eng{fond of read\textbf{ing}}. Ne dites jamais \emph{interested in to become}.
\end{attention}

\courssub{Exprimer la vocation et le projet de formation}

\begin{center}
\begin{tabularx}{\linewidth}{|X|X|}
\hline
\rowcolor{popblueL} English & Français \\
\hline
to develop an interest in a job & développer un intérêt pour un métier \\
\hline
to choose a career & choisir une carrière \\
\hline
to make up one's mind & se décider \\
\hline
to train as a nurse / as a mechanic & se former comme infirmier / mécanicien \\
\hline
to serve an apprenticeship & faire son apprentissage \\
\hline
to learn a trade & apprendre un métier manuel \\
\hline
to take up a trade / a profession & se lancer dans un métier / une profession \\
\hline
to enter a profession & entrer dans une profession \\
\hline
to qualify as an engineer & obtenir son diplôme d'ingénieur \\
\hline
to set up one's own business & créer sa propre entreprise \\
\hline
\end{tabularx}
\end{center}

\begin{exemplebox}[La vocation en contexte]
\eng{I developed an interest in teaching when I was in Form Five.} = J'ai développé un intérêt pour l'enseignement quand j'étais en Seconde.

\eng{After serving a five-year apprenticeship, he qualified as a carpenter.} = Après un apprentissage de cinq ans, il est devenu menuisier diplômé.

\eng{She has made up her mind: she will train as a nurse in Bamenda.} = Elle s'est décidée : elle se formera comme infirmière à Bamenda.

\eng{He took up his father's trade and now runs the family workshop.} = Il a repris le métier de son père et dirige maintenant l'atelier familial.
\end{exemplebox}

\coursec{Sharing an interest : échanger sur les métiers}

Parler de son projet professionnel, c'est aussi \emph{écouter} celui des autres, poser des questions, donner un avis ou un conseil. Voici les outils de conversation indispensables.

\courssub{Poser des questions sur les métiers}

\begin{center}
\begin{tabularx}{\linewidth}{|X|X|}
\hline
\rowcolor{popblueL} English & Français \\
\hline
What do you do (for a living)? & Que fais-tu dans la vie ? \\
\hline
What does your father do? & Que fait ton père (comme métier) ? \\
\hline
What would you like to be / to become? & Que voudrais-tu devenir ? \\
\hline
Why are you interested in this job? & Pourquoi ce métier t'intéresse-t-il ? \\
\hline
What attracts you most in this profession? & Qu'est-ce qui t'attire le plus dans cette profession ? \\
\hline
How long does the training take? & Combien de temps dure la formation ? \\
\hline
Is it a well-paid job? & Est-ce un métier bien payé ? \\
\hline
\end{tabularx}
\end{center}

\begin{attention}[What do you do?]
Pour demander le métier de quelqu'un, l'anglais dit \eng{What do you do?} (littéralement « que fais-tu ? »), jamais \emph{What is your job?} — grammaticalement correct mais peu naturel — et surtout pas \emph{What is your profession?} dans une conversation courante. Réponse : \eng{I am a teacher.} et non \emph{I am teacher} (l'article \eng{a} est obligatoire devant un métier au singulier).
\end{attention}

\courssub{Réagir, conseiller, encourager}

\begin{center}
\begin{tabularx}{\linewidth}{|X|X|}
\hline
\rowcolor{popblueL} English & Français \\
\hline
That sounds interesting! & Cela a l'air intéressant ! \\
\hline
I see. / Really? & Je vois. / Vraiment ? \\
\hline
Why don't you train as a nurse? & Pourquoi ne te formes-tu pas comme infirmier ? \\
\hline
You should learn a trade. & Tu devrais apprendre un métier. \\
\hline
You could work in your uncle's garage. & Tu pourrais travailler dans le garage de ton oncle. \\
\hline
Whatever you choose, work hard for it. & Quoi que tu choisisses, travaille dur pour y arriver. \\
\hline
I wish you success. & Je te souhaite de réussir. \\
\hline
\end{tabularx}
\end{center}

\begin{experience}[Jeu de rôle : the career interview]
Par groupes de deux :
\begin{enumerate}
\item L'élève A est journaliste pour le journal du lycée ; l'élève B est un professionnel invité (menuisier, infirmière, commerçante, ingénieur\ldots).
\item A pose au moins cinq questions avec les structures vues (\eng{What do you do? Why did you choose\ldots? How long\ldots?}).
\item B répond en utilisant \eng{I developed an interest in\ldots}, \eng{I trained as\ldots}, \eng{I earn my living by\ldots}.
\item Les rôles sont ensuite échangés. Le professeur veille à l'emploi de l'article \eng{a} devant les métiers et des prépositions correctes.
\end{enumerate}
\end{experience}

\coursec{Collocations et expressions du monde du travail}

Une \cle{collocation} est une association habituelle de mots : en anglais, on ne « fait » pas un métier, on ne « gagne » pas un salaire avec \eng{win}. Apprendre les mots par paires évite les calques du français.

\courssub{Les collocations essentielles}

\begin{center}
\begin{tabularx}{\linewidth}{|X|X|}
\hline
\rowcolor{popblueL} Collocation anglaise & Équivalent français \\
\hline
to look for a job / for work & chercher un emploi / du travail \\
\hline
to find a job & trouver un emploi \\
\hline
to lose one's job & perdre son emploi \\
\hline
to apply for a job & poser sa candidature, postuler \\
\hline
to get a job as a secretary & obtenir un poste de secrétaire \\
\hline
to do a job / to do one's work & faire un travail / son travail \\
\hline
to work hard & travailler dur \\
\hline
to earn a living / one's living & gagner sa vie \\
\hline
to earn a good salary & toucher un bon salaire \\
\hline
to be employed / unemployed & être employé / au chômage \\
\hline
to be self-employed & travailler à son compte \\
\hline
to retire (from one's job) & prendre sa retraite \\
\hline
a full-time / part-time job & un emploi à temps plein / partiel \\
\hline
a well-paid / badly-paid job & un emploi bien / mal payé \\
\hline
\end{tabularx}
\end{center}

\begin{attention}[Do ou make ? Win ou earn ?]
\begin{itemize}
\item Avec \eng{work} et \eng{job}, on utilise \eng{do}, jamais \eng{make} : \eng{He does his work well.} (Il fait bien son travail.) Mais on dit \eng{to make money} (gagner de l'argent, se faire de l'argent) et \eng{to make a decision}.
\item « Gagner sa vie » se dit \eng{to \textbf{earn} one's living}, jamais \emph{to win one's living} : \eng{win} s'emploie pour un match, un prix, une compétition (\eng{to win a match, to win a prize}).
\item « Chercher du travail » : \eng{to \textbf{look for} work} — la préposition \eng{for} est obligatoire, contrairement au français.
\end{itemize}
\end{attention}

\courssub{Expressions idiomatiques}

\begin{center}
\begin{tabularx}{\linewidth}{|X|X|}
\hline
\rowcolor{popblueL} Expression idiomatique & Sens \\
\hline
to learn the ropes & apprendre les ficelles du métier \\
\hline
to be the breadwinner & être le soutien de famille (celui qui fait vivre le foyer) \\
\hline
to work one's way up & gravir les échelons par son travail \\
\hline
to be out of work & être sans emploi, au chômage \\
\hline
to get the sack / to be sacked & se faire licencier (« renvoyer ») \\
\hline
a dead-end job & un emploi sans avenir, sans promotion possible \\
\hline
\end{tabularx}
\end{center}

\begin{exemplebox}[Idiomatismes en contexte]
\eng{As the eldest son, he became the breadwinner after his father's death.} = Fils aîné, il est devenu le soutien de famille après la mort de son père.

\eng{She started as a cleaner and worked her way up to manager.} = Elle a commencé comme femme de ménage et a gravi les échelons jusqu'au poste de directrice.

\eng{Many young graduates are out of work in our towns.} = Beaucoup de jeunes diplômés sont sans emploi dans nos villes.
\end{exemplebox}

\coursec{Faux amis et pièges des francophones}

\begin{attention}[Les faux amis du monde du travail]
\begin{itemize}
\item \eng{a \textbf{profession}} $\neq$ toujours « profession » au sens large : en anglais, le mot est réservé aux métiers intellectuels diplômés (médecin, avocat). Pour « métier » en général, dites \eng{job} ou \eng{occupation}.
\item \eng{a \textbf{career}} $\neq$ « carrière » au sens de piste : c'est le parcours professionnel. Attention à ne pas confondre avec \eng{a carrier} (transporteur).
\item \eng{a \textbf{formation}} n'existe pas : « une formation » se dit \eng{training} : \eng{a training centre} = un centre de formation ; \eng{vocational training} = la formation professionnelle.
\item \eng{a \textbf{stage}} n'existe pas pour « stage » : dites \eng{an internship} ou \eng{a work placement} ; le stagiaire est \eng{a trainee} ou \eng{an intern}.
\item \eng{a \textbf{diploma}} = un diplôme, mais « être diplômé » se dit \eng{to be qualified} ou \eng{to have a degree} (diplôme universitaire), pas \emph{to have a diploma} dans tous les contextes.
\item \eng{a \textbf{concours}} n'existe pas : « passer un concours » se dit \eng{to sit (for) a competitive examination} ou \eng{to take an entrance exam}.
\item \eng{a \textbf{patron}} : le patron se dit \eng{the boss} ou \eng{the employer} ; \eng{a patron} en anglais désigne un mécène ou un client habitué.
\item \eng{a \textbf{salary}} = un salaire (mensuel, souvent des employés qualifiés) ; \eng{wages} (pluriel) = le salaire (à l'heure ou à la semaine, souvent des ouvriers). Ne dites pas \emph{a wage} pour « un salaire » mensuel.
\end{itemize}
\end{attention}

\begin{center}
\begin{tabularx}{\linewidth}{|X|X|X|}
\hline
\rowcolor{popblueL} Français & English correct & Piège à éviter \\
\hline
chercher du travail & to look for work & \emph{to search work} \\
\hline
gagner sa vie & to earn one's living & \emph{to win one's life} \\
\hline
apprendre un métier & to learn a trade & \emph{to learn a job} \\
\hline
faire un apprentissage & to serve an apprenticeship & \emph{to do an apprenticeship} (courant mais moins correct) \\
\hline
un centre de formation & a training centre & \emph{a formation centre} \\
\hline
être au chômage & to be unemployed / out of work & \emph{to be in chomage} \\
\hline
un emploi bien payé & a well-paid job & \emph{a well-payed job} (orthographe : \eng{paid}) \\
\hline
\end{tabularx}
\end{center}

\begin{attention}[Orthographe : pay -- paid -- paid]
Le verbe \eng{to pay} est irrégulier dans son orthographe : prétérit et participe passé \eng{paid} (et non \emph{payed}). D'où \eng{a well-paid job}, \eng{He paid the worker yesterday.}
\end{attention}

\coursec{Mots de la même famille et réemploi en contexte}

\courssub{Familles de mots : employer, former, qualifier}

Construire la famille d'un mot, c'est multiplier son vocabulaire sans effort. Observez les suffixes : \cle{-er/-or} (celui qui fait), \cle{-ee} (celui qui subit), \cle{-ment}, \cle{-tion}, \cle{-ing} (noms d'action), \cle{-ed/-able} (adjectifs).

\begin{center}
\begin{tabular}{|l|l|l|}
\hline
\rowcolor{popblueL} Verbe & Nom (personne / action) & Adjectif \\
\hline
to employ & employer / employee ; employment & employed / unemployed \\
\hline
to train & trainer / trainee ; training & trained / untrained \\
\hline
to qualify & qualification & qualified / unqualified \\
\hline
to work & worker ; work & working / hardworking \\
\hline
to teach & teacher ; teaching & -- \\
\hline
to farm & farmer ; farming & -- \\
\hline
to trade & trader ; trade & -- \\
\hline
to apply & applicant ; application & -- \\
\hline
to retire & retirement & retired \\
\hline
\end{tabular}
\end{center}

\begin{exemplebox}[Une famille, trois phrases]
\eng{The factory employs two hundred workers.} (verbe) = L'usine emploie deux cents ouvriers.

\eng{Employment is a serious problem for young graduates.} (nom) = L'emploi est un problème grave pour les jeunes diplômés.

\eng{Many trained nurses are unemployed in our region.} (adjectifs) = Beaucoup d'infirmières diplômées sont au chômage dans notre région.
\end{exemplebox}

\begin{attention}[Trainee, pas « stagiaire » français]
\eng{A trainee} (« TRÉ-ni ») est un stagiaire en formation professionnelle ; le double \eng{e} final marque la personne qui subit l'action, comme \eng{employee}. Ne l'écrivez pas \emph{traine} et ne le confondez pas avec \eng{trainer} (le formateur).
\end{attention}

\courssub{Réemploi guidé : un court texte à trous}

\begin{exoresolu}[Texte à trous : My sister's choice]
Énoncé : complétez avec un mot de la famille indiquée entre parenthèses.

\eng{My sister has always been interested in nursing. After her Advanced Level, she decided to \ldots\ (training) as a nurse in a \ldots\ (train) centre in Bamenda. The \ldots\ (train) lasted three years. She is now a \ldots\ (qualify) nurse, but like many young graduates, she is still looking for \ldots\ (employ). She has \ldots\ (application) for a job at the regional hospital and hopes to be \ldots\ (employer) soon.}
\tcblower
Solution détaillée :

\eng{My sister has always been interested in nursing. After her Advanced Level, she decided to \textbf{train} as a nurse in a \textbf{training} centre in Bamenda. The \textbf{training} lasted three years. She is now a \textbf{qualified} nurse, but like many young graduates, she is still looking for \textbf{employment}. She has \textbf{applied} for a job at the regional hospital and hopes to be \textbf{employed} soon.}

Commentaires : après \eng{decided to}, il faut l'infinitif \eng{train} ; devant \eng{centre}, le nom composé exige \eng{training} ; devant le nom \eng{nurse}, l'adjectif \eng{qualified} ; après \eng{looking for}, le nom \eng{employment} ; \eng{has applied} (present perfect de \eng{to apply}, participe passé régulier avec \eng{y} $\rightarrow$ \eng{ied}) ; \eng{to be employed} (infinitif passif : « être employée »).
\end{exoresolu}

\begin{methode}[Réutiliser le lexique dans une composition (Probatoire / Bac)]
Pour un sujet du type \emph{“Write about the job you would like to do in the future”} :
\begin{enumerate}
\item \textbf{Nommer} le métier avec l'article : \eng{I would like to become a nurse.} (jamais \emph{I would like to become nurse}).
\item \textbf{Classer} : trade ou profession ? \eng{Nursing is a profession that requires long training.}
\item \textbf{Expliquer l'origine de l'intérêt} : \eng{I developed an interest in this job when\ldots}
\item \textbf{Donner des raisons} avec les bonnes prépositions : \eng{I am good at science; I am fond of helping people.}
\item \textbf{Décrire le projet de formation} : \eng{I hope to train as a nurse in\ldots; after qualifying, I would like to work in\ldots}
\item \textbf{Conclure} par une phrase de motivation : \eng{Whatever the difficulties, I will work hard to earn my living and serve my country.}
\end{enumerate}
\end{methode}

\begin{exercice}[Entraînement final]
\begin{enumerate}
\item Traduisez : (a) « Elle s'intéresse à la médecine depuis son enfance. » (b) « Mon oncle gagne bien sa vie comme soudeur. » (c) « Après son apprentissage, il a ouvert son propre atelier. »
\item Complétez avec la bonne préposition : (a) \eng{She is fond \ldots\ children.} (b) \eng{He is good \ldots\ mathematics.} (c) \eng{I am interested \ldots\ becoming an engineer.} (d) \eng{She applied \ldots\ a job at the hospital.}
\item Trouvez l'intrus et justifiez : \eng{to earn a living / to win a living / to make money / to earn a salary}.
\item Rédigez un court paragraphe (5 à 6 lignes) : \eng{The job I would like to do and why.}
\end{enumerate}
\end{exercice}

\begin{corrige}[Exercice « Entraînement final »]
1. (a) \eng{She has been interested in medicine since her childhood.} (b) \eng{My uncle earns a good living as a welder.} (c) \eng{After his apprenticeship, he opened his own workshop.}

2. (a) \eng{of} ; (b) \eng{at} ; (c) \eng{in} ; (d) \eng{for}.

3. L'intrus est \eng{to win a living} : cette expression n'existe pas ; « gagner sa vie » se dit \eng{to earn one's living}. \eng{Win} s'emploie pour un match ou un prix.

4. Modèle de production : \eng{The job I would like to do is teaching. Teaching is a noble profession, and I developed an interest in it when I was in Form Three. I am good at English and I am fond of explaining things to my classmates. After my Advanced Level, I hope to train as a teacher in a training college. When I qualify, I would like to teach in a government high school in my region. I know the salary is not very high, but I will be happy to earn my living by serving the youth of my country.}
\end{corrige}

\coursec{Developing an interest : exprimer l'intérêt et la vocation}

Dans la section précédente, nous avons dressé la carte du monde des métiers : \eng{trades}, \eng{jobs} et \eng{professions}. Mais un métier ne se choisit pas au hasard : il naît presque toujours d'un \cle{intérêt}, d'une curiosité, parfois d'une véritable \cle{vocation}. Comment dire en anglais qu'on « s'intéresse à » la mécanique, qu'on « rêve de » devenir ingénieur, qu'on « a décidé de » se former à la couture ? C'est tout l'objet de cette section : apprendre à exprimer la naissance, le développement et l'affirmation d'un intérêt professionnel, avec des structures grammaticalement exactes et des nuances de registre adaptées.

\courssub{Observation : un texte pour découvrir le vocabulaire en contexte}

Lisons d'abord le parcours de Ngong, un jeune élève de Bamenda. Soulignez mentalement toutes les expressions qui servent à parler de son intérêt pour la mécanique.

\begin{textebox}[Ngong's calling — a text to read]
Since childhood, Ngong has been keen on mechanics. While other boys in his quarter played football after school, he spent his holidays in his father's garage in Bamenda, watching the mechanics repair old cars and motorcycles. Little by little, he developed a real passion for engines. He became curious about every strange noise a car could make, and he took a growing interest in the tools hanging on the workshop wall.

At first, his mother was not pleased: she dreamed of her son becoming a doctor. But Ngong was determined. “I am passionate about machines,” he told her. “I don't want to sit in an office. I want to work with my hands.” When he turned fifteen, he announced that he had decided to train as a mechanic after his GCE Ordinary Level. His father, a skilled and respected craftsman, was proud: “My son is gifted and motivated. He will go far.”

Ngong's ambition is clear. He aspires to become one of the best mechanics in the North-West Region, and one day he dreams of opening his own garage. “Mechanics is not just a job for me,” he says with a smile. “It is a calling.”
\end{textebox}

\textbf{Glossaire}

\begin{center}
\begin{tabularx}{\linewidth}{|l|l|X|}
\hline
\rowcolor{popblueL} Mot anglais & Nature & Traduction française \\
\hline
to be keen on & expression & être passionné par, aimer beaucoup \\
\hline
quarter & nom & quartier (habitat) \\
\hline
garage & nom & garage (atelier de réparation) \\
\hline
to develop a passion for & expression & développer une passion pour \\
\hline
curious about & adjectif + prép. & curieux de \\
\hline
to train as & verbe + prép. & se former comme, suivre une formation de \\
\hline
craftsman & nom & artisan (plur. \eng{craftsmen}) \\
\hline
gifted & adjectif & doué \\
\hline
to aspire to become & expression & aspirer à devenir \\
\hline
calling & nom & vocation \\
\hline
\end{tabularx}
\end{center}

\textbf{Questions de compréhension}

\begin{enumerate}
\item What has Ngong been keen on since childhood?
\item Where did he spend his holidays?
\item What did his mother want him to become?
\item What has Ngong decided to do after his GCE Ordinary Level?
\item What is his long-term dream?
\end{enumerate}

\textit{Réponses attendues : 1. Mechanics. 2. In his father's garage in Bamenda. 3. A doctor. 4. To train as a mechanic. 5. Opening his own garage.}

\courssub{Exprimer l'intérêt : les verbes et structures de base}

Observons les phrases tirées du texte et de la vie quotidienne :

\eng{Ngong has been keen on mechanics since childhood.} « Ngong est passionné de mécanique depuis l'enfance. »

\eng{I am interested in nursing.} « Je m'intéresse au métier d'infirmier. »

\eng{She is fond of working with children.} « Elle aime beaucoup travailler avec les enfants. »

\eng{He dreams of becoming an engineer.} « Il rêve de devenir ingénieur. »

On remarque une constante : après l'expression d'intérêt, on trouve une \cle{préposition} (\eng{in}, \eng{on}, \eng{of}, \eng{about}, \eng{to}), suivie d'un \cle{nom} ou d'un \cle{verbe en -ing}.

\begin{propriete}[La règle d'or : préposition + nom ou verbe en -ing]
Après les expressions d'intérêt construites avec une préposition — \eng{interested in}, \eng{keen on}, \eng{fond of}, \eng{passionate about}, \eng{attracted to}, \eng{dream of}, \eng{curious about} — la préposition est \textbf{toujours} suivie :
\begin{itemize}
\item d'un \textbf{nom} : \eng{She is interested in medicine.} « Elle s'intéresse à la médecine. »
\item ou d'un \textbf{verbe en -ing} (gérondif) : \eng{I am interested in becoming a nurse.} « Je m'intéresse à l'idée de devenir infirmier. »
\end{itemize}
On ne dit jamais \eng{interested to become}. En anglais, une préposition ne peut jamais être suivie d'un infinitif : c'est une différence majeure avec le français (« je rêve \emph{de devenir} » → \eng{I dream \emph{of becoming}}).
\end{propriete}

\begin{exemplebox}[Les structures d'intérêt, du plus faible au plus fort]
\begin{itemize}
\item \eng{I am curious about carpentry.} « La menuiserie m'intrigue. » (curiosité légère)
\item \eng{Manka is interested in computer science.} « Manka s'intéresse à l'informatique. »
\item \eng{My brother is keen on repairing radios.} « Mon frère aime beaucoup réparer des radios. »
\item \eng{She is fond of designing clothes.} « Elle aime beaucoup dessiner des vêtements. »
\item \eng{Many students are attracted to journalism.} « Beaucoup d'élèves sont attirés par le journalisme. »
\item \eng{Etienne is passionate about cooking.} « Étienne est passionné de cuisine. »
\item \eng{Amina dreams of becoming a pilot.} « Amina rêve de devenir pilote. »
\end{itemize}
\end{exemplebox}

\begin{attention}[Piège n° 1 : « s'intéresser à » ne se traduit pas mot à mot]
Le verbe français « s'intéresser à » tente souvent les élèves d'écrire \eng{to interest in}, ce qui est \textbf{faux}. Les deux seules structures correctes sont :
\begin{itemize}
\item \eng{to be interested in} : \eng{I am interested in tailoring.} « Je m'intéresse à la couture. »
\item \eng{to take an interest in} : \eng{She took an interest in farming.} « Elle s'est mise à s'intéresser à l'agriculture. »
\end{itemize}
Attention aussi à la confusion entre \eng{interested} (celui qui ressent l'intérêt : « intéressé ») et \eng{interesting} (ce qui provoque l'intérêt : « intéressant »). \eng{I am interested in this interesting job.} « Ce métier intéressant m'intéresse. »
\end{attention}

\courssub{Développer un intérêt : la naissance d'une vocation}

Un intérêt n'apparaît pas toujours d'un coup : il se \cle{développe}, il \cle{grandit}. L'anglais dispose d'une série de verbes dynamiques pour raconter cette évolution.

\begin{definition}[Verbes du développement de l'intérêt]
\begin{itemize}
\item \eng{to develop an interest in} : développer un intérêt pour. \eng{She developed an interest in tailoring when she was twelve.} « Elle a développé un intérêt pour la couture à douze ans. »
\item \eng{to grow fond of} : se prendre d'affection pour, finir par aimer. \eng{He grew fond of carpentry in his uncle's workshop.} « Il a fini par aimer la menuiserie dans l'atelier de son oncle. »
\item \eng{to become curious about} : devenir curieux de. \eng{The students became curious about electricity.} « Les élèves sont devenus curieux de l'électricité. »
\item \eng{to discover a passion for} : découvrir une passion pour. \eng{During the holidays, I discovered a passion for photography.} « Pendant les vacances, j'ai découvert une passion pour la photographie. »
\item \eng{to take an interest in} : commencer à s'intéresser à. \eng{My father took an interest in my school results.} « Mon père s'est mis à s'intéresser à mes résultats scolaires. »
\end{itemize}
\end{definition}

\begin{popfigure}
\begin{center}
\begin{tikzpicture}[font=\small]
\draw[-{Stealth[length=3mm]}, very thick, popblue] (0,0) -- (12.6,0);
\foreach \x/\txt/\col in {0.4/{to be curious\\about}/popgreen, 3.0/{to be interested\\in}/popblue, 5.8/{to be keen\\on}/poppurple, 8.6/{to be passionate\\about}/poporange, 11.6/{to choose as\\a career}/popink}{
  \fill[\col] (\x,0) circle (0.09);
  \node[above, align=center, text=\col] at (\x,0.15) {\txt};
}
\node[below, align=center, text=popmuted] at (0.4,-0.15) {curiosité};
\node[below, align=center, text=popmuted] at (3.0,-0.15) {intérêt};
\node[below, align=center, text=popmuted] at (5.8,-0.15) {enthousiasme};
\node[below, align=center, text=popmuted] at (8.6,-0.15) {passion};
\node[below, align=center, text=popmuted] at (11.6,-0.15) {vocation};
\end{tikzpicture}
\end{center}
\legende{L'échelle de l'intérêt : de la simple curiosité au choix de carrière.}
\end{popfigure}

\courssub{Mots de la même famille : dérivation et prononciation}

Le lexique de l'intérêt et de la vocation se construit autour de quelques racines. Maîtriser les dérivés (noms, verbes, adjectifs, adverbes) enrichit l'expression écrite et orale.

\begin{center}
\begin{tabular}{|l|l|l|l|l|}
\hline
\rowcolor{popblueL} Verbe & Nom (chose/qualité) & Nom (personne) & Adjectif & Adverbe \\
\hline
\eng{to interest} & \eng{interest} & — & \eng{interested} / \eng{interesting} & \eng{interestingly} \\
\hline
\eng{to passion} (rare) & \eng{passion} & — & \eng{passionate} & \eng{passionately} \\
\hline
— & \eng{vocation} / \eng{calling} & — & \eng{vocational} & — \\
\hline
\eng{to ambition} (rare) & \eng{ambition} & — & \eng{ambitious} & \eng{ambitiously} \\
\hline
\eng{to aspire} & \eng{aspiration} & \eng{aspirant} & \eng{aspiring} & — \\
\hline
\eng{to motivate} & \eng{motivation} & — & \eng{motivated} / \eng{motivating} & \eng{motivatedly} \\
\hline
\eng{to determine} & \eng{determination} & — & \eng{determined} & \eng{determinedly} \\
\hline
\eng{to skill} (non usité) & \eng{skill} & \eng{craftsman} & \eng{skilled} / \eng{skilful} & \eng{skilfully} \\
\hline
\end{tabular}
\end{center}

\begin{attention}[Prononciation et pièges d'orthographe]
\begin{itemize}
\item \eng{Vocation} : se prononce « vo-kè-chène » (son \eng{/ʃ/} écrit \eng{ti}), pas « vo-ka-sion ».
\item \eng{Aspiration} : « as-pi-rè-chène » (son \eng{/ʃ/}).
\item \eng{Ambitious} : « am-bi-chi-eusse » (son \eng{/ʃ/} écrit \eng{ti}), accent sur \eng{bi}.
\item \eng{Skilful} (orthographe britannique : un seul \eng{l}) : « skille-fulle ».
\item \eng{Determined} : « di-tèr-mi-neud » (le \eng{e} final se prononce \eng{/ɪd/} après \eng{t/d}).
\end{itemize}
\end{attention}

\courssub{Décider et s'engager : les verbes de la décision}

Quand l'intérêt devient un projet, on passe aux \cle{verbes de décision et d'engagement}. Observez :

\eng{Ngong has decided to train as a mechanic.} « Ngong a décidé de se former comme mécanicien. »

\eng{She chose a career in teaching.} « Elle a choisi une carrière dans l'enseignement. »

\begin{propriete}[Deux constructions à ne pas confondre]
\begin{itemize}
\item \eng{to decide to} + \textbf{infinitif} : \eng{He decided to become an electrician.} « Il a décidé de devenir électricien. » (ici, pas de préposition, donc infinitif)
\item \eng{to dream of} / \eng{to think of} + \textbf{-ing} : \eng{He dreams of becoming an engineer.} « Il rêve de devenir ingénieur. »
\end{itemize}
Autres verbes d'engagement utiles :
\begin{itemize}
\item \eng{to take up a trade} : se lancer dans un métier manuel. \eng{After Form Five, she took up hairdressing.} « Après la classe de seconde, elle s'est lancée dans la coiffure. »
\item \eng{to go into} + domaine : entrer dans, se lancer dans. \eng{He wants to go into business.} « Il veut se lancer dans les affaires. »
\item \eng{to train as} + métier : se former comme. \eng{She is training as a nurse at the hospital in Buea.} « Elle se forme comme infirmière à l'hôpital de Buea. »
\item \eng{to aspire to become} : aspirer à devenir (soutenu). \eng{He aspires to become an engineer.} « Il aspire à devenir ingénieur. »
\end{itemize}
\end{propriete}

\courssub{Noms, adjectifs et questions : le lexique complet}

\begin{center}
\begin{tabularx}{\linewidth}{|X|X|X|}
\hline
\rowcolor{popblueL} English & Français & Exemple \\
\hline
interest & intérêt & \eng{Her interest in law is strong.} \\
\hline
passion & passion & \eng{He has a passion for engines.} \\
\hline
vocation / calling & vocation & \eng{Teaching is her calling.} \\
\hline
ambition & ambition & \eng{Her ambition is to become a surgeon.} \\
\hline
dream & rêve & \eng{His dream is to open a bakery.} \\
\hline
aspiration & aspiration & \eng{What are your career aspirations?} \\
\hline
motivation & motivation & \eng{His motivation is admirable.} \\
\hline
determination & détermination & \eng{She showed great determination.} \\
\hline
\end{tabularx}
\end{center}

\begin{center}
\begin{tabularx}{\linewidth}{|X|X|X|}
\hline
\rowcolor{popgreenL} Adjective & Français & Exemple traduit \\
\hline
interested & intéressé & \eng{I am interested in farming.} « L'agriculture m'intéresse. » \\
\hline
keen & enthousiaste & \eng{She is keen on computers.} « Elle adore l'informatique. » \\
\hline
passionate & passionné & \eng{He is passionate about music.} « Il est passionné de musique. » \\
\hline
ambitious & ambitieux & \eng{An ambitious student works hard.} « Un élève ambitieux travaille dur. » \\
\hline
determined & déterminé & \eng{Ngong was determined.} « Ngong était déterminé. » \\
\hline
motivated & motivé & \eng{Motivated learners succeed.} « Les apprenants motivés réussissent. » \\
\hline
gifted / talented & doué & \eng{She is a gifted tailor.} « C'est une couturière douée. » \\
\hline
skilled & qualifié, habile & \eng{A skilled craftsman} « un artisan qualifié » \\
\hline
skilful & habile (de ses mains) & \eng{A skilful driver} « un conducteur habile » \\
\hline
aspiring & aspirant, qui aspire à & \eng{An aspiring journalist} « un journaliste en herbe » \\
\hline
\end{tabularx}
\end{center}

\begin{attention}[Piège n° 2 : les faux amis du lexique de la vocation]
\begin{itemize}
\item \eng{ambitious} = ambitieux (correct), mais « une ambition » professionnelle se dit bien \eng{an ambition}, jamais \eng{a pretension} (\eng{pretension} = prétention, souvent péjoratif).
\item « Être doué » se dit \eng{gifted} ou \eng{talented} ; \eng{gifted} ne veut pas dire « offert en cadeau » dans ce contexte (participe passé de \eng{to gift}).
\item \eng{skilled} (qualifié, qui a un savoir-faire technique, ex. \eng{skilled labour}) ne se confond pas avec \eng{skilful} (habile, adroit dans l'exécution, orthographe UK : un \eng{l}).
\item « La vocation » : \eng{vocation} et \eng{calling} sont corrects ; attention, \eng{a calling} ne signifie pas « un appel téléphonique » (qui se dit \eng{a phone call}).
\item \eng{to pretend} ne signifie pas « prétendre » au sens de « revendiquer » ou « aspirer à », mais « faire semblant ». Pour « prétendre à un poste », on utilise \eng{to apply for} ou \eng{to claim}.
\end{itemize}
\end{attention}

\textbf{Interroger quelqu'un sur son intérêt et ses projets}

\begin{exemplebox}[Questions utiles à l'oral]
\begin{itemize}
\item \eng{What do you want to be?} « Que veux-tu faire plus tard ? »
\item \eng{What are you interested in?} « Qu'est-ce qui t'intéresse ? »
\item \eng{Why are you attracted to this job?} « Pourquoi ce métier t'attire-t-il ? »
\item \eng{Have you chosen a career yet?} « As-tu déjà choisi une carrière ? »
\item \eng{When did you become interested in this trade?} « Quand t'es-tu intéressé à ce métier ? »
\item \eng{What are your ambitions?} « Quelles sont tes ambitions ? »
\item \eng{What are your career aspirations?} « Quelles sont tes aspirations professionnelles ? » (plus formel)
\end{itemize}
\end{exemplebox}

\courssub{Nuances de registre : dire la même chose autrement}

Selon la situation — discussion entre amis, entretien, rédaction scolaire — on choisit un registre différent.

\begin{center}
\begin{tabularx}{\linewidth}{|l|X|X|}
\hline
\rowcolor{poppurpleL} Registre & English & Français \\
\hline
Familier & \eng{I'm crazy about football.} & Je suis fou de football. \\
\hline
Courant & \eng{I want to be a nurse.} & Je veux devenir infirmière. \\
\hline
Courant + & \eng{I'm keen on mechanics.} & La mécanique me passionne. \\
\hline
Soutenu & \eng{I aspire to become a lawyer.} & J'aspire à devenir avocat. \\
\hline
Soutenu & \eng{I have developed a deep interest in teaching.} & J'ai développé un vif intérêt pour l'enseignement. \\
\hline
\end{tabularx}
\end{center}

Dans une rédaction ou un entretien officiel, préférez \eng{I aspire to become...} ou \eng{I am passionate about...} ; réservez \eng{I'm crazy about...} aux conversations entre amis.

\courssub{Méthode : présenter son projet professionnel en trois étapes}

\begin{methode}[How to talk about your career plan]
\begin{enumerate}
\item \textbf{Annoncer le métier visé} : \eng{I want to be a mechanic.} / \eng{I aspire to become a nurse.}
\item \textbf{Expliquer l'origine de l'intérêt} : \eng{I became interested in mechanics when I helped my father in his garage.} / \eng{I discovered a passion for nursing when my grandmother was ill.}
\item \textbf{Annoncer la démarche concrète} : \eng{So I am going to train as a mechanic after my GCE.} / \eng{That is why I have decided to go into teaching.}
\end{enumerate}
Modèle complet : \eng{I want to be an electrician. I became interested in electricity when I watched my uncle repair the wiring in our house in Douala. So I have decided to train as an electrician after Form Five.} « Je veux devenir électricien. Je me suis intéressé à l'électricité en regardant mon oncle réparer l'installation de notre maison à Douala. J'ai donc décidé de me former comme électricien après la seconde. »
\end{methode}

\courssub{Exercice résolu et entraînement}

\begin{exoresolu}[Complétez avec la bonne forme du verbe]
Énoncé — Complétez chaque phrase avec le verbe entre parenthèses à la forme correcte (infinitif ou -ing) :
\begin{enumerate}
\item I am interested in \dots\ (become) a doctor.
\item She dreams of \dots\ (open) her own restaurant.
\item He decided \dots\ (train) as a plumber.
\item They are keen on \dots\ (learn) computer skills.
\item My sister aspires to \dots\ (become) a journalist.
\end{enumerate}
\tcblower
Solution détaillée :
\begin{enumerate}
\item \eng{becoming} — après la préposition \eng{in}, verbe en -ing : \eng{I am interested in becoming a doctor.}
\item \eng{opening} — après la préposition \eng{of} : \eng{She dreams of opening her own restaurant.}
\item \eng{to train} — \eng{to decide} est suivi de l'infinitif : \eng{He decided to train as a plumber.}
\item \eng{learning} — après la préposition \eng{on} : \eng{They are keen on learning computer skills.}
\item \eng{become} — \eng{to aspire to} est suivi de l'infinitif : \eng{My sister aspires to become a journalist.}
\end{enumerate}
Astuce : repérez d'abord s'il y a une \textbf{préposition} avant le verbe. Préposition → -ing ; pas de préposition → infinitif avec \eng{to}.
\end{exoresolu}

\begin{exercice}[À vous de jouer]
\begin{enumerate}
\item Traduisez en anglais : a) « Elle s'intéresse à la coiffure. » b) « Il rêve de devenir pilote. » c) « Nous avons décidé de nous former comme soudeurs. » d) « Elle a développé une passion pour la cuisine. »
\item Répondez en phrases complètes : \eng{What do you want to be? Why are you attracted to this job?}
\item Recopiez en corrigeant l'erreur : \eng{I am interesting to become a teacher.} / \eng{She dreams to be a nurse.}
\end{enumerate}
\end{exercice}

\begin{corrige}[Exercice « À vous de jouer »]
\begin{enumerate}
\item a) \eng{She is interested in hairdressing.} b) \eng{He dreams of becoming a pilot.} c) \eng{We have decided to train as welders.} d) \eng{She has developed a passion for cooking.}
\item Réponse libre, par exemple : \eng{I want to be an accountant because I am keen on mathematics and I am attracted to office work.}
\item \eng{I am interested in becoming a teacher.} (on ressent l'intérêt : \eng{interested} + préposition \eng{in} + -ing) ; \eng{She dreams of being a nurse.} ou \eng{She dreams of becoming a nurse.} (\eng{dream of} + -ing).
\end{enumerate}
\end{corrige}

\begin{aretenir}[L'essentiel de la section]
\begin{itemize}
\item Préposition + nom ou verbe en \textbf{-ing} : \eng{interested in}, \eng{keen on}, \eng{fond of}, \eng{passionate about}, \eng{dream of}, \eng{curious about}, \eng{attracted to}.
\item « S'intéresser à » = \eng{to be interested in} ou \eng{to take an interest in} (jamais \eng{to interest in}).
\item Développement : \eng{to develop an interest in}, \eng{to discover a passion for}, \eng{to grow fond of}, \eng{to become curious about}.
\item Décision : \eng{to decide to} + infinitif ; \eng{to train as}, \eng{to go into}, \eng{to take up a trade}, \eng{to aspire to become}.
\item Dérivation : \eng{interest} → \eng{interested/interesting} ; \eng{passion} → \eng{passionate} ; \eng{ambition} → \eng{ambitious} ; \eng{aspire} → \eng{aspiration/aspiring} ; \eng{skill} → \eng{skilled/skilful}.
\item Registres : \eng{I'm crazy about} (familier), \eng{I want to be} (courant), \eng{I aspire to become} (soutenu).
\end{itemize}
\end{aretenir}

\coursec{Sharing an interest : échanger sur les métiers}

Dans la section précédente, nous avons appris à exprimer notre propre intérêt pour un métier (\eng{I'm interested in becoming a tailor}, \eng{I've always dreamed of being a pilot}). Mais un intérêt ne se garde pas pour soi : il se partage. Dans la cour du lycée, à la maison, lors d'un \eng{career day}, on annonce son choix, on réagit au choix des autres, on pose des questions, on encourage. C'est exactement ce que l'examinateur attend de vous à l'oral : tenir une petite conversation naturelle sur les métiers. Cette section vous donne toutes les formules, les réactions, les questions et les expressions idiomatiques pour y parvenir.

\courssub{Observer : un dialogue après le « career day »}

\begin{textebox}[After the career day — a model dialogue]
Two students, Ngono and Ekane, are chatting after the career day organised at their high school in Yaoundé.

\textbf{Ngono:} Hi Ekane! Guess what! I've decided to take up plumbing.

\textbf{Ekane:} Really? Why plumbing?

\textbf{Ngono:} Well, I'm good with my hands, and plumbers are in great demand in Yaoundé. I'd like to tell you about my dream: I want to open my own workshop one day.

\textbf{Ekane:} That's a smart choice! How long is the training?

\textbf{Ngono:} About two years in a technical school. Then I'll learn the ropes with an experienced plumber.

\textbf{Ekane:} Is it well paid?

\textbf{Ngono:} Not at the beginning, but a skilled plumber earns good money. And you, are you still thinking of becoming a nurse?

\textbf{Ekane:} Yes! I want to follow in my mother's footsteps. She has been a nurse for twenty years.

\textbf{Ngono:} That's great! You are patient and kind — you are cut out for that job. You should go for it!

\textbf{Ekane:} Thanks! Why don't we talk to the professionals who came today? They can give us useful advice.

\textbf{Ngono:} Good idea. Let's go!

\textbf{Glossaire :} \emph{to take up} (v.) = se mettre à (un métier, une activité) ; \emph{good with my hands} (expr.) = habile de mes mains ; \emph{in great demand} (expr.) = très demandé ; \emph{workshop} (n.) = atelier ; \emph{training} (n.) = formation ; \emph{skilled} (adj.) = qualifié, compétent ; \emph{to earn money} (expr.) = gagner de l'argent ; \emph{kind} (adj.) = gentil, bienveillant ; \emph{advice} (n. invariable) = conseil(s).

\textbf{Questions de compréhension :} 1. What job has Ngono chosen, and why? 2. How long is the training to become a plumber? 3. Whose footsteps does Ekane want to follow in? 4. What advice does Ngono give Ekane?
\end{textebox}

\begin{corrige}[Questions de compréhension du dialogue]
1. \eng{Ngono has chosen plumbing because he is good with his hands and plumbers are in great demand in Yaoundé.} 2. \eng{The training lasts about two years in a technical school.} 3. \eng{She wants to follow in her mother's footsteps (her mother has been a nurse for twenty years).} 4. \eng{He tells her she is cut out for that job and that she should go for it.}
\end{corrige}

\courssub{Annoncer son choix aux autres}

Pour partager son intérêt, l'anglais dispose de formules d'annonce qu'il faut mémoriser telles quelles. Observez :

\eng{I'm thinking of becoming an electrician.} « Je pense devenir électricien. »

\eng{I'd like to tell you about my dream job.} « J'aimerais vous parler du métier de mes rêves. »

\eng{Guess what! I've decided to take up carpentry.} « Devine quoi ! J'ai décidé de me mettre à la menuiserie. »

\begin{propriete}[Les formules d'annonce]
\begin{itemize}
\item \cle{I'm thinking of + V-ing} : projet encore hésitant. \eng{I'm thinking of becoming a mechanic.}
\item \cle{I've decided to + base verbale} : décision ferme. \eng{I've decided to become a nurse.}
\item \cle{I'd like to tell you about...} : annonce plus formelle, utile à l'écrit et à l'oral d'examen.
\item \cle{Guess what!} : interjection familière pour créer la surprise avant une annonce.
\item \cle{to take up + nom de métier/activité} : se lancer dans. \eng{She took up tailoring last year.}
\end{itemize}
Attention à la structure : après \eng{think of}, le verbe est en \textbf{-ing} ; après \eng{decide}, on utilise \textbf{to + infinitif}.
\end{propriete}

\begin{attention}[Ne dites pas « I'm thinking to become »]
Faute très fréquente chez les francophones, par calque de « je pense devenir ». On dit \eng{I'm thinking \textbf{of becoming}} ou \eng{I'm thinking \textbf{about becoming}}. De même, \eng{to take up a job} signifie « se lancer dans un métier », et non « prendre un emploi » au sens de « être embauché » (qui se dit \eng{to get a job} ou \eng{to be hired}).
\end{attention}

\courssub{Réagir à l'intérêt de quelqu'un}

Quand un camarade annonce son choix, il faut réagir — un silence ou un simple « yes » est impoli en anglais ! Voici les réactions types :

\begin{center}
\begin{tabularx}{\linewidth}{|X|X|X|}
\hline
\rowcolor{popblueL} English & Français & Emploi \\
\hline
That's great! & C'est formidable ! & joie simple \\
\hline
How interesting! & Comme c'est intéressant ! & curiosité \\
\hline
Really? Why that job? & Vraiment ? Pourquoi ce métier ? & surprise + question \\
\hline
Good for you! & Tant mieux pour toi ! / Bravo ! & approbation \\
\hline
That suits you perfectly! & Cela te va parfaitement ! & compliment personnalisé \\
\hline
What a smart choice! & Quel choix judicieux ! & admiration \\
\hline
\end{tabularx}
\end{center}

Notez la construction exclamative : \cle{How + adjectif !} (\eng{How interesting!}) et \cle{What + (a) + nom !} (\eng{What a smart choice!}). Ne dites jamais \eng{What interesting!} : avec un adjectif seul, c'est \eng{how} qu'il faut.

\courssub{Poser des questions sur un métier}

Pour faire parler quelqu'un de son métier ou du métier qu'il vise, on utilise des questions fermées (\eng{Is it well paid?}) et surtout des questions en \eng{wh-}. Observez les structures :

\eng{What does a plumber do exactly?} « Que fait un plombier, exactement ? »

\eng{What qualifications do you need?} « Quels diplômes faut-il ? »

\eng{How long is the training?} « Combien de temps dure la formation ? »

\eng{Is it well paid?} « Est-ce bien payé ? »

\eng{What are the working conditions like?} « Comment sont les conditions de travail ? »

\begin{propriete}[La structure « What is ... like? »]
\cle{What + be + sujet + like ?} sert à demander une description : \eng{What is the job like?} « Comment est ce métier ? », \eng{What are the working hours like?} « Comment sont les horaires ? ». Piège majeur : le mot \eng{like} est ici une préposition (« comment »), pas le verbe « aimer ». \eng{What is he like?} = « Comment est-il (physiquement, moralement) ? », alors que \eng{What does he like?} = « Qu'est-ce qu'il aime ? ».
\end{propriete}

\courssub{Donner des informations sur un métier}

Pour décrire un métier, trois types d'informations sont attendus : les tâches, les qualités requises, la formation.

\begin{itemize}
\item \textbf{Les tâches} : présent simple + verbe d'action. \eng{A nurse looks after patients.} « Une infirmière s'occupe des patients. » \eng{A mechanic repairs cars.} \eng{A journalist writes articles and interviews people.}
\item \textbf{Les qualités requises} : \eng{You need to be patient and skilled.} « Il faut être patient et compétent. » \eng{You have to be hardworking and honest.}
\item \textbf{La formation} : \eng{It takes three years of training.} « Il faut trois ans de formation. » \eng{You need a certificate from a technical school.}
\end{itemize}

\begin{propriete}[La structure « It takes + durée »]
\cle{It takes + durée (+ to + infinitif)} exprime le temps nécessaire : \eng{It takes two years to become a welder.} « Il faut deux ans pour devenir soudeur. » Au passé : \eng{It took me six months to learn the ropes.} Ne traduisez pas par « il prend » : le sujet est le \eng{it} impersonnel, comme dans \eng{it rains}.
\end{propriete}

\courssub{Encourager et conseiller}

\begin{center}
\begin{tabularx}{\linewidth}{|X|X|X|}
\hline
\rowcolor{popblueL} English & Français & Structure \\
\hline
You should go for it! & Tu devrais te lancer ! & should + infinitif \\
\hline
Why don't you talk to a professional? & Pourquoi ne parles-tu pas à un professionnel ? & Why don't you + infinitif \\
\hline
You'd better start training early. & Tu ferais mieux de commencer tôt ta formation. & had better + infinitif \\
\hline
If I were you, I would apply now. & À ta place, je postulerais maintenant. & If I were you, I would... \\
\hline
\end{tabularx}
\end{center}

\begin{attention}[« You'd better » sans « to »]
Après \cle{had better}, le verbe est à la base verbale, \textbf{sans to} : \eng{You'd better start now} (et non \eng{to start}). De même, dans \eng{If I were you}, on garde \eng{were} même à la première personne : c'est le conditionnel irréel, jamais \eng{if I was you} dans un anglais soigné.
\end{attention}

\courssub{Expressions idiomatiques du monde du travail}

\begin{definition}[To follow in somebody's footsteps]
\cle{To follow in somebody's footsteps} = to do the same job as a parent or someone you admire — suivre la même voie que quelqu'un, marcher sur les traces de quelqu'un. Exemple : \eng{Amina followed in her mother's footsteps and became a nurse.} « Amina a suivi les traces de sa mère et est devenue infirmière. »
\end{definition}

\begin{exemplebox}[Idiomes en contexte]
\eng{He is cut out for teaching: he is patient and loves children.} « Il est fait pour l'enseignement : il est patient et aime les enfants. »

\eng{It took me six months to learn the ropes.} « Il m'a fallu six mois pour apprendre les ficelles du métier. »

\eng{My uncle is cut out for business; he sells everything!} « Mon oncle est fait pour le commerce ; il vend tout ! »

\eng{I don't want a dead-end job; I want a career with a future.} « Je ne veux pas d'un métier sans avenir ; je veux une carrière avec des perspectives. »
\end{exemplebox}

\begin{center}
\begin{tabularx}{\linewidth}{|X|X|}
\hline
\rowcolor{popblueL} Idiom & Meaning / Français \\
\hline
to follow in somebody's footsteps & suivre les traces de quelqu'un (même métier qu'un parent) \\
\hline
to learn the ropes & apprendre les ficelles du métier \\
\hline
to be cut out for (a job) & être fait pour (un métier) \\
\hline
a dead-end job & un métier sans avenir, sans promotion possible \\
\hline
to be in great demand & être très demandé (sur le marché du travail) \\
\hline
\end{tabularx}
\end{center}

\begin{attention}[« In great demand » ou « demanding » ?]
\eng{To be in great demand} signifie « être très demandé » : \eng{Plumbers are in great demand in Douala.} Ne confondez pas avec l'adjectif \eng{demanding}, qui signifie « exigeant, difficile » : \eng{Nursing is a demanding job.} « Le métier d'infirmière est exigeant. » Un métier peut être \eng{demanding} (difficile) et pourtant \eng{in great demand} (très recherché) !
\end{attention}

\courssub{Schéma d'un dialogue type}

\begin{popfigure}
\begin{center}
\begin{tikzpicture}[node distance=0.55cm,
box/.style={draw=popblue, fill=popblueL, rounded corners, align=center, text width=4.6cm, minimum height=0.9cm, font=\small},
arr/.style={-{Stealth[length=2.5mm]}, thick, popdark}]
\node[box, align=center] (a){1. Greeting\\ \eng{Hi! Guess what!}};
\node[box, below=of a, align=center] (b){2. Announcing one's choice\\ \eng{I've decided to take up plumbing.}};
\node[box, below=of b, align=center] (c){3. Reacting\\ \eng{Really? That's great! Why that job?}};
\node[box, below=of c, align=center] (d){4. Asking questions\\ training? salary? qualities?};
\node[box, below=of d, align=center] (e){5. Encouraging\\ \eng{You should go for it!}};
\node[box, below=of e, align=center] (f){6. Concluding\\ \eng{Thanks! Let's talk to a professional.}};
\draw[arr] (a) -- (b);
\draw[arr] (b) -- (c);
\draw[arr] (c) -- (d);
\draw[arr] (d) -- (e);
\draw[arr] (e) -- (f);
\end{tikzpicture}
\end{center}
\legende{Organigramme d'un dialogue type sur les métiers : six étapes à suivre à l'oral.}
\end{popfigure}

\begin{methode}[Réussir un dialogue oral sur les métiers]
\begin{enumerate}
\item Saluez et annoncez votre choix avec une formule d'annonce (\eng{I've decided to...}).
\item Réagissez toujours aux paroles de votre partenaire (\eng{That's great!}, \eng{Really?}).
\item Posez au moins deux questions : formation, salaire, conditions de travail.
\item Donnez des informations : tâches, qualités requises, durée de la formation.
\item Terminez par un encouragement ou un conseil (\eng{If I were you, I would...}).
\end{enumerate}
\end{methode}

\begin{exoresolu}[Compléter un dialogue]
Complétez ce dialogue avec les expressions étudiées : \textbf{A:} Guess what! I've decided to take up journalism. \textbf{B:} Really? (1) ............ that job? \textbf{A:} Because I love writing and meeting people. \textbf{B:} (2) ............ is the training? \textbf{A:} It takes three years. \textbf{B:} You are curious and hardworking — you are (3) ............ for it! You should (4) ............ !
\tcblower
\textbf{Solution :} (1) \eng{Why} — on demande la raison du choix (\eng{Why that job?}). (2) \eng{How long} — la réponse « three years » indique une durée. (3) \eng{cut out} — expression idiomatique : \eng{to be cut out for} = être fait pour. (4) \eng{go for it} — encouragement classique : « lance-toi ! ». Dialogue complet : \eng{B: Really? Why that job? ... B: How long is the training? ... B: You are cut out for it! You should go for it!}
\end{exoresolu}

\begin{exercice}[À vous de jouer]
1. Traduisez : « Ma sœur veut suivre les traces de notre père et devenir menuisière. » 2. Trouvez la question : \eng{............ ? — About two years in a technical school.} 3. Réagissez en anglais à cette annonce : \eng{I've decided to become a pilot.} (deux réactions différentes). 4. En binôme, construisez un dialogue de huit répliques après un \eng{career day}, en suivant l'organigramme ci-dessus.
\end{exercice}

\begin{corrige}[Exercice « À vous de jouer »]
1. \eng{My sister wants to follow in our father's footsteps and become a carpenter.} 2. \eng{How long is the training?} 3. Par exemple : \eng{That's great!} et \eng{Really? Why that job?} 4. Vérifier la présence des six étapes : salutation, annonce, réaction, questions, encouragement, conclusion.
\end{corrige}

\begin{savaistu}[Career days and career fairs]
Dans les pays anglophones (Grande-Bretagne, États-Unis, Nigeria, Ghana), les lycées organisent régulièrement des \eng{career fairs} (« foires aux métiers ») ou des \eng{career days} : des professionnels — médecins, ingénieurs, artisans, pilotes — viennent présenter leur métier aux élèves et répondre à leurs questions. Au Cameroun, de plus en plus de lycées, notamment à Yaoundé et à Douala, organisent des journées similaires. C'est le contexte idéal d'un dialogue d'examen : apprenez le dialogue modèle par cœur, vous pourrez l'adapter à n'importe quel métier !
\end{savaistu}

\coursec{Collocations et expressions du monde du travail}

Dans la section précédente, nous avons appris à \emph{partager} notre intérêt pour un métier : poser des questions, raconter un parcours, échanger lors d'un forum des métiers. Mais pour que ces échanges sonnent \emph{naturels} et non pas « traduits du français », il ne suffit pas de connaître des mots isolés comme \eng{job}, \eng{work} ou \eng{career}. Il faut savoir avec quels verbes et quelles prépositions ces mots s'associent. C'est précisément l'objet de cette section : les \cle{collocations}, c'est-à-dire les associations habituelles de mots en anglais.

\courssub{Qu'est-ce qu'une collocation ? Observation d'abord}

Lisons d'abord le court portrait suivant, puis observons les groupes de mots en gras.

\begin{textebox}[Clarisse's path to a successful career]
After \textbf{graduating from} a technical college, Clarisse \textbf{applied for a job} as an electrician. She was invited to an interview and, two weeks later, she \textbf{got the job}. She now \textbf{works for} a big construction company in Douala and \textbf{earns a good living}. Every morning, she \textbf{goes to work} at seven o'clock and \textbf{works hard} all day. She \textbf{is responsible for} a small team and \textbf{deals with} customers on building sites. She \textbf{acquired many skills} during her apprenticeship, and she now wants to \textbf{build a successful career} in her trade. Her colleagues say she will \textbf{succeed in} everything she undertakes.
\end{textebox}

\textbf{Glossaire :} \emph{to graduate from} (être diplômé de) ; \emph{to apply for} (postuler à) ; \emph{a construction company} (une entreprise de BTP) ; \emph{to earn a living} (gagner sa vie) ; \emph{an apprenticeship} (un apprentissage) ; \emph{a trade} (un métier manuel ou technique) ; \emph{to undertake} (entreprendre).

\textbf{Questions de compréhension :}
\begin{enumerate}
\item What job did Clarisse apply for?
\item Where does she work now?
\item What are her responsibilities at work?
\item What is her ambition for the future?
\end{enumerate}

Observons : on ne peut pas remplacer librement les verbes de ce texte. On dit \eng{to apply for a job} et jamais \eng{to demand a job} ; \eng{to earn a living} et jamais \eng{to win a living}. Ces associations sont figées par l'usage : ce sont des collocations.

\begin{propriete}[La collocation : une association habituelle de mots]
Une \cle{collocation} est une combinaison de mots qui « vont ensemble » naturellement en anglais : verbe + nom (\eng{to earn money}), verbe + préposition + nom (\eng{to apply for a job}), adjectif + nom (\eng{a well-paid job}), verbe + préposition (\eng{to depend on}).

\textbf{Règle d'or :} une collocation ne se traduit pas mot à mot. Elle s'apprend comme un bloc, exactement comme un mot nouveau.
\begin{itemize}
\item \eng{to earn a living} = gagner sa vie (et non \emph{to win a living})
\item \eng{to apply for a job} = postuler à un emploi (et non \emph{to apply a job})
\item \eng{to work hard} = travailler dur (et non \emph{to work hardly}, qui signifie « à peine travailler » !)
\end{itemize}
\end{propriete}

\courssub{Collocations avec \eng{job} : le nom des mille verbes}

Le mot \cle{job} (un emploi, un poste — nom \emph{comptable}) est au centre d'une constellation de verbes. Le français utilise souvent « emploi » ou « travail » avec des verbes différents ; l'anglais, lui, a ses propres couples verbe + nom.

\begin{exemplebox}[Job : les collocations essentielles]
\begin{itemize}
\item \eng{Many young Cameroonians \textbf{apply for jobs} in banks.} « Beaucoup de jeunes Camerounais postulent pour des emplois dans les banques. »
\item \eng{She has been \textbf{looking for a job} for six months.} « Elle cherche un emploi depuis six mois. »
\item \eng{He finally \textbf{found a job} in a hotel in Kribi.} « Il a finalement trouvé un emploi dans un hôtel à Kribi. »
\item \eng{Congratulations! I heard you \textbf{got the job}!} « Félicitations ! J'ai entendu que tu as décroché le poste ! »
\item \eng{My uncle \textbf{lost his job} when the factory closed.} « Mon oncle a perdu son emploi quand l'usine a fermée. »
\item \eng{She \textbf{quit her job} to start her own business.} « Elle a quitté son emploi pour créer sa propre entreprise. »
\item \eng{Teaching is a \textbf{full-time job}; selling phone credit is only a \textbf{part-time job} for him.} « L'enseignement est un emploi à temps plein ; vendre du crédit téléphonique n'est pour lui qu'un emploi à temps partiel. »
\item \eng{Nursing is a \textbf{steady job} and often a \textbf{well-paid job} abroad.} « Le métier d'infirmier est un emploi stable et souvent bien payé à l'étranger. »
\end{itemize}
\end{exemplebox}

\courssub{Collocations avec \eng{work} : verbe et nom incountable}

Rappel crucial : \cle{work} comme nom est \emph{incomptable} en anglais (on ne dit jamais \emph{a work} pour « un emploi » ; on dit \eng{a job}). Mais comme verbe, il forme des collocations très productives, notamment avec les prépositions \eng{as}, \eng{for} et \eng{with}.

\begin{propriete}[Work + prépositions : as, for, with]
\begin{itemize}
\item \eng{to work \textbf{as} + métier} : indique la fonction exercée. \eng{He works \textbf{as} a welder.} « Il travaille comme soudeur. »
\item \eng{to work \textbf{for} + employeur/entreprise} : indique pour qui on travaille. \eng{She works \textbf{for} a construction company.} « Elle travaille pour une entreprise de BTP. »
\item \eng{to work \textbf{with} + personnes/outils} : indique les collaborateurs ou les instruments. \eng{I work \textbf{with} a team of five plumbers.} « Je travaille avec une équipe de cinq plombiers. »
\end{itemize}
Autres collocations : \eng{to go to work} (aller au travail), \eng{to be at work} (être au travail), \eng{to be out of work} (être sans emploi, au chômage), \eng{to work hard} (travailler dur).
\end{propriete}

\begin{exemplebox}[Work en contexte]
\eng{He \textbf{works as} a welder \textbf{for} a construction company.} « Il travaille comme soudeur pour une entreprise de BTP. »

\eng{My father leaves home at six and \textbf{is at work} by seven.} « Mon père quitte la maison à six heures et est au travail à sept heures. »

\eng{Many graduates are \textbf{out of work} in big cities.} « Beaucoup de diplômés sont sans emploi dans les grandes villes. »
\end{exemplebox}

\courssub{Collocations avec \eng{career} : la trajectoire professionnelle}

\cle{Career} désigne une carrière, c'est-à-dire un parcours professionnel construit dans la durée. Attention : ce n'est pas un simple « job ».

\begin{exemplebox}[Career : choisir, poursuivre, bâtir]
\begin{itemize}
\item \eng{It is not easy to \textbf{choose a career} at sixteen.} « Il n'est pas facile de choisir une carrière à seize ans. »
\item \eng{She wants to \textbf{pursue a career in} medicine.} « Elle veut poursuivre une carrière en médecine. »
\item \eng{He \textbf{built a successful career in} business.} « Il a bâti une carrière réussie dans les affaires. »
\item \eng{A \textbf{career in} teaching requires patience.} « Une carrière dans l'enseignement exige de la patience. »
\end{itemize}
\end{exemplebox}

\begin{attention}[Faux ami : \eng{career} n'est pas « carrière » dans tous les sens]
\eng{Career} désigne uniquement le \emph{parcours professionnel}. Pour « faire carrière », on dit \eng{to make a career} ou \eng{to build a career}. En revanche, « carrière » au sens de lieu d'extraction de pierres se dit \eng{a quarry} ! Et le mot anglais \eng{carrier} (avec deux r) signifie « transporteur » ou « porteur » : ne confondez pas l'orthographe.
\end{attention}

\courssub{Collocations liées à la formation et à la qualification}

Avant d'exercer un métier, il faut se former. Voici les collocations du parcours de formation, très utiles pour raconter un itinéraire professionnel.

\begin{exemplebox}[Training : se former, se qualifier]
\begin{itemize}
\item \eng{All apprentices \textbf{undergo training} before they work alone.} « Tous les apprentis suivent une formation avant de travailler seuls. »
\item \eng{She \textbf{attended a training course} in computer repair.} « Elle a suivi un stage de formation en réparation informatique. »
\item \eng{He \textbf{served an apprenticeship} with a tailor in Bamenda.} « Il a fait son apprentissage chez un tailleur à Bamenda. »
\item \eng{During the course, they \textbf{acquire practical skills}.} « Pendant la formation, ils acquièrent des compétences pratiques. » (variante : \eng{to gain skills})
\item \eng{She \textbf{obtained a qualification} in accounting last year.} « Elle a obtenu un diplôme en comptabilité l'année dernière. »
\item \eng{He \textbf{graduated from} a technical college in Buea.} « Il a été diplômé d'un collège technique à Buea. »
\end{itemize}
\end{exemplebox}

\courssub{Collocations liées au revenu : gagner sa vie}

\begin{exemplebox}[Earning money : le lexique du revenu]
\begin{itemize}
\item \eng{He \textbf{earns a living} as a taxi driver in Yaoundé.} « Il gagne sa vie comme chauffeur de taxi à Yaoundé. »
\item \eng{She \textbf{earns money} by selling vegetables at the market.} « Elle gagne de l'argent en vendant des légumes au marché. »
\item \eng{Workers in this factory \textbf{get a salary} at the end of the month.} « Les ouvriers de cette usine reçoivent un salaire à la fin du mois. »
\item \eng{Some workers \textbf{are paid by the hour}; civil servants \textbf{are paid by the month}.} « Certains travailleurs sont payés à l'heure ; les fonctionnaires sont payés au mois. »
\end{itemize}
\end{exemplebox}

\begin{attention}[\eng{Earn}, \eng{make}, \eng{win} : trois verbes, trois réalités]
\begin{itemize}
\item \eng{to \textbf{earn} money} : gagner de l'argent \emph{par son travail}. \eng{She earns 150,000 francs a month.}
\item \eng{to \textbf{make} money} : gagner de l'argent \emph{en général}, souvent par une activité ou une affaire. \eng{He made a lot of money with his shop.}
\item \eng{to \textbf{win} money} : gagner de l'argent \emph{au jeu, à la loterie, à un concours}. \eng{He won money in a football bet.}
\end{itemize}
On ne peut donc jamais dire \emph{to win a salary} : un salaire se \emph{gagne} (\eng{earn}), il ne se « remporte » pas !
\end{attention}

\courssub{Verbes + prépositions à mémoriser}

Ces couples verbe + préposition sont indispensables pour parler du monde du travail. La préposition est \emph{obligatoire} et ne se devine pas depuis le français.

\begin{center}
\begin{tabularx}{\linewidth}{|X|X|X|}
\hline
\rowcolor{popblueL} Verbe + préposition & Exemple anglais & Traduction française \\
\hline
to depend on & Success \textbf{depends on} hard work. & Le succès dépend du travail acharné. \\
\hline
to specialise in & She \textbf{specialised in} car electronics. & Elle s'est spécialisée en électronique automobile. \\
\hline
to qualify as & He \textbf{qualified as} a nurse in 2022. & Il a obtenu son diplôme d'infirmier en 2022. \\
\hline
to succeed in & She \textbf{succeeded in} her final exam. & Elle a réussi à son examen final. \\
\hline
to be responsible for & He \textbf{is responsible for} ten workers. & Il est responsable de dix ouvriers. \\
\hline
to deal with & A receptionist \textbf{deals with} customers. & Un réceptionniste s'occupe des clients. \\
\hline
\end{tabularx}
\end{center}

\begin{attention}[La préposition change tout]
Ne dites jamais \emph{to succeed an exam} (calque de « réussir un examen ») : il faut \eng{to succeed \textbf{in} an exam} ou, plus simplement, \eng{to pass an exam}. De même, « dépendre de » se traduit \eng{to depend \textbf{on}}, jamais \emph{to depend of}. Et \eng{to deal with} (s'occuper de) n'a rien à voir avec « deal » au sens de « marché ».
\end{attention}

\courssub{Mots de la même famille : dérivation et morphologie}

Pour enrichir votre vocabulaire et éviter les répétitions, il est essentiel de connaître les familles de mots (nom, verbe, adjectif, personne) liées au monde du travail. L'anglais utilise beaucoup les suffixes pour changer la nature grammaticale.

\begin{propriete}[Familles de mots clés (Word Families)]
\begin{center}
\begin{tabularx}{\linewidth}{|l|X|X|X|}
\hline
\rowcolor{popblueL} Verbe & Nom (chose/action) & Nom (personne) & Adjectif \\
\hline
to employ & employment, unemployment & employer, employee & employed, unemployed \\
\hline
to work & work, workplace, workforce & worker, co-worker & working, workable \\
\hline
to qualify & qualification & qualifier & qualified, unqualified \\
\hline
to specialise & specialisation & specialist & specialised \\
\hline
to apprentice & apprenticeship & apprentice & apprentice (adj.) \\
\hline
to train & training, trainer & trainee & trained, untrained \\
\hline
to apply & application, applicant & applicant & applicable \\
\hline
to succeed & success & successor & successful, unsuccessful \\
\hline
\end{tabularx}
\end{center}

\textbf{Attention à l'accentuation (word stress) :} le stress peut changer de syllabe selon la forme.
\begin{itemize}
\item \eng{em\emph{ploy}} (verbe) $\rightarrow$ \eng{em\emph{ploy}ee} (nom personne, accent sur la dernière syllabe) $\rightarrow$ \eng{\emph{em}ployment} (nom chose, accent sur la 1ère).
\item \eng{qual\emph{i}fy} $\rightarrow$ \eng{qualific\emph{a}tion} (accent sur \emph{a}).
\item \eng{spec\emph{i}alise} $\rightarrow$ \eng{special\emph{i}sation}.
\end{itemize}
\end{propriete}

\courssub{Tableau récapitulatif des collocations essentielles}

\begin{center}
\begin{tabularx}{\linewidth}{|l|X|X|}
\hline
\rowcolor{popblueL} Verbe & Collocation & Traduction française \\
\hline
apply & to apply for a job & postuler à un emploi \\
\hline
look & to look for a job & chercher un emploi \\
\hline
lose & to lose one's job & perdre son emploi \\
\hline
quit & to quit a job & quitter/démissionner d'un emploi \\
\hline
work & to work as a driver & travailler comme chauffeur \\
\hline
work & to work for a company & travailler pour une entreprise \\
\hline
be & to be out of work & être sans emploi \\
\hline
build & to build a career & bâtir une carrière \\
\hline
pursue & to pursue a career in medicine & poursuivre une carrière en médecine \\
\hline
serve & to serve an apprenticeship & faire un apprentissage \\
\hline
gain & to gain skills & acquérir des compétences \\
\hline
earn & to earn a living & gagner sa vie \\
\hline
\end{tabularx}
\end{center}

\begin{popfigure}
\begin{tikzpicture}[mindmap, grow cyclic, every node/.style={concept, text=white, align=center, font=\small}, concept color=popblue, level 1/.append style={level distance=3.2cm, sibling angle=90}, level 2/.append style={level distance=2.4cm, sibling angle=45, concept color=poporange}]
\node{World of work}
  child[concept color=popgreen] { node {job} child { node {apply for} } child { node {lose / quit} } }
  child[concept color=poppurple] { node {work} child { node {work as / for} } child { node {out of work} } }
  child[concept color=popgold] { node {career} child { node {build / pursue} } child { node {career in} } }
  child[concept color=poppink] { node {training} child { node {serve an apprenticeship} } child { node {gain skills} } };
\end{tikzpicture}
\legende{Carte mentale des collocations du monde du travail : quatre familles autour de \eng{job}, \eng{work}, \eng{career} et \eng{training}.}
\end{popfigure}

\begin{savaistu}[Prononciation : pièges sonores du vocabulaire professionnel]
Quelques mots fréquents dont la prononciation piège les francophones (notation approximative entre guillemets français) :
\begin{itemize}
\item \eng{career} : « ca-ri-re » (accent sur le dernier \emph{i}, pas « ca-rière » à la française).
\item \eng{qualify} / \eng{qualification} : « quo-li-fa-ï » / « quo-li-fi-qué-ï-chone » (le \eng{ti} se prononce \eng{/ʃ/} « ch » dans \eng{-tion}).
\item \eng{apprenticeship} : « a-pren-ti-schip » (le \eng{ship} se prononce « chip », pas « chip » dur).
\item \eng{colleague} : « co-ligue » (le \eng{ea} se prononce \eng{/iː/}, le \eng{ue} final est muet).
\item \eng{salary} : « sa-le-ri » (trois syllabes, accent sur la première, pas « sa-la-ri »).
\item \eng{entrepreneur} : « on-tru-pre-neu » (proche du français mais accent sur \eng{neur} final).
\end{itemize}
\end{savaistu}

\courssub{Méthode : comment mémoriser durablement les collocations}

\begin{methode}[Apprendre les collocations en quatre étapes]
\begin{enumerate}
\item \textbf{N'apprends jamais un mot isolé.} Apprends le mot avec son verbe et sa préposition habituels : pas \eng{job} seul, mais \eng{to apply for a job} ; pas \eng{depend}, mais \eng{to depend on}.
\item \textbf{Crée une phrase personnelle} avec chaque collocation, en lien avec ta vie : \eng{My sister works as a hairdresser in Mvog-Ada.} Une phrase qui te concerne se retient dix fois mieux.
\item \textbf{Classe les collocations par familles} dans ton cahier (job / work / career / training / money), comme dans la carte mentale ci-dessus, et ajoute une ligne par semaine.
\item \textbf{Teste-toi à trous} : recopie tes phrases en cachant le verbe ou la préposition, puis vérifie. \eng{She works \ldots{} a nurse \ldots{} a clinic.} (\eng{as / for})
\end{enumerate}
\end{methode}

\courssub{Exercice résolu : texte à trous}

\begin{exoresolu}[Collocations en contexte]
Énoncé — Complétez le texte avec les collocations suivantes : \eng{earns a living}, \eng{applied for}, \eng{works as}, \eng{out of work}, \eng{gain skills}, \eng{responsible for}, \eng{graduated from}, \eng{successful career}.

« After he (1) \ldots{} secondary school, Divine (2) \ldots{} a job in a garage in Douala. He was (3) \ldots{} for almost a year before that, so he was very happy. He now (4) \ldots{} a mechanic and (5) \ldots{} by repairing cars and bendskins. He is (6) \ldots{} the workshop when the boss is away. Every day, he continues to (7) \ldots{} in electronics. His dream is to build a (8) \ldots{} in the automobile industry. »

\tcblower
Solution détaillée :
\begin{enumerate}
\item \eng{graduated from} — après « after », le verbe au passé convient : « après avoir été diplômé du secondaire ».
\item \eng{applied for} — collocation figée : postuler à un emploi (\eng{apply for a job}).
\item \eng{out of work} — « il était sans emploi depuis presque un an » ; la suite (« so he was very happy ») confirme le contraste.
\item \eng{works as} — \eng{work as + métier} : il travaille \emph{comme} mécanicien.
\item \eng{earns a living} — collocation figée : gagner sa vie.
\item \eng{responsible for} — \eng{to be responsible for + nom} : il est responsable de l'atelier.
\item \eng{gain skills} — acquérir des compétences ; l'infinitif suit \eng{continues to}.
\item \eng{successful career} — adjectif + nom : une carrière réussie.
\end{enumerate}
\end{exoresolu}

\begin{exercice}[À ton tour]
Complète avec la bonne préposition ou le bon verbe :
\begin{enumerate}
\item She works \ldots{} a secretary \ldots{} a bank in Bafoussam.
\item Many young people are looking \ldots{} jobs in the city.
\item He qualified \ldots{} an electrician last year.
\item A good nurse must deal \ldots{} patients kindly.
\item My brother earns his \ldots{} as a carpenter.
\end{enumerate}
\end{exercice}

\begin{corrige}[Exercice « À ton tour »]
\begin{enumerate}
\item \eng{as ; for} — elle travaille \emph{comme} secrétaire \emph{pour} une banque.
\item \eng{for} — \eng{to look for a job} : chercher un emploi.
\item \eng{as} — \eng{to qualify as} : obtenir sa qualification de.
\item \eng{with} — \eng{to deal with} : s'occuper de.
\item \eng{living} — \eng{to earn one's living} : gagner sa vie.
\end{enumerate}
\end{corrige}

\begin{aretenir}[L'essentiel de la section]
\begin{itemize}
\item Une collocation est un bloc figé : on l'apprend entière, jamais mot à mot.
\item \eng{Job} est comptable (\eng{a job}), \eng{work} est incomptable (jamais \emph{a work}).
\item Prépositions-clés : \eng{work as / for / with}, \eng{apply for}, \eng{depend on}, \eng{succeed in}, \eng{responsible for}, \eng{deal with}, \eng{qualify as}, \eng{specialise in}.
\item \eng{Earn} (par le travail) $\neq$ \eng{make} (en général) $\neq$ \eng{win} (au jeu).
\item Familles de mots : \eng{employ/employment/employee}, \eng{qualify/qualification/qualified}, \eng{specialise/specialisation/specialist}.
\item Stratégie : phrase personnelle + classement par familles + auto-test à trous.
\end{itemize}
\end{aretenir}

\coursec{Faux amis et pièges des francophones}

Dans la section précédente, nous avons appris à combiner correctement les mots du monde du travail (\emph{collocations}) : \eng{to apply for a job}, \eng{to earn a living}, \eng{to pursue a career}… Mais même avec de bonnes collocations, un élève francophone peut trahir sa pensée en traduisant mot à mot depuis le français. C'est le danger des \cle{faux amis} (\eng{false friends} ou \eng{false cognates}) : des mots qui se ressemblent dans les deux langues mais qui ne veulent pas dire la même chose. Cette section passe au crible les pièges les plus fréquents du vocabulaire des métiers et des professions.

\courssub{Qu'est-ce qu'un faux ami ? Observation d'abord}

Lisons d'abord ce court dialogue entre deux élèves d'un lycée de Yaoundé. L'un des deux commet plusieurs erreurs typiques de francophone. Sauras-tu les repérer ?

\begin{textebox}[At the career guidance office — a dialogue with mistakes]
\textbf{Marc:} Hello Sir, I am doing a formation in computer science.\\
\textbf{Mr Ebai:} You mean you are attending a \emph{training course}?\\
\textbf{Marc:} Yes Sir. And next year, I would like to do a stage in a company in Douala.\\
\textbf{Mr Ebai:} An \emph{internship}, you mean! A stage is part of a theatre!\\
\textbf{Marc:} Oh! And my father says I should assist the career fair on Saturday.\\
\textbf{Mr Ebai:} You should \emph{attend} it. To assist means to help somebody.\\
\textbf{Marc:} Thank you Sir. English is full of traps!\\
\textbf{Mr Ebai:} Indeed. But once you know the false friends, you will speak like a professional.
\end{textebox}

\begin{definition}[Faux ami (false friend)]
Un \cle{faux ami} est un mot d'une langue étrangère qui ressemble par sa forme à un mot français, mais dont le sens est différent. Exemple : \eng{a stage} ressemble à « un stage », mais signifie « une scène (de théâtre) » ou « une étape ». Les faux amis provoquent des erreurs de compréhension et d'expression ; il faut les apprendre par cœur comme du vocabulaire ordinaire.
\end{definition}

\courssub{« Une formation » n'est pas « a formation »}

C'est probablement l'erreur la plus répandue chez les élèves camerounais. En français, « je fais une formation » est une phrase banale. En anglais, \eng{a formation} désigne une \emph{formation géologique ou rocheuse} (des rochers, des montagnes) ou une \emph{formation militaire} (des soldats en rang). Pour parler d'études ou d'apprentissage, on utilise \cle{training} (nom indénombrable, sans article) ou \cle{a training course}.

\begin{attention}[« I am doing a formation » : erreur typique]
Ne dites jamais \eng{I am doing a formation} : un anglophone imaginerait que vous êtes en train de fabriquer un rocher ! Les phrases correctes sont :
\begin{itemize}
\item \eng{I am undergoing training.} « Je suis en formation. »
\item \eng{I am attending a training course.} « Je suis une formation. »
\item \eng{I am training to become a nurse.} « Je me forme pour devenir infirmier. »
\end{itemize}
Remarquez que \eng{training} est un nom \emph{indénombrable} : on ne dit jamais \eng{a training} ni \eng{trainings}. En revanche, on peut dire \eng{a training course} ou \eng{training courses}, car \eng{course} est dénombrable.
\end{attention}

\begin{exemplebox}[Exemples traduits]
\eng{She is following a six-month training course in catering.} « Elle suit une formation de six mois en restauration. »

\eng{The company offers free training to its workers.} « L'entreprise offre une formation gratuite à ses employés. »

\eng{We saw impressive rock formations in the North Region.} « Nous avons vu d'impressionnantes formations rocheuses dans la région du Nord. » (ici, \eng{formations} est correct car il s'agit de géologie !)
\end{exemplebox}

\courssub{« Un stage » n'est pas « a stage »}

Même mécanisme : le mot anglais \eng{stage} existe, mais il signifie « une scène » (théâtre) ou « une étape » (d'un projet, d'une course cycliste). Pour traduire « un stage » professionnel, on dit \cle{an internship} (terme surtout américain) ou \cle{a work placement} (terme surtout britannique). Le stagiaire est \cle{an intern}.

\begin{exemplebox}[Exemples traduits]
\eng{She did an internship in a hospital in Buea.} « Elle a fait un stage dans un hôpital à Buea. »

\eng{Our school organises work placements in local workshops.} « Notre école organise des stages dans des ateliers locaux. »

\eng{The actors appeared on stage.} « Les acteurs sont apparus sur scène. » (sens théâtral)

\eng{The first stage of the project is finished.} « La première étape du projet est terminée. » (sens « étape »)
\end{exemplebox}

\courssub{« Démissionner » n'est pas « to demission »}

Le verbe français « démissionner » n'a pas d'équivalent anglais en \emph{-mission}. Le verbe \eng{to demission} n'existe pas en anglais courant. Pour dire « démissionner », on utilise \cle{to resign} (souvent suivi de \eng{from}) ou, de façon plus familière, \cle{to quit one's job}.

\begin{exemplebox}[Exemples traduits]
\eng{He resigned from his job last year.} « Il a démissionné de son emploi l'année dernière. »

\eng{The headmistress resigned for personal reasons.} « La proviseure a démissionné pour des raisons personnelles. »

\eng{She quit her job to start her own business.} « Elle a démissionné pour créer sa propre entreprise. »

Nom correspondant : \eng{a resignation} « une démission » — \eng{He handed in his resignation.} « Il a remis sa démission. »

Attention à la prononciation de \eng{resign} : le \emph{s} se prononce « z » et le \emph{g} est muet : « ri-ZAÏNE ». Ne pas confondre avec \eng{to re-sign} (« signer à nouveau »), où le \emph{g} se prononce.
\end{exemplebox}

\courssub{« To assist » et « to attend » : le grand classique}

Voici un piège qui touche même les bons élèves. Le verbe français « assister à » signifie « être présent à ». Mais le verbe anglais \eng{to assist} signifie « aider » ! Pour « assister à » un cours, une réunion, une cérémonie, il faut utiliser \cle{to attend} (sans préposition).

\begin{attention}[To attend ≠ to assist]
\begin{itemize}
\item \eng{to attend} = assister à, être présent à (un cours, une réunion, un salon). \eng{I attended a career fair.} « J'ai assisté à un salon des métiers. »
\item \eng{to assist} = aider, prêter assistance. \eng{The nurse assists the doctor.} « L'infirmière aide le médecin. »
\end{itemize}
Attention aussi aux noms : \eng{an assistant} est « un assistant, un auxiliaire » (une personne qui aide), et \eng{assistance} signifie « aide », jamais « l'assistance » au sens de public présent. Pour « l'assistance » (le public), on dit \eng{the audience}.
\end{attention}

\courssub{Autres pièges fréquents du monde du travail}

\begin{propriete}[Pièges lexicaux à connaître par cœur]
\begin{itemize}
\item « Gagner sa vie » = \cle{to earn one's living} (et non \eng{to win one's life} : \eng{to win} signifie « gagner » un match, un prix, une compétition). \eng{He earns his living as a carpenter.} « Il gagne sa vie comme menuisier. »
\item « Être au chômage » = \cle{to be unemployed} ou \cle{to be out of work} (et non \eng{to be in unemployment}). \eng{Many young graduates are unemployed.} « Beaucoup de jeunes diplômés sont au chômage. »
\item « Un patron » = \cle{a boss} ou \cle{an employer} ; « un employé » = \cle{an employee}. Ne pas inverser : \eng{employer} = celui qui emploie, \eng{employee} = celui qui est employé. Attention : \eng{a patron} en anglais désigne un mécène ou un client habituel, pas un chef d'entreprise.
\item « La carrière » = \cle{a career}, correct ; mais \eng{a career} ne désigne jamais une carrière de pierres : celle-ci se dit \cle{a quarry}. Prononciation de \eng{career} : « ka-RIEUR » (accent sur la deuxième syllabe).
\item « Un diplôme » = \eng{a diploma}, mais attention à la prononciation : on dit « di-PLO-ma » (accent sur la deuxième syllabe), et non « di-plôm ».
\item « Une licence » (universitaire) = \eng{a bachelor's degree} ; \eng{a licence} désigne un permis officiel (comme \eng{a driving licence}, « un permis de conduire »).
\item « Une candidature » = \cle{an application} ; « postuler » = \cle{to apply for a job} (jamais \eng{to postulate}). \eng{She applied for a teaching post.} « Elle a postulé pour un poste d'enseignante. »
\end{itemize}
\end{propriete}

\begin{propriete}[Comment traduire « un métier » ?]
Il n'existe pas de mot anglais unique pour « métier ». Choisissez selon le contexte :
\begin{itemize}
\item \cle{trade} : métier manuel ou artisanal appris par la pratique (\eng{a carpenter's trade}, « le métier de menuisier »).
\item \cle{profession} : métier intellectuel exigeant des diplômes (\eng{the teaching profession}, « la profession enseignante »).
\item \cle{job} : un emploi précis, un poste (\eng{She found a job in Douala.}).
\item \cle{occupation} : terme général et administratif, utilisé dans les formulaires (\eng{Occupation: teacher}).
\end{itemize}
\end{propriete}

\courssub{Tableau récapitulatif des faux amis}

\begin{center}
\begin{tabularx}{\linewidth}{|X|X|X|}
\hline
\rowcolor{popblueL} \textbf{Mot français} & \textbf{Faux ami à éviter} & \textbf{Traduction correcte} \\
\hline
une formation & a formation (rochers, soldats) & training / a training course \\
\hline
un stage & a stage (scène, étape) & an internship / a work placement \\
\hline
démissionner & to demission (n'existe pas) & to resign / to quit one's job \\
\hline
assister à (une réunion) & to assist (aider) & to attend \\
\hline
gagner sa vie & to win one's life & to earn one's living \\
\hline
être au chômage & to be in unemployment & to be unemployed / out of work \\
\hline
un patron & a patron (un mécène) & a boss / an employer \\
\hline
une carrière de pierres & a career & a quarry \\
\hline
postuler & to postulate & to apply for \\
\hline
un métier (général) & a trade uniquement & job / trade / profession / occupation \\
\hline
\end{tabularx}
\end{center}

\begin{popfigure}
\begin{tikzpicture}[mindmap, grow cyclic,
  every node/.style={concept, align=center, font=\small},
  root concept/.append style={concept color=popblue, font=\small\bfseries},
  level 1 concept/.append style={level distance=3.4cm, sibling angle=72, font=\scriptsize},
  level 2 concept/.append style={level distance=2.6cm, sibling angle=50, font=\scriptsize}]
\node[root concept, align=center]{False friends\\ work vocabulary}
  child[concept color=poporange] { node {formation} child[concept color=poporangeL] { node[align=center]{training\\ course} } }
  child[concept color=popgreen] { node {stage} child[concept color=popgreenL] { node[align=center]{internship\\ placement} } }
  child[concept color=popgold] { node {assister à} child[concept color=popgoldL] { node {to attend} } }
  child[concept color=poppurple] { node {démissionner} child[concept color=poppurpleL] { node[align=center]{to resign\\ to quit} } }
  child[concept color=poppink] { node {gagner sa vie} child[concept color=poppinkL] { node[align=center]{to earn\\ one's living} } };
\end{tikzpicture}
\legende{Carte mentale : cinq faux amis à chasser de votre anglais et leurs traductions correctes.}
\end{popfigure}

\courssub{Texte de consolidation : « From a mistake to a lesson »}

\begin{textebox}[An original reading text — False friends at the job interview]
Last month, my cousin Sandrine went to a job interview in Douala. She had prepared carefully, but she was nervous. When the manager asked about her studies, she said: “I did a formation in accounting.” The manager smiled politely and asked: “A rock formation?” Sandrine blushed and corrected herself: “Sorry, I mean a training course.”

Then she explained that during her studies, she had done “a stage” at a bank in Bonaberi. The manager laughed gently: “On stage? Are you an actress?” Sandrine laughed too and replied: “No Sir, an internship! I worked as an intern for three months.”

The interview continued. Sandrine said she wanted to assist the training sessions organised by the company. The manager raised an eyebrow: “You want to help the trainers?” She quickly answered: “No, I want to attend them — to be present and learn.”

At the end, the manager shook her hand and said: “Your English is good, but be careful with false friends. Words like formation, stage and assist do not mean what French speakers think.” Sandrine got the job. Today, she earns her living as an accountant, and she keeps a small notebook where she writes every false friend she discovers. Her dream is to build a long career in banking, and maybe one day to become the boss of her own branch.

Her advice to all students is simple: “Never translate word for word. Learn whole expressions, and your English will sound natural.”
\end{textebox}

\textbf{Glossaire :} \emph{to blush} (v.) rougir ; \emph{to raise an eyebrow} (exp.) hausser un sourcil (montrer la surprise) ; \emph{an accountant} (n.) un comptable ; \emph{a branch} (n.) une succursale, une agence ; \emph{word for word} (exp.) mot à mot.

\textbf{Questions de compréhension :}
\begin{enumerate}
\item What three false friends did Sandrine use during her interview?
\item How did she correct each mistake?
\item What does she do today, and what is her dream?
\item What advice does she give to students?
\end{enumerate}

\begin{corrige}[Questions de compréhension]
1. She used \eng{a formation} (instead of \eng{a training course}), \eng{a stage} (instead of \eng{an internship}) and \eng{to assist} (instead of \eng{to attend}).
2. She corrected herself by saying \eng{a training course}, \eng{an internship} and \eng{to attend}.
3. She now works as an accountant and earns her living in banking; her dream is to build a long career and become the boss of her own branch.
4. She advises students never to translate word for word and to learn whole expressions.
\end{corrige}

\courssub{Exercice résolu et entraînement}

\begin{exoresolu}[Corrige les erreurs de francophone]
Énoncé — Les phrases suivantes contiennent chacune un faux ami. Corrige-les.
\begin{enumerate}
\item My brother is doing a formation to become an electrician.
\item I want to assist the meeting about jobs on Friday.
\item After two years, she demissioned and opened a shop.
\item He wins his life as a taxi driver in Bamenda.
\end{enumerate}
\tcblower
Solution détaillée :
\begin{enumerate}
\item \eng{My brother is attending a training course to become an electrician.} — \eng{a formation} est un faux ami ; pour des études professionnelles, on dit \eng{a training course} ou \eng{training}.
\item \eng{I want to attend the meeting about jobs on Friday.} — \eng{to assist} signifie « aider » ; « assister à » se dit \eng{to attend}, sans préposition.
\item \eng{After two years, she resigned and opened a shop.} — le verbe \eng{to demission} n'existe pas ; « démissionner » se dit \eng{to resign} (ou \eng{to quit her job}).
\item \eng{He earns his living as a taxi driver in Bamenda.} — « gagner sa vie » se dit \eng{to earn one's living} ; \eng{to win} s'emploie pour un match ou un prix.
\end{enumerate}
\end{exoresolu}

\begin{exercice}[À toi de jouer]
\begin{enumerate}
\item Traduis en anglais : « Elle a fait un stage dans une banque à Yaoundé. » / « Mon oncle est au chômage depuis janvier. » / « Le patron a embauché dix employés. »
\item Choisis le bon mot : \eng{The teacher (assists / attends) the students with their homework.} / \eng{We (assisted / attended) a career fair last Saturday.}
\item Trouve le faux ami dans cette phrase et corrige-le : \eng{My sister wants to win her life as a nurse.}
\end{enumerate}
\end{exercice}

\begin{corrige}[À toi de jouer]
1. \eng{She did an internship in a bank in Yaoundé.} / \eng{My uncle has been unemployed since January.} (present perfect avec \eng{since} : la situation continue) / \eng{The boss hired ten employees.}
2. \eng{The teacher assists the students with their homework.} (le professeur \emph{aide} les élèves) / \eng{We attended a career fair last Saturday.} (nous étions \emph{présents} au salon).
3. Le faux ami est \eng{to win her life} ; correction : \eng{My sister wants to earn her living as a nurse.}
\end{corrige}

\begin{experience}[Jeu de rôle : the job interview]
Par groupes de deux, jouez une interview d'embauche. L'élève A est le candidat et doit glisser volontairement deux faux amis dans ses réponses (\eng{formation}, \eng{stage}, \eng{assist}, \eng{win my life}…). L'élève B est le directeur : il doit repérer chaque erreur, demander une clarification avec humour (\eng{“A rock formation?”}), puis inviter le candidat à se corriger. Ensuite, échangez les rôles. La classe note les faux amis utilisés et leurs corrections au tableau.
\end{experience}

\begin{savaistu}[D'où vient le mot « vocation » ?]
Le mot anglais \eng{vocation} vient du latin \emph{vocare}, qui signifie « appeler ». Une vocation, c'est donc littéralement un « appel » : l'appel intérieur qui pousse quelqu'un vers un métier, comme la médecine, l'enseignement ou le sacerdoce. Le même mot latin a donné « vocabulaire » (l'ensemble des mots qu'on « appelle » pour parler) et « provoquer » (appeler vers soi). Quand tu dis \eng{Teaching is my vocation}, tu dis en réalité : « L'enseignement m'appelle » !
\end{savaistu}

\begin{aretenir}[L'essentiel de la section]
\begin{itemize}
\item « Une formation » = \eng{training} / \eng{a training course} ; « un stage » = \eng{an internship} / \eng{a work placement}.
\item « Démissionner » = \eng{to resign} ; « assister à » = \eng{to attend} ; « aider » = \eng{to assist}.
\item « Gagner sa vie » = \eng{to earn one's living} ; « être au chômage » = \eng{to be unemployed}.
\item « Postuler » = \eng{to apply for} ; « un patron » = \eng{a boss} / \eng{an employer}.
\item « Métier » se traduit selon le contexte : \eng{trade}, \eng{profession}, \eng{job} ou \eng{occupation}.
\item Règle d'or : ne jamais traduire mot à mot ; apprendre les expressions entières en contexte.
\end{itemize}
\end{aretenir}

\coursec{Mots de la même famille et réemploi en contexte}

Dans la section précédente, nous avons débusqué les faux amis et les calques du français qui guettent l'élève francophone (\eng{formation} n'est pas \eng{training}, \eng{profession} ne veut pas toujours dire « profession » au sens de métier manuel, etc.). Nous allons maintenant exploiter une arme très puissante pour enrichir votre vocabulaire : les \cle{familles de mots}. En anglais comme en français, un même radical donne naissance à tout un réseau de mots : verbe, nom de personne, nom abstrait, adjectif, adverbe. Connaître ces réseaux, c'est multiplier par quatre ou cinq votre vocabulaire actif sans effort supplémentaire de mémorisation — et c'est aussi la clé pour \emph{varier} vos productions écrites au lieu de répéter le même mot dix fois.

\courssub{Observer avant de mémoriser : un texte pour découvrir les familles}

\begin{textebox}[A career talk at Government High School, Buea]
Last Friday, our school invited a careers officer to talk to the Form One classes. She explained that \textbf{employers} in the South-West Region are looking for \textbf{skilled} and \textbf{qualified} young people, but that many school leavers are still \textbf{unemployed} because they lack practical \textbf{training}. “\textbf{Unemployment} is not an \textbf{inevitability},” she said. “If you \textbf{train} seriously, you become \textbf{employable}.”

She gave the example of her own brother. He \textbf{trained} as an electrician for two years. During his \textbf{training}, he was a \textbf{trainee} in a small workshop in Molyko, and his \textbf{trainer} taught him everything about wiring and safety. Today he is a respected \textbf{professional}; he loves his \textbf{occupation} and he even \textbf{employs} two \textbf{employees} of his own.

Her advice was clear: choose a trade you are cut out for, get the right \textbf{qualifications}, and never stop learning. “An \textbf{unskilled} worker complains about the job market,” she concluded, “but a well-\textbf{trained} worker creates his own opportunities.”
\end{textebox}

\textbf{Glossaire.} \emph{careers officer} (n.) : conseiller d'orientation ; \emph{Form One} (n.) : classe de 6\up{me} (premier cycle secondaire) ; \emph{school leaver} (n.) : jeune sorti de l'école ; \emph{inevitability} (n.) : fatalité, inévitabilité ; \emph{workshop} (n.) : atelier ; \emph{wiring} (n.) : câblage électrique ; \emph{to be cut out for} (expr.) : être fait pour ; \emph{job market} (n.) : marché de l'emploi.

\textbf{Questions de compréhension.}
\begin{enumerate}
\item What are employers in the South-West Region looking for?
\item Why are many school leavers still unemployed, according to the speaker?
\item How long did her brother's training last, and where did he do it?
\item What is the speaker's final piece of advice?
\end{enumerate}

\textbf{Observation.} Soulignez dans le texte tous les mots construits sur \eng{employ}, \eng{train}, \eng{skill}, \eng{qualify} et \eng{profession}. Vous constatez qu'un seul radical produit des verbes, des noms de personnes, des noms abstraits et des adjectifs. C'est exactement ce mécanisme que nous allons systématiser.

\courssub{Les grandes familles de mots du monde du travail}

\begin{definition}[Famille de mots]
Une \cle{famille de mots} (\eng{word family}) est un ensemble de mots construits sur le même radical et qui appartiennent à des classes grammaticales différentes : verbe, nom de personne, nom abstrait, adjectif, adverbe. Exemple : \eng{to employ} (verbe) $\rightarrow$ \eng{employer} (personne qui emploie) $\rightarrow$ \eng{employee} (personne employée) $\rightarrow$ \eng{employment} (nom abstrait) $\rightarrow$ \eng{unemployed} (adjectif).
\end{definition}

\begin{propriete}[La famille de \eng{employ}]
\begin{center}
\begin{tabular}{|l|l|l|}
\hline
\rowcolor{popblueL} \textbf{Mot} & \textbf{Classe} & \textbf{Traduction} \\
\hline
to employ & verbe & employer, embaucher \\
\hline
employer & nom (personne) & l'employeur, le patron \\
\hline
employee & nom (personne) & l'employé(e), le salarié \\
\hline
employment & nom abstrait & l'emploi, le travail \\
\hline
unemployment & nom abstrait & le chômage \\
\hline
unemployed & adjectif & au chômage, sans emploi \\
\hline
employable & adjectif & employable, apte à être embauché \\
\hline
self-employed & adjectif & travailleur indépendant, à son compte \\
\hline
\end{tabular}
\end{center}
Notez la logique des suffixes : \textbf{-er} marque celui qui fait l'action (\eng{employer} = celui qui emploie), \textbf{-ee} marque celui qui la subit (\eng{employee} = celui qui est employé). Même logique : \eng{trainer} / \eng{trainee}, \eng{interviewer} / \eng{interviewee}.
\end{propriete}

\begin{popfigure}
\begin{center}
\begin{tikzpicture}[
  grow=down,
  level distance=1.6cm,
  level 1/.style={sibling distance=3.4cm},
  level 2/.style={sibling distance=2.6cm},
  every node/.style={draw, rounded corners=2pt, align=center, font=\small, inner sep=3pt},
  edge from parent/.style={draw=popblue, thick, -{Stealth[length=2mm]}}
]
\node[fill=popblue, text=white, font=\small\bfseries, align=center]{to employ\\ \emph{employer (v.)}}
  child { node[fill=poporangeL, align=center]{employer\\ \emph{le patron}} }
  child { node[fill=poporangeL, align=center]{employee\\ \emph{le salarié}} }
  child { node[fill=popgreenL, align=center]{employment\\ \emph{l'emploi}}
    child { node[fill=poppinkL, align=center]{unemployment\\ \emph{le chômage}} }
    child { node[fill=poppinkL, align=center]{unemployed\\ \emph{au chômage}} }
  }
  child { node[fill=popgreenL, align=center]{employable\\ \emph{apte à l'embauche}} }
  child { node[fill=popgreenL, align=center]{self-employed\\ \emph{indépendant}} };
\end{tikzpicture}
\end{center}
\legende{Arbre lexical de la famille \eng{employ} : chaque branche porte la traduction française.}
\end{popfigure}

\begin{propriete}[Les autres familles essentielles de la leçon]
\begin{center}
\begin{tabularx}{\linewidth}{|l|X|X|}
\hline
\rowcolor{popblueL} \textbf{Radical} & \textbf{Famille} & \textbf{Traductions} \\
\hline
to train & trainer / trainee / training / trained & former / formateur / stagiaire / formation / formé \\
\hline
to occupy & occupation / occupational / occupant & métier, occupation / professionnel(le) / occupant \\
\hline
to qualify & qualification / qualified / unqualified / qualifying & diplôme, qualification / qualifié / non qualifié / qualificatif \\
\hline
skill (n.) & skilled / unskilled / skillful & compétence / qualifié / non qualifié / habile \\
\hline
profession & professional / professionally & profession / professionnel / professionnellement \\
\hline
\end{tabularx}
\end{center}
Deux collocations à retenir absolument : \eng{occupational hazards} = les risques du métier ; \eng{a skilled worker} = un ouvrier qualifié.
\end{propriete}

\begin{exemplebox}[Les familles en contexte]
\begin{itemize}
\item \eng{The employer hired three new employees.} « L'employeur a embauché trois nouveaux employés. »
\item \eng{Unemployment is a serious problem for young graduates.} « Le chômage est un problème sérieux pour les jeunes diplômés. »
\item \eng{She completed her training as a nurse in Douala.} « Elle a terminé sa formation d'infirmière à Douala. »
\item \eng{The trainer gives the trainees practical exercises every morning.} « Le formateur donne aux stagiaires des exercices pratiques chaque matin. »
\item \eng{You need good qualifications to apply for this job.} « Il faut de bons diplômes pour postuler à cet emploi. »
\item \eng{Fumes and noise are occupational hazards for welders.} « Les fumées et le bruit sont les risques du métier pour les soudeurs. »
\item \eng{He is a skilled carpenter; his brother is still unskilled.} « C'est un menuisier qualifié ; son frère est encore non qualifié. »
\item \eng{She behaves professionally in every situation.} « Elle se comporte professionnellement en toute situation. »
\end{itemize}
\end{exemplebox}

\courssub{Les suffixes des noms de métiers}

\begin{propriete}[Former un nom de métier : -er, -or, -ist, -ian, -ant]
L'anglais forme la plupart des noms de métiers en ajoutant un suffixe à un verbe ou à un nom :
\begin{itemize}
\item \textbf{-er} (le plus fréquent) : \eng{teach $\rightarrow$ teacher}, \eng{bake $\rightarrow$ baker}, \eng{farm $\rightarrow$ farmer}, \eng{work $\rightarrow$ worker} ;
\item \textbf{-or} : \eng{act $\rightarrow$ actor}, \eng{direct $\rightarrow$ director}, \eng{tailor} (tailleur), \eng{doctor} ;
\item \textbf{-ist} : \eng{art $\rightarrow$ artist}, \eng{journal $\rightarrow$ journalist}, \eng{dentist}, \eng{typist} ;
\item \textbf{-ian} : \eng{music $\rightarrow$ musician}, \eng{electric $\rightarrow$ electrician}, \eng{politics $\rightarrow$ politician} ;
\item \textbf{-ant / -ent} : \eng{account $\rightarrow$ accountant} (comptable), \eng{study $\rightarrow$ student}, \eng{assist $\rightarrow$ assistant}.
\end{itemize}
Attention : ces suffixes sont \textbf{inaccentués} ; l'accent tonique reste sur le radical : \eng{TEAcher}, \eng{ACTor}, \eng{ARTist}, mais \eng{elecTRIcian}, \eng{muSIcian} (l'accent se déplace vers la syllabe qui précède \emph{-ian}).
\end{propriete}

\begin{center}
\begin{tabularx}{\linewidth}{|l|X|X|}
\hline
\rowcolor{popblueL} \textbf{Suffixe} & \textbf{Exemples} & \textbf{Prononciation (repère français)} \\
\hline
-er & teacher, baker, driver, plumber & « eur » très bref, presque muet : « ti-tcheu(r) » \\
\hline
-or & actor, doctor, tailor, director & « eur » bref aussi : « ak-teu(r) » \\
\hline
-ist & artist, dentist, journalist, typist & « iste » : « ar-tiste » \\
\hline
-ian & electrician, musician, politician & « i-enn » accentué : « é-lek-tri-chienn » \\
\hline
-ant / -ent & accountant, student, assistant & « ent » bref : « a-kaon-tent » \\
\hline
\end{tabularx}
\end{center}

\begin{attention}[Orthographe : les pièges à éviter]
\begin{itemize}
\item \eng{employee} prend \textbf{deux « e »} à la fin (employ + ee). N'écrivez jamais \emph{employe} (calque du français « employé »).
\item \eng{skilled} : le radical \eng{skill} se termine déjà par deux \textbf{l} ; on ajoute simplement \emph{-ed} sans doubler davantage : \eng{skilled}. De même \eng{qualified} : qualify $\rightarrow$ le \emph{y} devient \emph{i} devant \emph{-ed} : \eng{qualified} (et non \emph{qualifyed}).
\item \eng{training} prend \textbf{un seul « n »} (train + ing), car « ai » est un digramme vocalique (voyelle longue) : pas de consonne à doubler. Mais \eng{beginner}, \eng{runner} doublent la consonne (voyelle courte unique accentuée).
\item \eng{occupation} : deux « c », un seul « p ». \eng{professional} : un « f », deux « s », un « l » final pour l'adjectif ; \eng{professionally} en prend deux.
\end{itemize}
\end{attention}

\begin{exoresolu}[Complétez avec le mot de la famille qui convient]
Énoncé — Complétez chaque phrase avec la forme correcte du mot entre parenthèses :
\begin{enumerate}
\item Many young \dots\ (employ) people sell phone credit in the streets of Yaoundé.
\item Her \dots\ (train) as a midwife lasted three years.
\item A \dots\ (skill) tailor can sew a complete suit in two days.
\item He works \dots\ (profession) as an accountant in a bank.
\item The \dots\ (employ) of this factory treats his workers well.
\end{enumerate}
\tcblower
Solution détaillée :
\begin{enumerate}
\item \textbf{unemployed} — il faut un adjectif qualifiant \eng{people} ; le sens (« vendent du crédit téléphonique pour survivre ») exige le préfixe négatif \emph{un-}.
\item \textbf{training} — après le possessif \eng{her}, il faut un nom ; \eng{training} (formation) est le nom abstrait de la famille de \eng{train}.
\item \textbf{skilled} — devant le nom \eng{tailor}, il faut l'adjectif : « un tailleur qualifié ». Orthographe : skill + ed = \eng{skilled}.
\item \textbf{professionally} — l'adverbe modifie le verbe \eng{works} : « il travaille professionnellement ».
\item \textbf{employer} — le sujet de \eng{treats} est une personne qui emploie : le patron. Ne pas confondre avec \eng{employee} (le salarié).
\end{enumerate}
\end{exoresolu}

\courssub{Réemploi en contexte : rédiger « My dream job »}

\begin{methode}[Enrichir une production écrite grâce aux familles de mots]
Pour obtenir une bonne note au Bac, évitez de répéter le même mot. Procédez ainsi :
\begin{enumerate}
\item Choisissez deux ou trois familles de mots de la leçon (\eng{employ}, \eng{train}, \eng{skill}, \eng{qualify}).
\item Dans chaque famille, utilisez des classes différentes : un verbe, un nom de personne, un nom abstrait, un adjectif.
\item Modèle de variation : \eng{She trained as a nurse. Her training lasted three years. She is now a well-trained professional.} « Elle s'est formée comme infirmière. Sa formation a duré trois ans. C'est maintenant une professionnelle bien formée. » — Trois occurrences de la même famille, trois classes grammaticales, zéro répétition.
\item Ajoutez au moins une collocation de la leçon : \eng{occupational hazards}, \eng{in great demand}, \eng{self-employed}, \eng{to be cut out for}.
\item Relisez en vérifiant l'orthographe des mots pièges : \eng{employee}, \eng{skilled}, \eng{qualified}, \eng{training}.
\end{enumerate}
\end{methode}

\begin{textebox}[Model paragraph — My dream job]
My dream job is to become an electrician. I developed an interest in electricity when I watched my uncle repair generators in our quarter. Electricians are in great demand in Cameroon, and the training takes only two years. I know I am cut out for this trade because I am skilled with my hands. After my training, I will apply for a job or become self-employed.
\end{textebox}

\textbf{Analyse du modèle.} Ce court paragraphe réemploie : \eng{job}, \eng{developed an interest in} (section 2), \eng{in great demand} et \eng{cut out for} (collocations de la section 4), \eng{training} (nom abstrait), \eng{skilled} (adjectif), \eng{trade} (champ lexical), \eng{apply for a job} et \eng{self-employed} (famille de \eng{employ}). Soit huit éléments de la leçon en cinq phrases : c'est exactement la densité attendue.

\begin{experience}[Production guidée : My dream job]
Rédigez un paragraphe de 8 à 10 lignes intitulé \eng{My dream job}. Consignes obligatoires :
\begin{enumerate}
\item Réemployez au moins \textbf{huit mots ou expressions de la leçon}, dont au moins trois familles de mots différentes (par exemple \eng{employ}, \eng{train}, \eng{qualify}).
\item Utilisez au moins une fois un verbe, un nom de personne et un nom abstrait de la même famille (ex. \eng{to train} / \eng{trainee} / \eng{training}).
\item Insérez au moins deux collocations : \eng{in great demand}, \eng{occupational hazards}, \eng{skilled worker}, \eng{self-employed}, \eng{apply for a job}, \eng{be cut out for}.
\item Échangez ensuite votre paragraphe avec un voisin : chacun souligne les mots de la leçon et vérifie l'orthographe de \eng{employee}, \eng{skilled}, \eng{qualified} et \eng{training}.
\end{enumerate}
\end{experience}

\begin{exercice}[Pour s'entraîner]
\begin{enumerate}
\item Donnez la famille complète (verbe, nom de personne, nom abstrait, adjectif) de : \eng{employ}, \eng{train}, \eng{qualify}.
\item Formez le nom de métier : \eng{to teach}, \eng{to act}, \eng{art}, \eng{music}, \eng{electricity}, \eng{to bake}, \eng{to account}.
\item Traduisez : « Les ouvriers qualifiés sont très demandés, mais beaucoup de jeunes sans formation restent au chômage. »
\item Complétez : \eng{A t\dots\ trains t\dots\ during their t\dots.} (famille de \eng{train}, trois formes différentes).
\end{enumerate}
\end{exercice}

\begin{corrige}[Exercice « Pour s'entraîner »]
\begin{enumerate}
\item \eng{employ} : to employ / employer, employee / employment / employed, unemployed, employable. \eng{train} : to train / trainer, trainee / training / trained. \eng{qualify} : to qualify / — / qualification / qualified, unqualified.
\item teacher, actor, artist, musician, electrician, baker, accountant.
\item \eng{Skilled workers are in great demand, but many untrained young people remain unemployed.}
\item \eng{A \textbf{trainer} trains \textbf{trainees} during their \textbf{training}.} « Un formateur forme des stagiaires pendant leur formation. »
\end{enumerate}
\end{corrige}

\begin{aretenir}[L'essentiel de la section]
\begin{itemize}
\item Une famille de mots = un radical + des classes grammaticales variées : \eng{to employ $\rightarrow$ employer / employee / employment / unemployment / unemployed / employable}.
\item Le suffixe \textbf{-er} désigne celui qui agit, \textbf{-ee} celui qui subit : \eng{employer/employee}, \eng{trainer/trainee}.
\item Les suffixes de métiers : \textbf{-er} (teacher), \textbf{-or} (actor), \textbf{-ist} (artist), \textbf{-ian} (electrician), \textbf{-ant} (accountant).
\item Orthographe piège : \eng{employee} (deux « e »), \eng{qualified} (y $\rightarrow$ i), \eng{training} (un seul « n »).
\item Pour enrichir vos écrits : variez les classes grammaticales au sein d'une même famille au lieu de répéter un seul mot.
\end{itemize}
\end{aretenir}

\newpage

\coursec{Activité d'intégration}

\courssub{Objectifs de l'activité}

À la fin de cette activité d'intégration, tu seras capable de :
\begin{itemize}
\item réemployer à l'oral le vocabulaire des métiers, des qualifications et des conditions de travail dans un entretien authentique ;
\item utiliser correctement au moins cinq \cle{collocations} (par exemple \eng{to earn a living}, \eng{to apply for a job}, \eng{to gain experience}) et deux \cle{expressions idiomatiques} (par exemple \eng{to climb the career ladder}, \eng{a dead-end job}) ;
\item poser des questions fermées et ouvertes sur un métier avec les structures interrogatives correctes ;
\item rédiger une courte affiche informative et attractive en anglais, en employant les mots de la même famille (\eng{employ}, \eng{employer}, \eng{employee}, \eng{employment}) ;
\item justifier un choix à l'oral avec \eng{I am keen on... because...}, \eng{I am interested in...}, \eng{What attracts me is...}.
\end{itemize}

\courssub{Situation}

Ton lycée organise, comme dans les écoles britanniques et américaines, un \cle{Career Day} : une journée où des professionnels de la région viennent présenter leur métier aux élèves. Cette année, le proviseur a demandé au club d'anglais de préparer l'événement entièrement en anglais : interviews des invités et affiches \eng{Job of the Week} pour le hall du lycée. Tu es membre du club et tu vas jouer deux rôles : celui d'un professionnel camerounais et celui d'un élève journaliste.

\begin{textebox}[Poster displayed in the school hall — Government High School, Buea]
CAREER DAY 2025 — FIND YOUR PATH!

Are you interested in a trade or a profession? Do you want to know what a nurse, a carpenter, an engineer or a shopkeeper really does every day?

Come to the school hall on Friday 14th March, from 10 a.m. to 2 p.m.

Local professionals will answer your questions:
-- What does your job consist of?
-- What qualifications do you need?
-- Is it well paid?
-- Why did you choose this trade?

Students of the English Club will interview the guests and design “Job of the Week” posters. At the end of the day, vote for the most attractive job and say why: “I am keen on this job because...”

Don't miss it! Your future career starts here.
\end{textebox}

\begin{definition}[Glossaire du document]
\begin{itemize}
\item \eng{a trade} (nom) : un métier manuel ou artisanal (menuisier, coiffeur, mécanicien).
\item \eng{a career} (nom) : une carrière, un parcours professionnel.
\item \eng{well paid} (adjectif) : bien payé.
\item \eng{to design a poster} (verbe) : concevoir une affiche.
\end{itemize}
\end{definition}

\courssub{Consignes}

Travail par binômes, en quatre étapes guidées.

\begin{enumerate}
\item \textbf{Préparation des rôles (10 min).} Choisissez chacun un métier dans la liste : \eng{carpenter, nurse, engineer, shopkeeper, tailor, farmer, teacher, mechanic, hairdresser, bank clerk}. L'élève A prépare sa fiche de professionnel : nom du métier, tâches quotidiennes, formation, qualités requises, avantages. L'élève B prépare au moins six questions d'interview en s'aidant du poster et des structures vues en classe (\eng{What does your job consist of? How long have you been working as...? What do you like most about...?}).

\item \textbf{L'interview (10 min).} Jouez la scène deux fois en inversant les rôles. Le professionnel doit employer \textbf{au moins cinq collocations} (par exemple \eng{to earn a living}, \eng{to work long hours}, \eng{to gain experience}, \eng{to run a business}, \eng{to take pride in one's work}) et \textbf{deux expressions idiomatiques} (par exemple \eng{to learn the ropes}, \eng{to climb the career ladder}, \eng{to get one's foot in the door}). Le journaliste prend des notes en anglais.

\item \textbf{L'affiche “Job of the Week” (20 min).} Ensemble, rédigez une affiche de 8 à 10 lignes présentant l'un des deux métiers. Elle doit contenir : le nom du métier, les tâches principales, la formation et les qualifications, les qualités requises, et un argument d'attraction. Utilisez au moins trois mots de la même famille (par exemple \eng{employ / employer / employee / employment} ou \eng{train / trainer / training}).

\item \textbf{Le vote (10 min).} Les affiches sont exposées au mur. Chaque binôme visite les affiches des autres, puis vote à voix haute : \eng{I am keen on this job because it is rewarding and well paid.} ou \eng{What attracts me is the chance to help people.} Justifiez toujours votre choix avec \eng{because}, \eng{since} ou \eng{as}.
\end{enumerate}

\begin{attention}[Pièges à éviter pendant l'activité]
\begin{itemize}
\item Ne dis pas \eng{I want to make my studies} : dis \eng{to do one's studies} ou \eng{to study}.
\item \eng{A formation} n'existe pas en anglais : dis \eng{training} ou \eng{qualifications}.
\item \eng{Actually} signifie « en fait », pas « actuellement » : dis \eng{at present} ou \eng{nowadays}.
\item Attention à la préposition : \eng{to be interested IN}, \eng{to be keen ON}, \eng{to apply FOR a job}, \eng{to depend ON}.
\end{itemize}
\end{attention}

\courssub{Ressources à mobiliser}

\begin{itemize}
\item \textbf{Lexique des métiers :} \eng{a trade, a profession, an occupation, a skilled worker, an apprentice, qualifications, training, wages, a salary, working conditions}.
\item \textbf{Exprimer l'intérêt :} \eng{to be interested in, to be keen on, to be fond of, to be passionate about, to dream of becoming, What attracts me is...}.
\item \textbf{Collocations clés :} \eng{to earn a living, to apply for a job, to gain experience, to work long hours, to run a business, to take pride in one's work, to look for a job}.
\item \textbf{Expressions idiomatiques :} \eng{to learn the ropes} (apprendre les ficelles du métier), \eng{to climb the career ladder} (gravir les échelons), \eng{to get one's foot in the door} (mettre un pied dans l'entreprise), \eng{a dead-end job} (un métier sans avenir).
\item \textbf{Familles de mots :} \eng{employ / employer / employee / employment / unemployed} ; \eng{train / trainer / trainee / training} ; \eng{profession / professional / professionally}.
\item \textbf{Grammaire :} questions avec \eng{do/does/did} et les mots interrogatifs \eng{what, why, how long, how much} ; le présent simple pour décrire les tâches habituelles ; \eng{since} et \eng{for} avec le present perfect (\eng{I have been a nurse for ten years}).
\end{itemize}

\courssub{Proposition de production et corrigé commenté}

\begin{exoresolu}[Interview et affiche modèles]
Rédigez l'interview et l'affiche “Job of the Week” pour le métier d'infirmière à l'hôpital régional de Bamenda.
\tcblower
\textbf{Interview (extrait modèle)}

\eng{Journalist: Good morning, Mrs Ngum. What does your job consist of?}

\eng{Nurse: Good morning. As a nurse, I look after patients, give them their medicine and help the doctors. I often work long hours, but I take pride in my work.}

\eng{Journalist: What qualifications do you need?}

\eng{Nurse: You need a nursing diploma and a lot of patience. When I started, I had to learn the ropes very quickly!}

\eng{Journalist: Is it well paid?}

\eng{Nurse: It is not the best-paid job in Cameroon, but I earn a living and the work is rewarding. Why did I choose this profession? Because I have always been keen on helping people.}

\eng{Journalist: Thank you very much, Mrs Ngum.}

\textbf{Affiche “Job of the Week” (modèle)}

\eng{JOB OF THE WEEK: NURSE}

\eng{A nurse looks after sick people in hospitals and health centres. He or she gives medicine, takes temperatures and comforts patients. To become a nurse, you need professional training and a diploma. A good nurse must be patient, kind and hard-working. Nurses often work long hours, but they earn a living and gain useful experience every day. It is a rewarding profession: you can really climb the career ladder and become a head nurse. Choose nursing if you are passionate about helping others!}

\textbf{Commentaires sur les choix de langue}
\begin{itemize}
\item Collocations réemployées : \eng{to look after patients}, \eng{to work long hours}, \eng{to take pride in one's work}, \eng{to earn a living}, \eng{to gain experience}, \eng{to climb the career ladder} : le contrat des cinq collocations est rempli.
\item Expression idiomatique : \eng{to learn the ropes}, bien placée dans une réponse orale naturelle.
\item Familles de mots exploitées : \eng{profession / professional}, \eng{train / training}, \eng{reward / rewarding}.
\item Grammaire : présent simple pour les tâches habituelles (\eng{I look after, I give}) ; questions correctement formées avec \eng{does... consist of} et \eng{do you need} ; prépositions correctes : \eng{keen on}, \eng{passionate about}.
\item Aucun faux ami : \eng{training} et non « formation », \eng{diploma} correctement employé.
\end{itemize}
\end{exoresolu}

\courssub{Bilan de l'activité}

Grâce à ce Career Day, tu as réemployé dans une situation de communication réelle tout le lexique de la leçon : les noms de métiers, les verbes de l'intérêt et de la vocation, les collocations du monde du travail et les familles de mots. Tu as aussi pratiqué la grammaire de l'interview (questions ouvertes et fermées) et la rédaction d'un texte court informatif et persuasif. Vérifie que tu sais maintenant :
\begin{itemize}
\item nommer au moins dix métiers et décrire leurs tâches en anglais ;
\item dire ce qui t'attire dans un métier avec \eng{to be keen on} et \eng{to be interested in} ;
\item employer cinq collocations et deux expressions idiomatiques sans hésiter ;
\item éviter les faux amis \eng{formation}, \eng{actually} et \eng{to make studies}.
\end{itemize}
Pour aller plus loin, rédige à la maison un court paragraphe (6 à 8 lignes) intitulé \eng{My dream job}, en réutilisant au moins quatre éléments du lexique de la leçon.

\newpage

\coursec{Méthodes et exercices résolus}

\courssub{Savoir-faire 1 : Identifier et employer le vocabulaire des métiers}

\begin{textebox}[Reading : Choosing a path]
At the career fair of Government High School Buea, many students discovered new occupations. Nadine, who has always been interested in electricity, wants to learn a trade and become an electrician like her uncle. Her friend Tabi dreams of becoming a journalist because he is keen on writing and reporting the news. ``Choosing a profession is not easy,'' says their teacher, ``but if you work hard and follow your passion, you will earn a good living and serve your community.''
\end{textebox}

\begin{definition}[Trade, job, profession, occupation, career]
\begin{itemize}
\item \cle{trade} : métier manuel ou artisanal, appris par la pratique (carpenter, tailor, plumber). Collocation : \eng{to learn a trade}.
\item \cle{job} : emploi concret, poste précis, souvent temporaire ou rémunéré (a part-time job).
\item \cle{profession} : métier qualifié exigeant des études supérieures (doctor, lawyer, teacher).
\item \cle{occupation} : terme général et soutenu, employé sur les formulaires administratifs.
\item \cle{career} : carrière, c'est-à-dire le parcours professionnel dans la durée.
\end{itemize}
\end{definition}

\begin{methode}[Identifier et employer le vocabulaire des métiers]
\begin{enumerate}
\item \textbf{Repérer le suffixe du métier.} La plupart des noms de métiers se forment avec un suffixe : \cle{-er}/\cle{-or} (teacher, actor), \cle{-ist} (dentist, journalist), \cle{-ian} (electrician, musician), \cle{-ant}/\cle{-ent} (accountant, student). Le suffixe indique souvent « celui qui fait l'action ».
\item \textbf{Classer le mot dans la bonne catégorie.} Distingue \cle{trade} (métier manuel/artisanal : carpenter, tailor), \cle{job} (emploi concret, poste), \cle{profession} (métier qualifié exigeant des études : doctor, lawyer) et \cle{occupation} (terme général, soutenu).
\item \textbf{Mémoriser le mot avec sa collocation.} N'apprends jamais un mot seul : \eng{to learn a trade}, \eng{to apply for a job}, \eng{to practise a profession}, \eng{to take up a career}.
\item \textbf{Vérifier l'article.} Pour dire le métier de quelqu'un, l'anglais exige l'article indéfini : \eng{She is a nurse.} (jamais \emph{She is nurse}).
\item \textbf{Réemployer dans une phrase personnelle.} Fabrique une phrase vraie sur toi ou ton entourage : \eng{My uncle is a carpenter in Bafoussam.}
\end{enumerate}
\end{methode}

\begin{exemplebox}[Quelques métiers fréquents et leur prononciation]
\begin{center}
\begin{tabularx}{\linewidth}{|X|X|X|}\hline
\rowcolor{popblueL} English & Français & Prononciation \\ \hline
a carpenter & un menuisier & « ka-peun-teur » \\ \hline
an electrician & un électricien & « è-lek-tri-cheun » \\ \hline
a tailor & un tailleur & « té-leur » \\ \hline
a nurse & une infirmière, un infirmier & « neu-seu » \\ \hline
an accountant & un comptable & « e-kaoun-teunt » \\ \hline
a journalist & un journaliste & « djeu-ne-liste » \\ \hline
a farmer & un agriculteur & « fa-meur » \\ \hline
a vet (veterinarian) & un vétérinaire & « vète » \\ \hline
\end{tabularx}
\end{center}
\end{exemplebox}

\begin{exoresolu}[Jobs, trades or professions ?]
\textbf{Énoncé.} Complete each sentence with the correct word from the box: \emph{trade, job, profession, occupation, career}.\\
1. Teaching is a noble \dots\ which requires patience and training.\\
2. My brother learnt the \dots\ of shoemaking in a workshop in Douala.\\
3. She found a part-time \dots\ as a waitress during the holidays.\\
4. On the form, write your name, age and \dots .\\
5. He had a brilliant \dots\ as a footballer before becoming a coach.
\tcblower
\textbf{Raisonnement et solution.}
\begin{enumerate}
\item \textbf{profession} — l'enseignement exige des études et une qualification ; on dit \eng{a noble profession}. (« Teaching is a noble profession which requires patience and training. »)
\item \textbf{trade} — la cordonnerie est un métier manuel, artisanal ; collocation : \eng{to learn a trade}. (« My brother learnt the trade of shoemaking in a workshop in Douala. »)
\item \textbf{job} — un emploi concret, ici temporaire ; collocation : \eng{a part-time job}. (« She found a part-time job as a waitress during the holidays. »)
\item \textbf{occupation} — sur un formulaire administratif, le terme officiel est \eng{occupation}. (« On the form, write your name, age and occupation. »)
\item \textbf{career} — une carrière, c'est le parcours professionnel dans la durée ; collocation : \eng{a brilliant career}. (« He had a brilliant career as a footballer before becoming a coach. »)
\end{enumerate}
\textbf{Conclusion.} Le choix dépend du sens : \emph{profession} = métier qualifié, \emph{trade} = métier manuel, \emph{job} = poste précis, \emph{occupation} = terme administratif, \emph{career} = parcours dans le temps.
\end{exoresolu}

\courssub{Savoir-faire 2 : Exprimer l'intérêt et la vocation}

\begin{propriete}[Les structures de l'intérêt : préposition + verbe en -ing]
Après une préposition (\emph{in, on, of, about}), le verbe anglais prend toujours la forme en \cle{-ing} (gérondif) :
\begin{itemize}
\item \eng{to be interested in + nom / V-ing} : \eng{I am interested in nursing.} (« Je m'intéresse au métier d'infirmier. »)
\item \eng{to be keen on + V-ing} : \eng{She is keen on becoming an engineer.} (« Elle tient à devenir ingénieure. »)
\item \eng{to be passionate about + nom / V-ing} : \eng{He is passionate about farming.} (« Il est passionné d'agriculture. »)
\item \eng{to be fond of + nom / V-ing} : \eng{They are fond of teaching.} (« Ils aiment enseigner. »)
\item \eng{to dream of + V-ing} : \eng{He dreams of becoming a pilot.} (« Il rêve de devenir pilote. »)
\end{itemize}
\end{propriete}

\begin{methode}[Developing an interest : exprimer l'intérêt pour un métier]
\begin{enumerate}
\item \textbf{Choisir la bonne structure.} Trois structures à maîtriser : \eng{to be interested in}, \eng{to be keen on}, \eng{to dream of becoming} : \eng{He dreams of becoming a pilot.}
\item \textbf{Respecter la préposition.} \cle{interested in}, \cle{keen on}, \cle{passionate about}, \cle{fond of}. Après ces prépositions, le verbe est toujours en \cle{-ing}.
\item \textbf{Exprimer l'origine de la vocation.} \eng{Since I was a child, I have always wanted to be a doctor.} / \eng{My interest in farming grew when I visited my grandmother's farm in Bafut.}
\item \textbf{Exprimer le développement de l'intérêt.} \eng{to develop an interest in}, \eng{to grow fond of}, \eng{to discover a passion for} : \eng{During the holidays, I developed an interest in tailoring.}
\item \textbf{Conclure avec un projet.} \eng{That is why I intend to study medicine at university.}
\end{enumerate}
\end{methode}

\begin{exoresolu}[Express your interest]
\textbf{Énoncé.} Rewrite each sentence using the word in brackets without changing the meaning.\\
1. I like computers very much and I want to work in this field. (keen)\\
2. She started to like carpentry when she helped her father. (developed)\\
3. He wants to become a journalist. (dreams)\\
4. Since my childhood, I have always loved animals. That is why I want to be a vet. (interested)
\tcblower
\textbf{Raisonnement et solution.}
\begin{enumerate}
\item La structure \eng{to be keen on + V-ing} remplace « like very much » : \eng{I am keen on working with computers.} Attention : \emph{on} + gérondif, jamais l'infinitif.
\item \eng{to develop an interest in + nom} exprime le début progressif d'un goût : \eng{She developed an interest in carpentry when she helped her father.}
\item \eng{to dream of + V-ing} : \eng{He dreams of becoming a journalist.} La préposition \emph{of} est obligatoire et impose le gérondif \emph{becoming}.
\item \eng{to be interested in + nom} : \eng{I have always been interested in animals, so I want to become a vet.} On garde le present perfect (\emph{have always been}) car l'intérêt commence dans le passé et continue aujourd'hui.
\end{enumerate}
\textbf{Conclusion.} Pour exprimer l'intérêt, retiens le trio : \emph{interested in}, \emph{keen on}, \emph{dream of}, toujours suivis d'un nom ou d'un verbe en \emph{-ing}.
\end{exoresolu}

\courssub{Savoir-faire 3 : Échanger sur les métiers à l'oral}

\begin{methode}[Sharing an interest : parler des métiers en dialogue]
\begin{enumerate}
\item \textbf{Poser la question du métier correctement.} \eng{What do you do?} (présent, pour un adulte) ; \eng{What do you want to be?} ou \eng{What would you like to do in the future?} (pour un élève). Évite \emph{What is your job?} en conversation directe : c'est abrupt et administratif.
\item \textbf{Répondre avec une structure complète.} \eng{I want to be a tailor because I love designing clothes.} Toujours justifier avec \cle{because} ou \cle{since}.
\item \textbf{Demander des précisions.} \eng{Why did you choose that trade?} / \eng{How long does the training take?} / \eng{Is it difficult to find a job in that field?}
\item \textbf{Réagir et encourager.} \eng{That sounds interesting!} / \eng{That's a wonderful choice!} / \eng{I am sure you will succeed.}
\item \textbf{Parler des qualités requises.} \eng{To be a good nurse, you need patience and kindness.} Structure : \cle{you need + nom} ou \cle{it requires + nom}.
\end{enumerate}
\end{methode}

\begin{exoresolu}[A dialogue at the career fair]
\textbf{Énoncé.} Complete the dialogue between two students, Amina and Peter, with appropriate questions or answers.\\
Amina : (1) \dots\ in the future, Peter?\\
Peter : I would like to become an electrician.\\
Amina : (2) \dots\ that trade?\\
Peter : Because I have always been interested in electricity, and my uncle has a workshop in Bamenda.\\
Amina : (3) \dots\ the training take?\\
Peter : About three years in a technical school.\\
Amina : (4) \dots\ ! I am sure you will be a great electrician.
\tcblower
\textbf{Raisonnement et solution.}
\begin{enumerate}
\item La réponse « I would like to become an electrician » indique une question sur le projet d'avenir : \eng{What would you like to do} (ou \eng{What do you want to be}) \eng{in the future, Peter?}
\item La réponse commence par « Because », donc la question porte sur la raison : \eng{Why did you choose that trade?} (prétérit, car le choix est fait).
\item La réponse « About three years » indique une durée : \eng{How long does the training take?} Attention à l'auxiliaire \emph{does} et à la base verbale \emph{take}.
\item Une réaction positive de conclusion : \eng{That sounds interesting!} ou \eng{That's a wonderful choice!}
\end{enumerate}
\textbf{Conclusion.} Dans un dialogue sur les métiers, chaque réponse donne l'indice de la question : « because » appelle \emph{why}, une durée appelle \emph{how long}, un projet appelle \emph{what would you like to do}.
\end{exoresolu}

\begin{experience}[Role play : at the career guidance office]
Par groupes de deux : l'élève A est conseiller d'orientation (\eng{careers adviser}), l'élève B est un élève de Première qui hésite entre deux métiers. B explique ses centres d'intérêt avec \eng{I am interested in...}, \eng{I am keen on...} ; A pose au moins trois questions (\eng{Why...?}, \eng{How long...?}, \eng{What qualities do you need...?}) puis conseille un métier avec \eng{I think you should...} ou \eng{Have you thought of becoming a...?}. Échangez ensuite les rôles.
\end{experience}

\courssub{Savoir-faire 4 : Employer les collocations du monde du travail}

\begin{methode}[Collocations : les couples de mots qui vont ensemble]
\begin{enumerate}
\item \textbf{Apprendre les verbes qui accompagnent « job ».} \eng{to look for a job}, \eng{to find a job}, \eng{to apply for a job}, \eng{to get a job}, \eng{to lose one's job}, \eng{to do a job} (accomplir une tâche).
\item \textbf{Apprendre les verbes qui accompagnent « career » et « trade ».} \eng{to pursue a career}, \eng{to build a career}, \eng{to make a career in}, \eng{to learn a trade}, \eng{to master a trade}.
\item \textbf{Employer les expressions du travail.} \eng{to work hard}, \eng{to earn a living} (gagner sa vie), \eng{to be unemployed} (être au chômage), \eng{to be self-employed} (travailler à son compte), \eng{to retire} (prendre sa retraite).
\item \textbf{Éviter les calques.} On ne dit pas \emph{win one's life} mais \eng{earn one's living} ; on ne dit pas \emph{make a job} mais \eng{do a job} ou \eng{have a job}.
\item \textbf{Tester la collocation dans une phrase.} \eng{After his training, Samuel applied for a job at the Douala port.}
\end{enumerate}
\end{methode}

\begin{exoresolu}[Choose the right verb]
\textbf{Énoncé.} Choose the correct verb to complete each sentence.\\
1. Many young people (win / earn) their living as motorbike taxi drivers in Yaoundé.\\
2. After secondary school, she wants to (make / learn) a trade.\\
3. He (applied for / applied to) a job at the hospital last week.\\
4. My grandfather (retreated / retired) last year after forty years of teaching.\\
5. It is not easy to (do / make) a career in music.
\tcblower
\textbf{Raisonnement et solution.}
\begin{enumerate}
\item \textbf{earn} — la collocation correcte est \eng{to earn one's living} (gagner sa vie) ; \emph{win one's life} est un calque du français « gagner sa vie ». (« Many young people earn their living as motorbike taxi drivers in Yaoundé. »)
\item \textbf{learn} — collocation : \eng{to learn a trade} (apprendre un métier) ; \emph{make a trade} n'existe pas dans ce sens.
\item \textbf{applied for} — on postule \emph{à} un emploi avec \eng{to apply for a job} ; \emph{apply to} s'emploie pour une institution (\eng{to apply to a university}).
\item \textbf{retired} — \eng{to retire} = prendre sa retraite ; \emph{to retreat} signifie « battre en retraite, se retirer » (sens militaire). Faux ami à éviter.
\item \textbf{make} — collocation : \eng{to make a career in} (faire carrière dans). On dit aussi \eng{to pursue a career}.
\end{enumerate}
\textbf{Conclusion.} Les collocations ne se traduisent pas mot à mot : apprends chaque verbe avec son nom partenaire (\emph{earn a living}, \emph{learn a trade}, \emph{apply for a job}, \emph{retire}, \emph{make a career}).
\end{exoresolu}

\courssub{Savoir-faire 5 : Éviter les faux amis et réemployer les familles de mots}

\begin{attention}[Faux amis du monde du travail]
\begin{center}
\begin{tabularx}{\linewidth}{|X|X|X|}\hline
\rowcolor{popblueL} Mot anglais & Sens réel & Piège français \\ \hline
actually & en fait, en réalité & « actuellement » = \emph{currently}, \emph{at present} \\ \hline
training & formation (professionnelle) & \emph{formation} = structure, disposition \\ \hline
to retire & prendre sa retraite & \emph{to retreat} = battre en retraite \\ \hline
a vocation & un appel profond & \emph{vacation} = vacances (américain) \\ \hline
a degree & un diplôme universitaire & « degré » = \emph{degree} uniquement au sens de mesure \\ \hline
to achieve & accomplir, réaliser & « acheter » = \emph{to buy} \\ \hline
\end{tabularx}
\end{center}
\end{attention}

\begin{methode}[Faux amis et familles de mots]
\begin{enumerate}
\item \textbf{Méfie-toi des mots qui ressemblent au français.} Pour « formation professionnelle », dis \eng{training} ou \eng{vocational training} ; \eng{actually} = « en fait », pas « actuellement » (\eng{currently}).
\item \textbf{Profession vs vocation vs vacation.} \eng{profession} = métier ; \eng{vocation} = appel profond ; \eng{vacation} = vacances (anglais américain). Trois mots à ne pas confondre.
\item \textbf{Construis les familles de mots.} À partir d'un verbe, trouve le nom de la personne et le nom de l'activité : \emph{to teach} $\rightarrow$ \emph{teacher} (personne) $\rightarrow$ \emph{teaching} (activité) ; \emph{to farm} $\rightarrow$ \emph{farmer} $\rightarrow$ \emph{farming} ; \emph{to build} $\rightarrow$ \emph{builder} $\rightarrow$ \emph{building}.
\item \textbf{Attention à l'orthographe des dérivés.} \emph{to employ} $\rightarrow$ \emph{employer} (employeur) / \emph{employee} (employé) / \emph{employment} ; \emph{to apply} $\rightarrow$ \emph{application} (le \emph{y} devient \emph{i} devant le suffixe).
\item \textbf{Vérifie la prononciation des suffixes.} \emph{-er}, \emph{-or}, \emph{-ist} sont atones : \emph{teacher} se prononce « ti-tcheur », \emph{doctor} « do-kteur », \emph{artist} « a-tiste ».
\end{enumerate}
\end{methode}

\begin{exoresolu}[Word families and false friends]
\textbf{Énoncé.} A. Complete the table with the missing words.\\[2pt]
\begin{tabular}{|l|l|l|}\hline
Verb & Person & Activity \\ \hline
to teach & \dots\ & teaching \\ \hline
to farm & farmer & \dots\ \\ \hline
\dots\ & actor & acting \\ \hline
to employ & employer / \dots\ & employment \\ \hline
\end{tabular}\\[4pt]
B. Correct the false friend in each sentence.\\
1. I want to follow a formation in computer science.\\
2. My father is actually in Douala for business. (meaning: at the moment)
\tcblower
\textbf{Raisonnement et solution.}\\
\textbf{A.} 1. \textbf{teacher} — suffixe \emph{-er} ajouté au verbe \emph{teach}. 2. \textbf{farming} — le nom de l'activité se forme en \emph{-ing} : \emph{farming} (l'agriculture). 3. \textbf{to act} — le verbe de base de la famille \emph{actor / acting}. 4. \textbf{employee} — l'employeur est \emph{employer} (celui qui emploie), l'employé est \emph{employee} (celui qui est employé) ; ne pas inverser !\\[4pt]
\textbf{B.} 1. \eng{I want to take a training course in computer science.} — « formation » au sens d'études se dit \emph{training (course)} ; \emph{formation} en anglais désigne une structure ou une disposition. 2. \eng{My father is currently in Douala for business.} — \emph{actually} signifie « en fait, en réalité » ; « actuellement » se traduit \emph{currently} ou \emph{at present}.\\
\textbf{Conclusion.} Pour chaque famille de mots, retiens les trois colonnes : verbe, personne, activité. Et face à un mot qui ressemble au français, vérifie toujours son sens réel en anglais.
\end{exoresolu}

\begin{savaistu}[Apprenticeship in Cameroon and abroad]
Au Cameroun, beaucoup de jeunes apprennent un métier par l'apprentissage informel chez un maître artisan : menuiserie, couture, mécanique. En anglais, l'apprenti se dit \eng{an apprentice} et l'apprentissage \eng{an apprenticeship}. Dans les pays anglophones, les \eng{vocational schools} (écoles de formation professionnelle) et les \eng{career fairs} (forums des métiers) aident les élèves à choisir leur voie, exactement comme le career fair du texte de Buea.
\end{savaistu}

\begin{aretenir}[L'essentiel de la leçon]
\begin{itemize}
\item \emph{trade} = métier manuel, \emph{job} = poste, \emph{profession} = métier qualifié, \emph{occupation} = terme administratif, \emph{career} = parcours professionnel.
\item Intérêt : \eng{interested in}, \eng{keen on}, \eng{passionate about}, \eng{fond of}, \eng{dream of} + nom ou verbe en \emph{-ing}.
\item Collocations clés : \eng{learn a trade}, \eng{apply for a job}, \eng{earn one's living}, \eng{make a career}, \eng{retire}.
\item Familles de mots : verbe $\rightarrow$ personne $\rightarrow$ activité (\emph{teach / teacher / teaching}).
\item Faux amis : \emph{actually} $\neq$ « actuellement » ; \emph{training} = « formation » ; \emph{retire} $\neq$ \emph{retreat}.
\end{itemize}
\end{aretenir}

\begin{attention}[Les 5 erreurs qui coûtent des points au Bac]
\begin{enumerate}
\item \textbf{Oublier l'article devant le métier.} Écris \eng{She is a teacher.} et jamais \emph{She is teacher}. En anglais, tout nom de métier au singulier prend \emph{a/an}.
\item \textbf{Traduire « gagner sa vie » par \emph{win one's life}.} La bonne expression est \eng{to earn one's living}. \emph{Win} s'emploie pour un match, un prix, une victoire.
\item \textbf{Mettre l'infinitif après une préposition.} Écris \eng{I am interested in becoming a doctor.} et jamais \emph{interested in to become}. Après \emph{in, on, of, about}, le verbe prend toujours \emph{-ing}.
\item \textbf{Confondre les faux amis.} \emph{actually} = « en fait » (pas « actuellement » : \emph{currently}) ; \emph{training} = « formation » (pas \emph{formation}) ; \emph{to retire} = « prendre sa retraite » (pas \emph{to retreat}).
\item \textbf{Inverser employer et employee.} L'\emph{employer} est le patron (celui qui emploie), l'\emph{employee} est le salarié (celui qui est employé). Une inversion change complètement le sens de la phrase.
\end{enumerate}
\end{attention}

\begin{exercice}[Pour s'entraîner]
1. Complete with a word from the lesson: \eng{After his \dots\ (training) as a mechanic, Ali \dots\ (applied for) a \dots\ in a garage in Douala. He wants to \dots\ (earn) his living and later be \dots\ (self-employed).}\\
2. Rewrite with the word in brackets: \eng{She loves children, so she wants to be a teacher.} (keen)\\
3. Translate into English: « Mon frère s'intéresse à la menuiserie et rêve d'ouvrir son atelier à Bamenda. »\\
4. Find the person and the activity: \emph{to drive}, \emph{to sing}, \emph{to write}.
\end{exercice}

\begin{corrige}[Pour s'entraîner]
1. \eng{After his training as a mechanic, Ali applied for a job in a garage in Douala. He wants to earn his living and later be self-employed.}\\
2. \eng{She is keen on working with children, so she wants to be a teacher.} (\emph{keen on} + gérondif).\\
3. \eng{My brother is interested in carpentry and dreams of opening his workshop in Bamenda.} (\emph{interested in} + nom ; \emph{dreams of} + \emph{-ing}).\\
4. \emph{to drive} $\rightarrow$ \emph{driver} / \emph{driving} ; \emph{to sing} $\rightarrow$ \emph{singer} / \emph{singing} ; \emph{to write} $\rightarrow$ \emph{writer} / \emph{writing}.
\end{corrige}

\newpage

\coursec{Exercices}

\courssub{Je vérifie mes connaissances}

\begin{exercice}[QCM : le bon mot]
Choisis la réponse correcte (a, b, c ou d) pour compléter chaque phrase.

\begin{enumerate}
\item \eng{My brother is learning plumbing; he is doing an \_\_\_\_\_\_ in a workshop.}\\
(a) internship \quad (b) stage \quad (c) formation \quad (d) career
\item \eng{She works for herself; she is \_\_\_\_\_\_.}\\
(a) employed \quad (b) self-employed \quad (c) jobless \quad (d) retired
\item \eng{After three years of \_\_\_\_\_\_, he became a qualified electrician.}\\
(a) formation \quad (b) training \quad (c) information \quad (d) education
\item \eng{Many young people in Douala \_\_\_\_\_\_ as motorcycle-taxi drivers.}\\
(a) win their life \quad (b) earn their living \quad (c) gain their life \quad (d) make their life
\item \eng{I am very \_\_\_\_\_\_ computer science; I want to become a programmer.}\\
(a) interested by \quad (b) keen on \quad (c) fond to \quad (d) attracted with
\item \eng{You need a \_\_\_\_\_\_ to practise as a nurse in Cameroon.}\\
(a) qualification \quad (b) quality \quad (c) quantity \quad (d) question
\end{enumerate}
\end{exercice}

\begin{exercice}[Vrai ou faux ? Justifie ta réponse]
Lis le texte suivant, puis dis si chaque affirmation est \eng{true} ou \eng{false}. Justifie chaque réponse par une courte citation du texte.

\begin{textebox}[A nurse from Bamenda]
\eng{Grace grew up in Bamenda. As a child, she was fascinated by the nurses at the local hospital. After her Advanced Level, she applied for a place in a nursing school in Buea. The training lasted three years and was very demanding, but Grace never gave up. She obtained her qualification in 2019 and was immediately hired by a private clinic. Today she earns a decent living and hopes to open her own health centre one day. She often tells students: “Choose a job you love, and work hard for it.”}
\end{textebox}

\begin{enumerate}
\item \eng{Grace decided to become a nurse when she was an adult.}
\item \eng{She studied nursing in Bamenda.}
\item \eng{Her training was easy.}
\item \eng{She found a job as soon as she got her qualification.}
\item \eng{She is self-employed at the moment.}
\item \eng{She advises young people to choose a job they are passionate about.}
\end{enumerate}
\end{exercice}

\begin{exercice}[Vocabulaire : devine le métier]
Pour chaque définition, trouve le métier correspondant dans la liste : \eng{carpenter, tailor, midwife, mechanic, pharmacist, journalist}.

\begin{enumerate}
\item \eng{A person who repairs cars and engines.}
\item \eng{A person who makes or mends clothes.}
\item \eng{A person who makes furniture from wood.}
\item \eng{A person who helps women when their babies are born.}
\item \eng{A person who prepares and sells medicines.}
\item \eng{A person who writes articles for newspapers or reports the news on the radio.}
\end{enumerate}
\end{exercice}

\begin{exercice}[Conjugaison à compléter]
Complète chaque phrase avec le verbe entre parenthèses au temps et à la forme qui conviennent.

\begin{enumerate}
\item \eng{My uncle \_\_\_\_\_\_ (work) as a welder in Edea for twenty years now.}
\item \eng{Last year, Amina \_\_\_\_\_\_ (apply) for a job at the town council.}
\item \eng{When I finish my training, I \_\_\_\_\_\_ (open) my own workshop.}
\item \eng{While the apprentices \_\_\_\_\_\_ (learn) their trade, the master was watching them.}
\item \eng{He \_\_\_\_\_\_ (not / be) interested in farming when he was younger.}
\item \eng{Look! The tailor \_\_\_\_\_\_ (sew) a beautiful dress for the wedding.}
\end{enumerate}
\end{exercice}

\courssub{Je m'entraîne}

\begin{exercice}[Traduction]
Traduis les phrases suivantes en anglais, en utilisant le vocabulaire étudié dans la leçon.

\begin{enumerate}
\item Mon frère fait un apprentissage chez un menuisier à Nkongsamba.
\item Elle gagne sa vie comme coiffeuse depuis cinq ans.
\item Je suis passionné par la mécanique et je voudrais en faire mon métier.
\item Il faut une formation sérieuse pour devenir infirmier.
\item Beaucoup de jeunes Camerounais rêvent de travailler à leur propre compte.
\item Quel métier t'intéresse le plus, et pourquoi ?
\end{enumerate}
\end{exercice}

\begin{exercice}[Texte à trous : un apprenti de Nkongsamba]
Complète le texte suivant avec les dix mots de la banque. Chaque mot ne peut être utilisé qu'une seule fois.

\begin{center}
\begin{tabular}{|l|l|l|l|l|}
\hline
\rowcolor{popblueL} \eng{trade} & \eng{training} & \eng{skilled} & \eng{earn a living} & \eng{apply for} \\
\hline
\eng{keen on} & \eng{employer} & \eng{qualification} & \eng{self-employed} & \eng{career} \\
\hline
\end{tabular}
\end{center}

\begin{textebox}[Junior's choice]
\eng{Junior is eighteen and lives in Nkongsamba. He has always been (1) \_\_\_\_\_\_ electricity, so after his BEPC he decided to learn a (2) \_\_\_\_\_\_ instead of going to high school. He joined a workshop where he is receiving practical (3) \_\_\_\_\_\_ from a (4) \_\_\_\_\_\_ electrician. His (5) \_\_\_\_\_\_ is satisfied with his work and pays him a small allowance. In two years, Junior hopes to obtain a (6) \_\_\_\_\_\_ that will allow him to (7) \_\_\_\_\_\_ jobs in big companies. Later, he would like to be (8) \_\_\_\_\_\_ and open his own business, because he wants to (9) \_\_\_\_\_\_ honestly and build a successful (10) \_\_\_\_\_\_.}
\end{textebox}
\end{exercice}

\begin{exercice}[Correction d'erreurs de francophones]
Chaque phrase contient une erreur typique d'un élève francophone (faux ami, calque ou mauvaise préposition). Souligne l'erreur et réécris la phrase correctement.

\begin{enumerate}
\item \eng{I am doing a formation to become a nurse.}
\item \eng{He won his life as a tailor in Bafoussam.}
\item \eng{My sister is very interested by medicine.}
\item \eng{She works at her own account.}
\item \eng{I want to make the same job as my father.}
\item \eng{He learned me how to repair engines.}
\item \eng{After my studies, I will search a job in Douala.}
\item \eng{The carpenters are fabricating chairs in the workshop.}
\end{enumerate}
\end{exercice}

\begin{exercice}[Appariement : dialogue après le \eng{career day}]
Deux élèves discutent après la journée des métiers organisée par leur lycée. Remets les répliques de B dans le bon ordre pour reconstituer le dialogue. Écris simplement les lettres dans l'ordre.

\begin{itemize}
\item \eng{A: Did you enjoy the career day yesterday?}
\item \eng{B (1): \_\_\_\_\_\_}
\item \eng{A: Which profession attracted you most?}
\item \eng{B (2): \_\_\_\_\_\_}
\item \eng{A: Really? Why that one?}
\item \eng{B (3): \_\_\_\_\_\_}
\item \eng{A: And what training does it require?}
\item \eng{B (4): \_\_\_\_\_\_}
\item \eng{A: That sounds great. Good luck with your project!}
\item \eng{B (5): \_\_\_\_\_\_}
\end{itemize}

Répliques à classer :
\begin{itemize}
\item[(a)] \eng{Because I love engines, and skilled mechanics earn a good living here.}
\item[(b)] \eng{Thank you! And don't forget to tell me about your own choice.}
\item[(c)] \eng{Yes, very much! It was so interesting to meet real professionals.}
\item[(d)] \eng{You need about three years of training in a technical school, plus an apprenticeship.}
\item[(e)] \eng{I think I was most attracted by car mechanics.}
\end{itemize}
\end{exercice}

\courssub{Je me prépare au Bac}

\begin{exercice}[Compréhension écrite et questions de langue]
Lis attentivement le texte suivant, puis réponds aux questions en anglais, avec des phrases complètes.

\begin{textebox}[From the classroom to the workshop]
\eng{In many Cameroonian towns, a quiet revolution is taking place. More and more young people are turning to technical trades such as plumbing, welding, tailoring and hairdressing. For a long time, parents believed that success only meant becoming a doctor, a lawyer or a civil servant. Today, attitudes are changing.}

\eng{Take the example of Sandra, a nineteen-year-old girl from Mbanga. After her Probatoire, she told her parents that she wanted to become a plumber. Her father was shocked: “That is a man's job!” he said. But Sandra was determined. She enrolled in a technical college, completed a six-month internship with a building company in Douala, and obtained her qualification last year. She now works on construction sites and earns more money than some of her friends who went to university.}

\eng{Experts say that skilled workers are badly needed in Cameroon. A good electrician or a reliable mechanic can build a solid career, become self-employed and even create jobs for others. As Sandra puts it: “A trade in your hands is better than a diploma in your pocket.”}
\end{textebox}

\textbf{A. Comprehension questions}
\begin{enumerate}
\item What change is happening among young Cameroonians, according to the text?
\item What did parents traditionally consider as signs of success? Give two examples.
\item Why was Sandra's father shocked?
\item Mention the three steps of Sandra's training.
\item Explain in your own words the meaning of Sandra's last sentence.
\end{enumerate}

\textbf{B. Language work}
\begin{enumerate}
\item Find in the text: (a) two words from the same family as \eng{employ}; (b) one synonym of \eng{job}; (c) one adjective meaning \eng{qualified and experienced}.
\item Put the verbs in brackets in the correct tense: \eng{Sandra (enrol) in a technical college after she (tell) her parents about her project.}
\item Turn into reported speech: \eng{Her father said: “That is a man's job!”}
\end{enumerate}
\end{exercice}

\begin{exercice}[Traduction]
Traduis le passage suivant en anglais. Soigne le choix du vocabulaire des métiers et évite les calques du français.

\begin{textebox}[Texte à traduire]
Mon cousin Étienne a toujours été passionné par la menuiserie. Après le BEPC, il est entré comme apprenti dans l'atelier d'un artisan de Bafoussam. Sa formation a duré quatre ans, mais aujourd'hui il est très habile de ses mains. Il fabrique des meubles qu'il vend au marché, et il gagne bien sa vie. Son rêve est d'ouvrir sa propre entreprise et de former à son tour de jeunes apprentis. Je suis fier de lui, car il a choisi un métier qu'il aime vraiment.
\end{textebox}
\end{exercice}

\begin{exercice}[Composition type Baccalauréat]
Rédige un paragraphe en anglais de \textbf{80 à 100 mots} sur le sujet suivant :

\begin{textebox}[Topic]
\eng{Write a paragraph about the job you would like to do in the future. Say why you are interested in it, what training it requires and what qualities are needed.}
\end{textebox}

Conseils méthodologiques :
\begin{enumerate}
\item Annonce clairement le métier choisi dès la première phrase (\eng{In the future, I would like to become...}).
\item Donne au moins deux raisons de ton intérêt (\eng{I have always been keen on... / What attracts me most is...}).
\item Précise la formation nécessaire (\eng{This job requires... years of training / an apprenticeship / a qualification in...}).
\item Termine par les qualités indispensables (\eng{To succeed in this profession, one needs to be patient, skilled and hard-working.}).
\item Vérifie l'orthographe, les temps verbaux et le nombre de mots.
\end{enumerate}
\end{exercice}



\newpage

\coursec{Corrigés des exercices}

\coursec{Lesson 24 — Vocabulary: Words and expressions related to developing and sharing interest in a trade/job/profession}

\courssub{Je vérifie mes connaissances}

\begin{corrige}[Exercice 1 — QCM : le bon mot]
\begin{enumerate}
\item \textbf{(a) \cle{internship}}. \eng{He is doing an internship in a workshop.} Faux amis : \eng{a stage} désigne une scène de théâtre ou une étape ; \eng{formation} ne s'emploie pas pour une formation professionnelle ; \eng{a career} désigne une carrière entière, pas un stage.
\item \textbf{(b) \cle{self-employed}}. \eng{She is self-employed.} « À son propre compte ». \eng{Employed} = salarié ; \eng{jobless} = sans emploi ; \eng{retired} = à la retraite.
\item \textbf{(b) \cle{training}}. \eng{After three years of training, he became a qualified electrician.} \eng{Education} désigne l'éducation générale, pas une formation technique.
\item \textbf{(b) \cle{earn their living}}. \eng{Many young people in Douala earn their living as motorcycle-taxi drivers.} « Gagner sa vie » = \eng{to earn one's living} ; \eng{win one's life} est un calque du français.
\item \textbf{(b) \cle{keen on}}. \eng{I am very keen on computer science.} Prépositions correctes : \eng{interested in}, \eng{fond of}, \eng{keen on}.
\item \textbf{(a) \cle{qualification}}. \eng{You need a qualification to practise as a nurse in Cameroon.} \eng{Quality} = qualité ; \eng{quantity} = quantité.
\end{enumerate}
\end{corrige}

\begin{corrige}[Exercice 2 — Vrai ou faux ?]
\begin{enumerate}
\item \textbf{False.} \eng{“As a child, she was fascinated by the nurses at the local hospital.”} — Elle a choisi ce métier enfant, pas adulte.
\item \textbf{False.} \eng{“She applied for a place in a nursing school in Buea.”} — C'est à Buea, pas à Bamenda.
\item \textbf{False.} \eng{“The training lasted three years and was very demanding.”} — La formation était exigeante, pas facile.
\item \textbf{True.} \eng{“She obtained her qualification in 2019 and was immediately hired by a private clinic.”} — \eng{immediately} = « aussitôt, dès que ».
\item \textbf{False.} Elle est salariée d'une clinique privée ; elle \emph{espère} seulement ouvrir son propre centre : \eng{“She hopes to open her own health centre one day.”}
\item \textbf{True.} \eng{“Choose a job you love, and work hard for it.”} — Elle conseille de choisir un métier qu'on aime.
\end{enumerate}
\end{corrige}

\begin{corrige}[Exercice 3 — Devine le métier]
\begin{enumerate}
\item \eng{a mechanic} (un mécanicien)
\item \eng{a tailor} (un tailleur)
\item \eng{a carpenter} (un menuisier-charpentier)
\item \eng{a midwife} (une sage-femme ; pluriel : \eng{midwives})
\item \eng{a pharmacist} (un pharmacien)
\item \eng{a journalist} (un journaliste)
\end{enumerate}
\end{corrige}

\begin{corrige}[Exercice 4 — Conjugaison à compléter]
\begin{enumerate}
\item \eng{has worked} (ou \eng{has been working}) : present perfect, car l'action a commencé dans le passé et dure encore (\eng{for twenty years now}).
\item \eng{applied} : prétérit simple, marqué par \eng{last year}. Orthographe : \eng{apply} $\rightarrow$ \eng{applied} (y $\rightarrow$ i + ed).
\item \eng{will open} (ou \eng{am going to open}) : futur. Attention : après \eng{when}, on garde le présent (\eng{when I finish}, jamais \eng{when I will finish}).
\item \eng{were learning} : prétérit continu, action longue en arrière-plan (\eng{while}).
\item \eng{was not} (ou \eng{wasn't}) : prétérit de \eng{be} à la forme négative.
\item \eng{is sewing} : présent continu, annoncé par \eng{Look!}
\end{enumerate}
\end{corrige}

\courssub{Je m'entraîne}

\begin{corrige}[Exercice 5 — Traduction]
\begin{enumerate}
\item \eng{My brother is doing an apprenticeship with a carpenter in Nkongsamba.} (ou \eng{is an apprentice to a carpenter})
\item \eng{She has been earning her living as a hairdresser for five years.} (present perfect continu avec \eng{for five years})
\item \eng{I am keen on mechanics and I would like to make it my trade.} (ou \eng{to make a career of it})
\item \eng{You need serious training to become a nurse.} (ou \eng{Proper training is required to become a nurse.})
\item \eng{Many young Cameroonians dream of being self-employed.} (ou \eng{of working for themselves})
\item \eng{Which job are you most interested in, and why?} (ou \eng{Which trade interests you most, and why?})
\end{enumerate}
Points de vigilance : « gagner sa vie » = \eng{earn one's living} ; « à leur propre compte » = \eng{self-employed} ; « passionné par » = \eng{keen on} (jamais \eng{interested by}).
\end{corrige}

\begin{corrige}[Exercice 6 — Texte à trous : Junior's choice]
\begin{enumerate}
\item \eng{keen on} \quad 2. \eng{trade} \quad 3. \eng{training} \quad 4. \eng{skilled} \quad 5. \eng{employer}
\item[6.] \eng{qualification} \quad 7. \eng{apply for} \quad 8. \eng{self-employed} \quad 9. \eng{earn a living} \quad 10. \eng{career}
\end{enumerate}
Texte reconstitué : \eng{Junior is eighteen and lives in Nkongsamba. He has always been keen on electricity, so after his BEPC he decided to learn a trade instead of going to high school. He joined a workshop where he is receiving practical training from a skilled electrician. His employer is satisfied with his work and pays him a small allowance. In two years, Junior hopes to obtain a qualification that will allow him to apply for jobs in big companies. Later, he would like to be self-employed and open his own business, because he wants to earn a living honestly and build a successful career.}\\
Astuce de méthode : repère les prépositions et articles qui encadrent le trou (\eng{been \_\_ electricity} appelle un adjectif suivi de \eng{on} ; \eng{a \_\_} appelle un nom comptable singulier ; \eng{to \_\_ jobs} appelle un verbe suivi de \eng{for}).
\end{corrige}

\begin{corrige}[Exercice 7 — Correction d'erreurs de francophones]
\begin{enumerate}
\item \eng{I am doing a \cle{training course} to become a nurse.} — \eng{formation} est un faux ami ; on dit \eng{training} ou \eng{a training course}.
\item \eng{He \cle{earns his living} as a tailor in Bafoussam.} — « Gagner sa vie » = \eng{to earn one's living}, jamais \eng{win one's life}.
\item \eng{My sister is very interested \cle{in} medicine.} — La préposition correcte est \eng{in} (ou \eng{keen on}).
\item \eng{She \cle{is self-employed}.} (ou \eng{works for herself}) — « À son compte » ne se traduit pas mot à mot.
\item \eng{I want to \cle{do} the same job as my father.} — Avec \eng{job}, on emploie \eng{do}, pas \eng{make}.
\item \eng{He \cle{taught} me how to repair engines.} — « Apprendre à quelqu'un » = \eng{teach} (prétérit \eng{taught}) ; \eng{learn} signifie « apprendre soi-même ».
\item \eng{After my studies, I will \cle{look for} a job in Douala.} — « Chercher quelque chose » = \eng{look for} (ou \eng{search for}) ; \eng{search} seul porte sur un lieu (\eng{search the house}).
\item \eng{The carpenters are \cle{making} chairs in the workshop.} — \eng{fabricate} évoque la fabrication industrielle ou « inventer (un mensonge) » ; pour un artisan, on dit \eng{make}.
\end{enumerate}
\end{corrige}

\begin{corrige}[Exercice 8 — Appariement : dialogue après le career day]
Ordre correct : (1) \textbf{c} — (2) \textbf{e} — (3) \textbf{a} — (4) \textbf{d} — (5) \textbf{b}.

Dialogue reconstitué :
\begin{itemize}
\item \eng{A: Did you enjoy the career day yesterday?}
\item \eng{B: Yes, very much! It was so interesting to meet real professionals.}
\item \eng{A: Which profession attracted you most?}
\item \eng{B: I think I was most attracted by car mechanics.}
\item \eng{A: Really? Why that one?}
\item \eng{B: Because I love engines, and skilled mechanics earn a good living here.}
\item \eng{A: And what training does it require?}
\item \eng{B: You need about three years of training in a technical school, plus an apprenticeship.}
\item \eng{A: That sounds great. Good luck with your project!}
\item \eng{B: Thank you! And don't forget to tell me about your own choice.}
\end{itemize}
Indices : \eng{Did you enjoy...?} appelle \eng{Yes, very much!} ; \eng{Why that one?} appelle une réponse en \eng{Because} ; \eng{what training does it require?} appelle \eng{You need about three years of training...} ; \eng{Good luck!} appelle \eng{Thank you!}.
\end{corrige}

\courssub{Je me prépare au Bac}

\begin{corrige}[Exercice 9 — Compréhension écrite et questions de langue]
\textbf{A. Comprehension}
\begin{enumerate}
\item \eng{More and more young Cameroonians are turning to technical trades such as plumbing, welding, tailoring and hairdressing.}
\item \eng{Parents traditionally believed that success meant becoming a doctor, a lawyer or a civil servant.} (Deux exemples suffisent)
\item \eng{He was shocked because Sandra wanted to become a plumber, which he considered a man's job.}
\item \eng{She enrolled in a technical college, completed a six-month internship with a building company in Douala, and obtained her qualification.}
\item \eng{She means that a practical skill you can use is more valuable than a diploma that does not lead to a job.} (Toute formulation personnelle équivalente est acceptée)
\end{enumerate}
\textbf{B. Language work}
\begin{enumerate}
\item (a) \eng{self-employed} (famille de \eng{employ}) ; (b) \eng{trade} ou \eng{career} (synonymes de \eng{job}) ; (c) \eng{skilled} (« qualifié et expérimenté »).
\item \eng{Sandra \cle{enrolled} in a technical college after she \cle{had told} her parents about her project.} (Plus-que-parfait pour l'antériorité ; \eng{told} seul est accepté après \eng{after}.)
\item \eng{Her father said (that) that was a man's job.} — Recul du présent au prétérit : \eng{is} $\rightarrow$ \eng{was}.
\end{enumerate}
Barème indicatif (sur 10) : A. 5 points (1 point par question, réponse en phrase complète exigée) ; B. 5 points (question 1 : 3 points ; question 2 : 1 point ; question 3 : 1 point). Pénaliser le mot-à-mot du texte sans reformulation aux questions 4 et 5.
\end{corrige}

\begin{corrige}[Exercice 10 — Traduction]
\eng{My cousin Etienne has always been keen on carpentry. After his BEPC, he became an apprentice in the workshop of a craftsman in Bafoussam. His training lasted four years, but today he is very good with his hands. He makes furniture which he sells at the market, and he earns a good living. His dream is to open his own business and, in turn, to train young apprentices. I am proud of him, because he has chosen a trade he really loves.}

Points de vigilance :
\begin{itemize}
\item « passionné par » = \eng{keen on} ou \eng{passionate about} (jamais \eng{passionate by}) ;
\item « habile de ses mains » = \eng{good with his hands} (pas de calque) ;
\item « gagner bien sa vie » = \eng{earn a good living} ;
\item « former » = \eng{train} (jamais \eng{form}, qui signifie « constituer ») ;
\item « fier de lui » = \eng{proud of him}.
\end{itemize}
Barème indicatif (sur 10) : exactitude du lexique des métiers (4 points), temps verbaux et grammaire (3 points), absence de calques (2 points), orthographe et fluidité (1 point).
\end{corrige}

\begin{corrige}[Exercice 11 — Composition type Baccalauréat]
\textbf{Proposition de rédaction modèle} (98 mots) :

\eng{In the future, I would like to become an electrician. I have always been keen on electricity: as a child, I enjoyed repairing small radios with my uncle. What attracts me most is that skilled electricians are badly needed in Cameroon, so I will easily earn a good living. This trade requires about three years of training in a technical college, followed by an apprenticeship in a workshop to obtain a qualification. Later, I hope to be self-employed and open my own business in Yaoundé. To succeed in this profession, one needs to be careful, patient and hard-working, because electricity can be dangerous. I am determined to work hard for my dream.}

\textbf{Points de vigilance} :
\begin{itemize}
\item annoncer le métier dès la première phrase (\eng{I would like to become...}) ;
\item employer le lexique de la leçon : \eng{keen on}, \eng{training}, \eng{apprenticeship}, \eng{qualification}, \eng{self-employed}, \eng{earn a good living} ;
\item respecter le nombre de mots (80--100) et vérifier les temps (présent pour les goûts, futur pour les projets).
\end{itemize}

\textbf{Barème indicatif} (sur 10) : respect du sujet et de la consigne de longueur (2 points), richesse et exactitude du vocabulaire des métiers (3 points), correction grammaticale et temps verbaux (3 points), organisation et cohérence du paragraphe (1 point), orthographe et présentation (1 point).
\end{corrige}

\newpage

\coursec{Fiche bilan}

\begin{popfigure}
\begin{tikzpicture}[
  mindmap,
  grow cyclic,
  every node/.style={concept, text=white, font=\small\bfseries},
  root concept/.append style={concept color=popblue, font=\bfseries, minimum size=3.2cm, align=center},
  level 1 concept/.append style={level distance=4.6cm, sibling angle=60, font=\footnotesize\bfseries},
  level 2 concept/.append style={level distance=2.6cm, font=\scriptsize}
]
\node[root concept]{Jobs, Trades \& Professions}
  child[concept color=poporange]{ node {Trades, jobs, professions}
    child { node[concept color=poporange!70] {trade, craft, occupation} }
    child { node[concept color=poporange!70] {skilled / unskilled} }
  }
  child[concept color=popgreen]{ node {Developing an interest}
    child { node[concept color=popgreen!70] {to be keen on, to be drawn to} }
    child { node[concept color=popgreen!70] {vocation, calling, passion} }
  }
  child[concept color=poppurple]{ node {Sharing an interest}
    child { node[concept color=poppurple!70] {career talk, job fair} }
    child { node[concept color=poppurple!70] {to recommend, to advise} }
  }
  child[concept color=popgold]{ node {Collocations}
    child { node[concept color=popgold!70] {to take up a trade} }
    child { node[concept color=popgold!70] {to earn a living} }
  }
  child[concept color=poppink]{ node {Faux amis}
    child { node[concept color=poppink!70] {formation = training} }
    child { node[concept color=poppink!70] {a stage = an internship} }
  }
  child[concept color=popblue!60]{ node {Word families}
    child { node[concept color=popblue!40] {employ, employer, employee} }
    child { node[concept color=popblue!40] {suffixes -er, -or, -ist, -ian} }
  };
\end{tikzpicture}
\legende{Carte mentale de la leçon : le lexique des métiers et de l'intérêt professionnel.}
\end{popfigure}

\begin{aretenir}[L'essentiel de la leçon]
\textbf{1. Lexique clé (anglais -- français)}
\begin{itemize}
  \item \eng{a trade} : un métier manuel ou artisanal ; \eng{a craft} : un artisanat, un savoir-faire ; \eng{a profession} : une profession qualifiée (médecin, avocat, enseignant) ; \eng{an occupation} : un métier, une occupation (terme neutre, souvent administratif).
  \item \eng{a vocation / a calling} : une vocation ; \eng{an apprentice} : un apprenti ; \eng{an apprenticeship} : un apprentissage.
  \item \eng{a skilled / an unskilled worker} : un ouvrier qualifié / non qualifié ; \eng{a trainee} : un stagiaire, un élève en formation.
  \item \eng{to be keen on} : être passionné par ; \eng{to be interested in} : s'intéresser à ; \eng{to be drawn to} : être attiré par ; \eng{to be cut out for} : être fait pour.
  \item \eng{to take up a trade} : se lancer dans un métier ; \eng{to learn a trade} : apprendre un métier ; \eng{to earn a living} : gagner sa vie ; \eng{to make a career in} : faire carrière dans.
  \item \eng{a career talk} : une conférence d'orientation ; \eng{a job fair} : un salon de l'emploi ; \eng{job prospects} : les perspectives d'emploi.
\end{itemize}

\textbf{2. Règles à connaître par cœur}

\begin{center}
\begin{tabularx}{\linewidth}{|X|X|}\hline
\rowcolor{popblueL} \textbf{Règle} & \textbf{Exemple} \\ \hline
\eng{interested in} + nom ou verbe en \emph{-ing} & \eng{She is interested in becoming a nurse.} « Elle souhaite devenir infirmière. » \\ \hline
\eng{to look forward to} + \emph{-ing} (ici \eng{to} est une préposition, pas la marque de l'infinitif) & \eng{I look forward to starting my apprenticeship.} « J'ai hâte de commencer mon apprentissage. » \\ \hline
\eng{keen on} + \emph{-ing} ; \eng{fond of} + \emph{-ing} & \eng{He is keen on repairing engines.} « Il adore réparer les moteurs. » \\ \hline
\eng{to advise} + quelqu'un + \eng{to} + infinitif & \eng{The teacher advised me to learn a trade.} « Le professeur m'a conseillé d'apprendre un métier. » \\ \hline
Suffixes de métiers : \emph{-er} (\eng{teacher}), \emph{-or} (\eng{actor}), \emph{-ist} (\eng{dentist}), \emph{-ian} (\eng{electrician}) & \eng{My uncle is an electrician in Douala.} « Mon oncle est électricien à Douala. » \\ \hline
Famille d'\eng{employ} : \eng{employer} (employeur), \eng{employee} (employé), \eng{employment} (emploi), \eng{unemployment} (chômage) & \eng{Youth unemployment is a serious issue.} « Le chômage des jeunes est un problème grave. » \\ \hline
\end{tabularx}
\end{center}

\textbf{3. Faux amis à éviter}
\begin{itemize}
  \item \eng{training} = la formation professionnelle. Le mot anglais \eng{formation} existe mais signifie « formation géologique » ou « mise en forme » : ne jamais l'employer pour « formation professionnelle ».
  \item \eng{a career} = une carrière (sur la durée) ; \eng{a job} = un emploi précis ; \eng{work} est indénombrable : on dit \eng{I am looking for work.} et jamais \eng{a work}.
  \item \eng{actually} = en fait, réellement ; « actuellement » se dit \eng{currently} ou \eng{at present}.
  \item « Un stage » ne se dit pas \eng{a stage} (qui signifie « une scène, une étape ») mais \eng{an internship} ou \eng{a work placement}.
  \item \eng{information} est indénombrable : jamais de « s » (\eng{some useful information}).
\end{itemize}

\textbf{4. Méthodes express}
\begin{itemize}
  \item Pour présenter un métier : nommer le métier, dire en quoi il consiste (\eng{to be in charge of}, \eng{to deal with}), puis donner son avis (\eng{What I like about it is...}).
  \item Pour exprimer une vocation : \eng{I have always been fascinated by...} / \eng{Ever since I was a child, I have wanted to...}
  \item Pour échanger sur les métiers : poser des questions ouvertes (\eng{What made you choose this career? What qualifications do you need?}).
\end{itemize}
\end{aretenir}

\begin{aretenir}[Checklist avant l'examen]
\begin{itemize}
  \item $\square$ Je sais distinguer \eng{job}, \eng{work}, \eng{trade}, \eng{profession} et \eng{occupation}.
  \item $\square$ Je connais au moins dix noms de métiers avec leur suffixe (\emph{-er}, \emph{-or}, \emph{-ist}, \emph{-ian}).
  \item $\square$ J'emploie correctement \eng{interested in} + \emph{-ing} et \eng{look forward to} + \emph{-ing}.
  \item $\square$ Je sais exprimer une vocation : \eng{calling}, \eng{vocation}, \eng{to be cut out for}.
  \item $\square$ Je maîtrise les collocations : \eng{take up a trade}, \eng{earn a living}, \eng{make a career}.
  \item $\square$ Je n'écris jamais \eng{a work} ni \eng{informations} avec un « s ».
  \item $\square$ Je ne traduis pas « formation » par \eng{formation} mais par \eng{training}.
  \item $\square$ Je connais la famille d'\eng{employ} : \eng{employer}, \eng{employee}, \eng{employment}, \eng{unemployed}.
  \item $\square$ Je sais poser trois questions pour interviewer quelqu'un sur son métier.
  \item $\square$ Je sais prononcer \eng{career} (« ké-ri-eur ») et \eng{entrepreneur} (« an-tre-pre-neur »).
  \item $\square$ Je peux rédiger un court paragraphe sur le métier de mes rêves avec trois raisons.
  \item $\square$ Je révise régulièrement le lexique en le réemployant dans mes propres phrases.
\end{itemize}
\end{aretenir}

\findecours

\end{document}
